\documentclass[11pt]{amsart}
\usepackage{amsmath,amssymb,amsthm,mathtools,mathrsfs}
\usepackage[margin=1.08in]{geometry}
\usepackage{booktabs}
\usepackage{array}
\usepackage{longtable}
\usepackage{tabularx}
\usepackage{hyperref}
\hypersetup{hidelinks}
\usepackage[final,expansion=false]{microtype}
\setlength{\parskip}{4pt plus 1pt}
\emergencystretch=2.5em

\newtheorem{theorem}{Theorem}[section]
\newtheorem{proposition}[theorem]{Proposition}
\newtheorem{lemma}[theorem]{Lemma}
\newtheorem{corollary}[theorem]{Corollary}
\theoremstyle{definition}
\newtheorem{definition}[theorem]{Definition}
\newtheorem{convention}[theorem]{Convention}
\theoremstyle{remark}
\newtheorem{remark}[theorem]{Remark}

\newcommand{\C}{\mathbb C}
\newcommand{\R}{\mathbb R}
\newcommand{\N}{\mathbb N}
\newcommand{\Z}{\mathbb Z}
\newcommand{\Q}{\mathbb Q}
\newcommand{\A}{\mathbb A}
\newcommand{\dd}{\,d}
\newcommand{\dstar}{\,d^{*}}
\newcommand{\Tr}{\operatorname{Tr}}
\newcommand{\Rel}{\operatorname{Re}}
\newcommand{\Pic}{\operatorname{Pic}}
\newcommand{\Det}{\operatorname{Det}}
\newcommand{\Picb}{\operatorname{Pic}_{b}}
\newcommand{\Hom}{\operatorname{Hom}}
\newcommand{\Cone}{\operatorname{Cone}}
\newcommand{\supp}{\operatorname{supp}}
\newcommand{\vol}{\operatorname{vol}}
\newcommand{\sgn}{\operatorname{sgn}}
\newcommand{\id}{\operatorname{id}}
\newcommand{\HH}{\operatorname{HH}}
\newcommand{\Sh}{\operatorname{Sh}}
\newcommand{\Perf}{\operatorname{Perf}}
\newcommand{\Aut}{\operatorname{Aut}}
\newcommand{\Res}{\operatorname{Res}}
\newcommand{\Inf}{\operatorname{Inf}}
\newcommand{\Rest}{\operatorname{Res}}
\newcommand{\Eff}{\operatorname{Eff}}
\newcommand{\Th}{\widehat{T}}
\newcommand{\Wh}{\widehat{W}}
\newcommand{\Ch}{\widehat{C}_S}
\newcommand{\DA}{D_A}
\newcommand{\Atr}{\mathcal A_{\mathrm{tr}}}
\newcommand{\Corrtr}{\mathrm{Corr}^{\mathrm{tr}}_{A}}
\newcommand{\Nrad}{\mathcal N}
\newcommand{\BT}{B_T}
\newcommand{\NAd}{N_A^{\mathrm{dist}}}
\newcommand{\Log}{\operatorname{Log}}
\newcolumntype{Y}{>{\raggedright\arraybackslash}X}

\begin{document}

\title[A correspondence algebra on the idele class group]
{A Correspondence Algebra with a Canonical Trace\\
on the Idele Class Group,\\
its Tate Symbol and its Numerical Form}

\author{David Alejandro Trejo Pizzo}
\thanks{Independent Researcher. \texttt{dtrejopizzo@gmail.com}}

\date{\today}

\begin{abstract}
Let $S$ be a finite set of places of $\Q$ containing $\infty$, let
$A_S=\R\times\prod_{p\in S_f}\Q_p$, let $\Gamma_S$ be the group of positive
$S$-units and $C_S=\Gamma_S\backslash A_S^{\times}$ the semilocal idele class
group.  On the square of the quotient $X_S=[A_S/\Gamma_S]$ we build a category
of correspondences generated by the graph bimodules of idele classes, prove
two-sided excision for it, and equip it with an involution, the half-density
Hilbert adjoint; the product formula fixes the square root of the modulus
that this requires.  A splitting of the valuation map identifies $C_S$ with
$\R^{\times}\times\prod_{p\in S_f}\Z_p^{\times}$ and turns the smeared
correspondences into the commutative convolution $\dagger$-algebra
$D_A=C_c^{\infty}(C_S)$.

On $D_A$ we define a trace $T$ as a sum of local fixed-point distributions
with densities $\tau_v(x)=|x|_v^{1/2}/|1-x|_v$.  It converges absolutely away
from the identity class, and it is unique there up to a distribution
supported at the identity.  Its Fourier symbol on the dual group is
real-analytic, bounded above and tends to $-\infty$ along every component;
each local symbol is a Mittag--Leffler sum over poles at height $\pm\frac12$
with residue $\mp i$.  Consequently the Hermitian form
$I_A(F,G)=T(F\star G^{\dagger})$ is nondegenerate on $D_A$, and it has
infinite index of both signs.

The two values $\widehat F(\mp i/2)$ define a surjection of $D_A$ onto a
hyperbolic plane with nonzero kernel.  It has no module splitting inside
$D_A$, but in a separated completion the two rulings appear as limits of
spread Tate characters, they span a canonical complement of the primitive
part, and the scalar form vanishes on them.  We also construct a Banach
$\dagger$-algebra on which $T$ and the determinant $\exp T$ are continuous
and the determinant germ is holomorphic and multiplicative, a proper
effective cone, and a Riemann--Roch biextension with values in pairs of based
lines, whose second component is the hyperbolic form of the Tate degrees; it
has the shape, formally, of Weil's decomposition $D\cdot D'=P\cdot
P'+d_1d_2'+d_2d_1'$, as a pair rather than a sum.  At a fixed $S$ the second
component cannot be dropped, because the scalar determinant is trivial on the
rulings.  An appendix, which uses Mertens' theorem and is not part of the
construction, shows that in the limit over all $S$ the scalar form becomes
hyperbolic on the rulings with structure constant $1$.
\end{abstract}

\maketitle

\tableofcontents

\section{Introduction}\label{sec:intro}

On a smooth projective curve $C$ over a finite field, the correspondences on
$C\times C$ carry the structure of Weil's theory of correspondences
\cite{Weil1948}: a composition (pull, tensor and push
over $C^{3}$), a unit (the diagonal), a transpose, an intersection pairing, a
notion of effectivity, and a Riemann--Roch determinant.  The two rulings
$C\times\{pt\}$ and $\{pt\}\times C$ span a hyperbolic plane with Gram matrix
$\left(\begin{smallmatrix}0&1\\1&0\end{smallmatrix}\right)$ in the numerical
group, and every correspondence decomposes into a ruling part, recorded by
its bidegree, and a primitive part.

In this paper we construct a theory with the same operations on an adelic
space.  The space is the quotient $X_S=[A_S/\Gamma_S]$ of a semilocal adelic
space by the positive $S$-units, the correspondences are the graph bimodules
of idele classes acting on it, and the intersection pairing is an analytic
trace: a sum of local fixed-point distributions whose densities are the
Fourier orbital integrals of the dilations.  Apart from the product formula,
no arithmetic input enters the construction; in particular no explicit
formula and no positivity statement is used.

Three features of this theory have no counterpart for curves, and they
determine its shape.  First, $C_S$ is a continuum, so there are infinitely
many pairwise non-isomorphic generators, and the pairing is a Hermitian form
on a nuclear space rather than a bilinear form on a finitely generated group.
Second, the Hochschild invariants of a graph bimodule do not see the
dilation.  At a finite place the fixed point of a dilation is the origin of
$\Q_p$; the coefficient algebra is totally disconnected, so this fixed point
carries no transverse data, and its Hochschild contribution is a rank one
class independent of the dilation (Theorem~\ref{thm:shadow}).  Ranks, Euler
characteristics and determinants of these complexes are therefore either
constant along the relevant families or undefined
(Corollary~\ref{cor:blind}).  The local statements are proved; assembling them
over the places of $S$ uses a twisted form of Wodzicki's excision theorem,
which we isolate as a hypothesis (Theorem~\ref{thm:excision-input}).  What does depend on the dilation is the
normalization of that rank one class by the half-density and the Haar
measure, that is, the analytic trace.  Third, the two ruling directions are
not represented by elements of the test algebra.  They are eigenvectors for
the whole convolution algebra with eigenvalue $\widehat a(\mp i/2)$, a value
of the Fourier transform off the unitary axis, whereas the transform of a
test function is only seen on the axis (Theorem~\ref{thm:nosplit}).  They
exist only in a completion, as limits of test functions whose mass escapes to
infinity at the rate at which a single Tate weight carries all of it
(Theorem~\ref{thm:rulings}).

\subsection*{Main results}

Fix a finite set $S$ of places of $\Q$ with $\infty\in S$ and $S_f\ne
\emptyset$, and write $A_S=\R\times\prod_{p\in S_f}\Q_p$ and
$\Gamma_S=\Z_S^{\times}$ for the group of positive rationals whose prime
factorization uses only the primes of $S_f$.  The first result is elementary,
and it is what makes the construction canonical.

\begin{theorem}[Lemma~\ref{thm:split} and
Corollary~\ref{cor:trivial}]\label{thm:introA}
The map $\varphi(x)=\prod_{p\in S_f}p^{v_p(x_p)}$ is a continuous
homomorphism $A_S^{\times}\to\Gamma_S$ restricting to the identity on
$\Gamma_S$.  Hence
\[
A_S^{\times}=\Gamma_S\times\ker\varphi,\qquad
C_S:=\Gamma_S\backslash A_S^{\times}\cong\ker\varphi=
\R^{\times}\times\prod_{p\in S_f}\Z_p^{\times},
\]
and $|x|_S=|x_{\infty}|_{\infty}$ on $\ker\varphi$.  The resulting
homomorphic section of the class map trivializes the Fell bundle of idele
classes, so that $D_A:=C_c^{\infty}(C_S)$, with convolution and
$F^{\dagger}(u)=\overline{F(u^{-1})}$, is a commutative $\dagger$-algebra,
and $F\mapsto E_F$ is an injective monoid map onto smeared correspondences.
\end{theorem}

For a place $v$ of $\Q$ and $x\in\Q_v^{\times}$, $x\ne1$, put
\[
\tau_v(x):=\frac{|x|_v^{1/2}}{|1-x|_v}.
\]
At a finite place $\tau_p(p^ku)=p^{-|k|/2}$ for $k\ne0$, the half-density of
the dilation, and $\tau_p(u)=|1-u|_p^{-1}$ for a unit $u\ne1$; for
$u\in1+q_p\Z_p$ ($q_p=p$, $q_2=4$) this is the number of $u$-orbit classes of
characters on a valuation shell divided by $1-p^{-1}$, and for the other units
there is no such interpretation (Lemma~\ref{lem:tau},
Proposition~\ref{prop:haarrank}).  The trace is
\[
T(F)=\sum_{v\in S}W_v(F)+c\,F(1),
\]
where $W_p$ is the shell integral plus a regularized integral over the units,
$W_{\infty}$ is the archimedean fixed-point distribution, and $c$ is a real
constant (Definition~\ref{def:trace}).  On test functions vanishing near the
identity class, $T$ is an absolutely convergent sum of integrals
(Theorem~\ref{thm:absconv}), and it is the unique such functional up to a
distribution supported at the identity class (Theorem~\ref{thm:unique}).
The central computation of the paper is the Fourier symbol of $T$.

\begin{theorem}[Theorems~\ref{thm:symbol} and \ref{thm:ML}]\label{thm:introB}
Write $\widehat C_S=\{\chi=(\varepsilon,\tau;(\kappa_q))\}$, a countable
disjoint union of copies of $\{0,1\}\times\R$.  Then
$T(H)=\int_{\widehat C_S}\widehat H(\chi)(\widehat T(\chi)+c)\,d\chi$, where
$\widehat T=\widehat W_{\infty}+\sum_{p\in S_f}\widehat W_p$ is
real-valued, real-analytic on each component, bounded above, and
$\widehat T(\varepsilon,\tau;(\kappa_q))\to-\infty$ as $|\tau|\to\infty$.  At
a finite place
\[
\widehat W_p(\chi)=
\begin{cases}
2(\log p)\,\Rel\bigl(\zeta_p/(1-\zeta_p)\bigr), & \kappa_p=1,\\[2pt]
-c(\kappa_p)\log p, & \kappa_p\neq1,
\end{cases}
\qquad \zeta_p=p^{-1/2}\overline{\eta_p(\chi)},
\]
and each $\widehat W_p(\kappa_p=1)$ extends to a meromorphic
$2\pi/\log p$-periodic function of $\tau$ whose poles are simple, lie at
height $\pm\tfrac12$, and have residue $\mp i$; the even archimedean symbol is
meromorphic with simple poles at $\pm i(2n+\tfrac12)$ and the same residues.
On the principal component every place of $S$ contributes one pole at $\pm
i/2$, so the total residue there is $\mp i(|S_f|+1)$.
\end{theorem}

Only three properties of $\widehat T$ are needed for the next theorem: it is
real-analytic, its finite-place terms are bounded, and its archimedean term
tends to $-\infty$.

\begin{theorem}[Theorem~\ref{thm:radical}]\label{thm:introC}
Let $\Nrad=\{F\in D_A: T(F\star G^{\dagger})=0\text{ for all }G\in D_A\}$.
Then $\Nrad=\{0\}$ for every value of $c$.  Hence the numerical quotient
$N_A=D_A/\Nrad$ is $D_A$ itself, and the Hermitian form
$I_A(F,G)=T(F\star G^{\dagger})$ is nondegenerate on the whole test algebra.
\end{theorem}

By the product formula the modulus $|\cdot|_S$ is trivial on $\Gamma_S$, so
$\lambda_{\mp}(F)=\int_{C_S}F(u)|u|_S^{\mp1/2}\dstar u$ are well-defined
characters of the convolution algebra.  They are the values $\widehat
F(\mp i/2)$ of the entire continuation of $\widehat F$, and they are the
residue contributions of the symbol at its poles $\pm i/2$.  With
$M=(\lambda_-,\lambda_+)$, $B_T=\C_-\oplus\C_+$ and
$\beta_T(z,w)=z_-\overline{w_+}+z_+\overline{w_-}$ there is a canonical exact
sequence of topological $D_A$-modules
\[
0\longrightarrow D_A^{\mathrm{prim}}\longrightarrow N_A
\stackrel{M}{\longrightarrow}B_T\longrightarrow 0,
\]
in which the involution exchanges the two basis vectors of $B_T$, and whose
kernel is nonzero and carries a nondegenerate induced pairing
(Theorem~\ref{thm:tateseq}).  That the kernel is nonzero is shown by an
explicit computation: two test functions supported on cosets
$u(1+p^m\Z_p)$, $u'(1+p^m\Z_p)$ with $v_p(1-u)=j\ne j'=v_p(1-u')$ have the
same image under $M$, and their traces differ by $\nu(\log p)(p^{j}-p^{j'})$
(Theorem~\ref{thm:separate}).  The sequence has no $D_A$-module section
(Theorem~\ref{thm:nosplit}).  In the separated completion
$N_A^{\mathrm{dist}}$, however, the rulings appear
(Theorems~\ref{thm:rulings}, \ref{thm:splitting}, \ref{thm:rulnull}): the
completion splits canonically as the sum of a primitive part and the ruling
plane, and the ruling plane is exactly the radical of the scalar component of
the pairing.  The signature of the form is as follows.

\begin{theorem}[Theorem~\ref{thm:signature}]\label{thm:introD}
$I_A(F,F)=\int_{\widehat C_S}|\widehat F(\chi)|^{2}(\widehat T(\chi)+c)\,
d\chi$ is real.  For every $c$ the form has infinite negative index; it has
infinite positive index if and only if $c>-m$, where $m=\sup\widehat T$, and
$m\ge\log\pi+\gamma+3\log2+\pi/2>0$.  In particular for $c=0$ the form is
indefinite with infinite index on both sides.
\end{theorem}

For the determinant we introduce the norm
$\lVert F\rVert_{\mathrm{tr}}=\Lambda(F)+N(F)$, where $\Lambda$ weights $F$ by
$|w|_S^{1/2}+|w|_S^{-1/2}$ and $N$ weights $\widehat F$ by $|\widehat
T|+|c|$.  The completion $\Atr$ is a commutative Banach $\dagger$-algebra on
which $T$, $\exp T$ and both Tate characters are continuous, and on which the
determinant germ
$\log\Delta_F(z)=\sum_{r\ge1}(-1)^{r+1}z^{r}T(F^{\star r})/r$ is holomorphic
for $|z|<1/\lVert F\rVert_{L^1}$ and satisfies
$\Delta_F\Delta_G=\Delta_{F+G+zF\star G}$ (Theorem~\ref{thm:banach}).  The
map $H\mapsto E_H$ extends to $\Atr$ as a bounded, faithful,
$\dagger$-compatible representation, and the category of presented complexes
over $\Atr$ has the same formal properties as over $D_A$
(Theorem~\ref{thm:Atrcat}).

The scalar determinant $L_S(F,G)=(\C,e^{I_A(F,G)})$ is trivial on the
rulings: along every admissible approximation of $e_{\pm}$ the scalar pairing
tends to $0$, while $\beta_T$ tends to $1$ (Theorems~\ref{thm:rulnullclass}
and \ref{thm:scalartrivial}).  The determinant must therefore take values in
$\Picb(B_T)$, the Picard groupoid of based invertible $B_T$-modules, and the
Riemann--Roch biextension is the pair
$\lambda_{RR}=(L_S,\lambda_{\mathrm{per}}\circ(M\times M))$.  Its logarithm is
the vector-valued pairing $(I_A,\beta_T\circ(M\times M))$, a pair which has
the shape of Weil's decomposition $D\cdot D'=P\cdot P'+d_1d_2'+d_2d_1'$; the
first component is the full $I_A(F,G)$, not a primitive pairing, so the
analogy is formal (Theorem~\ref{thm:RR}).

The need for a second component is a feature of a fixed finite set $S$.  The
sets $S$ form a filtered system with retractions, along which the trace
stabilizes and becomes invariant (Theorems~\ref{thm:chartsystem} and
\ref{thm:colimtrace}); in the limit the scalar form is hyperbolic on the
rulings, with structure constant $1$, which is the logarithmic density of
prime powers given by Mertens' theorem (Theorem~\ref{thm:mertens}).  That
statement uses the von Mangoldt function and Mertens' theorem, which the rest
of the paper does not; it is placed in Appendix~\ref{app:mertens} and is not
part of the construction.

Sections~\ref{sec:carrier}--\ref{sec:shadow} construct the space, the test
algebra, the correspondence category with its involution, and compute the
Hochschild invariants of graph bimodules.  Sections~\ref{sec:trace} and
\ref{sec:symbol} define the trace and compute its symbol.
Section~\ref{sec:numerical} studies the pairing, the Tate sequence and the
completion.  Sections~\ref{sec:presented}--\ref{sec:charts} treat presented
complexes, the Banach algebra and the change of $S$;
Sections~\ref{sec:eff}--\ref{sec:classical} treat effectivity, the
Riemann--Roch biextension and the comparison with curves; and
Section~\ref{sec:signature} computes the signature and the limit over $S$.

\section{The semilocal space and the test algebra}\label{sec:carrier}

Throughout, $S$ is a finite set of places of $\Q$ with $\infty\in S$ and
$S_f:=S\setminus\{\infty\}\ne\emptyset$.  We write
\[
A_S:=\R\times\prod_{p\in S_f}\Q_p,\qquad
\Gamma_S:=\Z_S^{\times}=\Bigl\{\textstyle\prod_{p\in S_f}p^{e_p}: e_p\in\Z\Bigr\}
\subset\Q_{>0},
\]
$\Gamma_S$ being embedded diagonally in $A_S^{\times}$.  It is a free abelian
group of rank $|S_f|$, and it acts on $A_S$ by multiplication.  Our basic
space is the quotient groupoid
\begin{equation}\label{eq:carrier}
X_S:=[A_S/\Gamma_S],\qquad Y_S:=X_S\times X_S .
\end{equation}
The two projections of $Y_S$ are distinct anchor maps, so a kernel on $Y_S$
records both actions.

We equip $A_S$ with the additive product Haar measure $\mu$ normalized so that
the unit ball at each finite place has volume one, and we write
$|a|_S=\prod_{v\in S}|a_v|_v$ for the module of $a\in A_S^{\times}$.

\begin{lemma}[Product formula]\label{lem:prodformula}
$|\gamma|_S=1$ for every $\gamma\in\Gamma_S$.  Consequently multiplication by
$\gamma$ preserves $\mu$.
\end{lemma}

\begin{proof}
A positive $S$-unit $\gamma$ has $|\gamma|_v=1$ at every place $v\notin S$, so
the global product formula $\prod_{v}|\gamma|_v=1$ reduces to
$\prod_{v\in S}|\gamma|_v=1$.
\end{proof}

\subsection{Coefficients and the source algebra}

At a finite place $p$ let $F_p$ be the nuclear Fr\'echet algebra of functions
on $\Q_p$ with rapidly decaying shell--character Fourier coefficients: writing
$j\in\Z$ for the valuation shell $p^{j}\Z_p^{\times}$ and $\kappa$ for the
characters of $\Z_p^{\times}$ with conductor exponent $c_p(\kappa)$, the
seminorms are
\begin{equation}\label{eq:seminorms}
v_{l,N}(f)=\sum_{j,\kappa}p^{l|j|}\,c_p(\kappa)^{N}\,|f_{j,\kappa}|<\infty ,
\qquad l,N\ge0 .
\end{equation}
Elements of $F_p$ vanish to infinite order at $x_p=0$ and as $x_p\to\infty$.
Put $E_p:=\C\cdot\chi_p\oplus F_p$ with $\chi_p=\mathbf 1_{\Z_p}$, the algebra
unitized at the zero coordinate; $\chi_p$ carries the value at the origin.  At
the archimedean place, $E_{\R}$ is the algebra of smooth functions $f$ on
$\R^{\times}$ such that, on each of the two half-lines and in the logarithmic
coordinate $t=\log|x_{\infty}|$,
$\sup_t e^{l|t|}|D^{m}f|<\infty$ for all $l,m$; thus $E_{\R}$ is the direct
sum of two copies of the algebra of rapidly decreasing smooth functions of
$t\in\R$, one for each sign of $x_{\infty}$.  An idele $a$ acts on it by
$f\mapsto f(a_{\infty}^{-1}\,\cdot\,)$, that is, by translation by
$\log|a_{\infty}|$ in $t$, composed with the exchange of the two half-lines
when $a_{\infty}<0$; for $a_{\infty}\ne1$ this action has no fixed point.  The
common coefficient algebra is the completed projective tensor product
\begin{equation}\label{eq:coeff}
E_P:=E_{\R}\,\widehat\otimes\,\bigotimes_{p\in S_f}E_p ,
\end{equation}
nuclear, with tempered $\Gamma_S$-action.  The \emph{source algebra} is the
rapid crossed product
\begin{equation}\label{eq:source}
U:=(E_P\oplus\C 1)\rtimes_{\mathrm{rapid}}\Gamma_S ,
\end{equation}
a unital complete Fr\'echet algebra with gauge unitaries $U_{\gamma}$,
$\gamma\in\Gamma_S$.  We write $N:=E_P\rtimes_{\mathrm{rapid}}\Gamma_S$ for
its nonunital coefficient ideal.  The seminorms \eqref{eq:seminorms} satisfy
$v_{l,N}(fg)\le v_{l,N}(f)v_{l,N}(g)$.

Here \emph{rapid} refers to the height $H(a/b):=\max(a,b)$ of a positive
rational written in lowest terms; $H$ is symmetric,
$H(\gamma^{-1})=H(\gamma)$, and submultiplicative,
$H(\gamma\delta)\le H(\gamma)H(\delta)$.  For a complete Fr\'echet algebra $A$
with an action $\alpha$ of $\Gamma_S$, the rapid crossed product
$A\rtimes_{\mathrm{rapid}}\Gamma_S$ consists of the sums
$\sum_{\gamma}a_{\gamma}U_{\gamma}$ with
$\sum_{\gamma}H(\gamma)^{k}\nu(a_{\gamma})<\infty$ for every $k\ge0$ and every
continuous seminorm $\nu$ of $A$, with the product determined by
$U_{\gamma}aU_{\gamma}^{-1}=\alpha_{\gamma}(a)$, where
$\alpha_{\gamma}(a)(x)=a(\gamma^{-1}x)$.  The action on $E_P$ is
\emph{tempered}: for every continuous seminorm $\nu$ there are $\nu'$, $C$ and
$M$ with $\nu(\alpha_{\gamma}f)\le C\,H(\gamma)^{M}\nu'(f)$.  Indeed, at a
finite place $\gamma$ shifts the shell index by $v_p(\gamma)$, which costs at
most $p^{l|v_p(\gamma)|}\le H(\gamma)^{l}$ in \eqref{eq:seminorms}, and moves
$\chi_p$ by the difference of two indicator functions supported on
$|v_p(\gamma)|$ shells; at the real place it translates $t$ by
$\log\gamma$, which costs at most $e^{l|\log\gamma|}\le H(\gamma)^{l}$.
Temperedness and submultiplicativity of $H$ make the product continuous.

Modules are unital complete Fr\'echet $U$-bimodules; an exact sequence is
\emph{admissible} if it splits continuously and linearly; $D^{-}_{\mathrm{rel}}(U\text{-}U)$
is the bounded-above relative derived category.  Free bimodules are relatively
projective and bounded-above complexes of relative projectives are K-flat on
either side.  We write
\begin{equation}\label{eq:comp}
M\star N:=Q(M)\,\widehat\otimes_U\,Q(N),\qquad \Delta:=U,
\end{equation}
$Q$ denoting a relatively projective resolution, and
\begin{equation}\label{eq:shadow}
\Sh(M):=\mathrm{Tot}\bigl[M\,\widehat\otimes\,U^{\widehat\otimes n}\bigr]_{n\ge0},
\qquad \Sh(M\star N)\cong\Sh(N\star M)
\end{equation}
for the relative Hochschild shadow.  The symbol $\Sh$ is reserved for this
categorical invariant; it is never identified with the analytic trace of
Section~\ref{sec:trace}.

\subsection{The valuation splitting}

\begin{lemma}\label{thm:split}
Define
\begin{equation}\label{eq:phi}
\varphi:A_S^{\times}\longrightarrow\Gamma_S,\qquad
\varphi(x):=\prod_{p\in S_f}p^{\,v_p(x_p)} .
\end{equation}
Then $\varphi$ is a continuous group homomorphism with
$\varphi|_{\Gamma_S}=\id$, and therefore
\begin{equation}\label{eq:splitting}
A_S^{\times}=\Gamma_S\times\ker\varphi,\qquad
\ker\varphi=\R^{\times}\times\prod_{p\in S_f}\Z_p^{\times},\qquad
C_S:=\Gamma_S\backslash A_S^{\times}\cong\ker\varphi .
\end{equation}
Moreover $|x|_S=|x_{\infty}|_{\infty}$ for $x\in\ker\varphi$.
\end{lemma}

\begin{proof}
Each $v_p$ is a homomorphism $\Q_p^{\times}\to\Z$, so $\varphi$ is a
homomorphism, and it is continuous because $v_p$ is locally constant on
$\Q_p^{\times}$.  If $\gamma=\prod_{p\in S_f}p^{e_p}\in\Gamma_S$ is embedded
diagonally then $v_p(\gamma)=e_p$ for $p\in S_f$, so $\varphi(\gamma)=\gamma$.
A homomorphism from an abelian group onto a subgroup restricting to the
identity on that subgroup splits it as the direct product of the subgroup and
the kernel.  The kernel is $\{x: v_p(x_p)=0\text{ for all }p\in S_f\}$, which
is \eqref{eq:splitting}.  For such $x$ one has $|x_p|_p=1$ at every finite
place of $S$, so $|x|_S=|x_{\infty}|$.
\end{proof}

We fix once and for all the resulting section
\begin{equation}\label{eq:section}
\sigma: C_S\;\xrightarrow{\ \cong\ }\;\ker\varphi\;\hookrightarrow\;A_S^{\times}
\end{equation}
of the class map; it is a homomorphism.  We normalize Haar measure
$\dstar w$ on $C_S\cong\R^{\times}\times\prod_{p\in S_f}\Z_p^{\times}$ so that
each compact factor $\Z_p^{\times}$ has total mass one and the archimedean
factor carries
\begin{equation}\label{eq:archmeasure}
\dstar x_{\infty}=\frac{dx}{2|x|},
\end{equation}
so that the two signs together have total mass one per unit of logarithmic
length.  In the coordinate $t=\log|w|_S$ we therefore have
$\dstar w=dt\otimes\mu_{\mathrm{cpt}}$ with $\mu_{\mathrm{cpt}}$ a probability
measure on the compact group $\{\pm1\}\times\prod_{p\in S_f}\Z_p^{\times}$.

\begin{remark}\label{rem:normdict}
Two normalizations of the involution occur in this subject.  The
\emph{unnormalized} one is $f^{\sharp}(x)=x^{-1}\overline{f(x^{-1})}$, paired
with the local density $|1-u|_v^{-1}$; the \emph{half-density} one, used
throughout this paper, is $F^{\dagger}(u)=\overline{F(u^{-1})}$, paired with
$\tau_v=|x|_v^{1/2}|1-x|_v^{-1}$.  The substitution
$F\mapsto|\cdot|_S^{1/2}F$ is a topological algebra isomorphism carrying the
first convention to the second.  All formulas below are in the half-density
convention.
\end{remark}

\subsection{The test algebra}

\begin{definition}\label{def:DA}
$\DA:=C_c^{\infty}(C_S)$ is the space of functions on $C_S$ which are smooth
in the $\R^{\times}$ coordinate, locally constant in the
$\prod_p\Z_p^{\times}$ coordinates, and compactly supported.  It carries the
inductive limit of the Fr\'echet spaces of functions supported in a fixed
compact set and invariant under a fixed open subgroup of the compact part; this
topology is nuclear and of type LF.  We equip $\DA$ with the convolution
\[
(F\star G)(w)=\int_{C_S}F(u)G(u^{-1}w)\dstar u
\]
and the involution $F^{\dagger}(u)=\overline{F(u^{-1})}$.  We write
\[
\DA^{\mathrm{off}}:=\{F\in\DA:\ F\equiv0\ \text{on a neighbourhood of }1\}.
\]
\end{definition}

$\DA$ is a commutative $\dagger$-algebra: associativity and commutativity of
convolution are Fubini's theorem on the locally compact abelian group $C_S$,
and $(F\star G)^{\dagger}=F^{\dagger}\star G^{\dagger}$ follows from
inversion-invariance of Haar measure.  A \emph{bounded} subset of $\DA$ lies in
a single Fr\'echet step: there are a compact $K\subseteq C_S$ and an integer
$N\ge0$ such that all its members are supported in $K$, are invariant under
$U_N:=\prod_{q\in S_f}(1+q^{N}\Z_q)$, and have uniformly bounded step
seminorms.

\begin{remark}\label{rem:approxid}
$\DA$ has an approximate identity: the normalized characteristic functions of
a neighbourhood basis of $1$ by sets of the form (interval) $\times$
(compact-open subgroup) lie in $\DA$, and $F\star e_{\varepsilon}\to F$ in the
fixed-support Fr\'echet step determined by a fixed compact neighbourhood of
$\supp F$.
\end{remark}

\section{The correspondence category and excision}\label{sec:cat}

\subsection{Graph bimodules}

An external idele $a\in A_S^{\times}$ acts on $A_S$ by multiplication,
commuting with $\Gamma_S$, and induces a map $m_{[a]}$ on $X_S$ depending only
on the class $[a]\in C_S$.  Its graph is a closed subgroupoid of $Y_S$ and its
structure kernel is the bimodule
\begin{equation}\label{eq:graphbimod}
K_{[a]}\cong E_a:={}_{1}U_{\beta_a},
\end{equation}
namely $U$ with the right action twisted by the automorphism $\beta_a$ induced
by $a$.  The composition of two graphs has no excess:
\begin{equation}\label{eq:starmult}
E_u\star E_v\cong E_{uv},\qquad x\otimes y\mapsto x\,\beta_u(y),
\end{equation}
with inverse $z\mapsto z\otimes 1$ and a strict pentagon; and $E_1=\Delta=U$.

The morphisms between graph bimodules are classified as follows.  Here, for
$a\in A_S^{\times}$, $\beta_a$ is the automorphism of $U$ which acts on $E_P$
by $f\mapsto f(a^{-1}\,\cdot\,)$ and fixes every $U_{\gamma}$; it commutes with
the lattice action because multiplication is commutative, so it is an
automorphism of the crossed product.

\begin{proposition}\label{prop:hom}
For $a,b\in A_S^{\times}$,
\begin{equation}\label{eq:homclass}
\Hom_{U\text{-}U}(E_a,E_b)=
\begin{cases}
\C\cdot U_{\gamma}, & a=\gamma b\ \text{with}\ \gamma\in\Gamma_S,\\
0, & \text{otherwise},
\end{cases}
\end{equation}
the element $U_{\gamma}$ acting by right multiplication; since $\Gamma_S$ acts
freely on $A_S^{\times}$, at most one $\gamma$ occurs.  Consequently every
bounded complex of finite sums of the $E_a$ is homotopy equivalent, in the
category $G$ of Definition~\ref{def:G} below, to a finite sum of shifts of
generators $E_a$.
\end{proposition}

\begin{proof}
Let $D:=\R^{\times}\times\prod_{p\in S_f}\Q_p$.  The elements of $F_p$ are
continuous functions on $\Q_p$ vanishing at $x_p=0$, their shell expansions
converging uniformly, and the elements of $E_{\R}$ are continuous on
$\R^{\times}$; so the elements of $E_P\oplus\C1$ are continuous functions on
$D$.  The map $E_P\oplus\C1\to C(D)$ is injective: on $E_P$ this follows from
the injectivity of each factor map and the approximation property of the
nuclear factors, and a nonzero constant does not lie in the image of $E_P$,
whose elements tend to $0$ as $|\log|x_{\infty}||\to\infty$.  Let
$\Omega:=\R^{\times}\times\prod_{p\in S_f}\Q_p^{\times}$, a dense open subset
stable under $A_S^{\times}$.  The algebra $E_P$ separates the points of
$\Omega$: $E_{\R}$ contains $C_c^{\infty}(\R^{\times})$, $E_p$ contains the
indicator function of every compact open subset of $\Q_p^{\times}$ (a finite
combination of shell idempotents times characters of $\Z_p^{\times}$), and
products of such functions separate the points of the product.

A bimodule map $f:E_a\to E_b$ satisfies $f(x)=xz$ with $z:=f(1)$, by left
linearity; and $f(x\beta_a(y))=f(x)\beta_b(y)$ for all $x,y$ says exactly
that
\begin{equation}\label{eq:intertw}
\beta_a(y)\,z=z\,\beta_b(y)\qquad(y\in U).
\end{equation}
Conversely every $z$ with \eqref{eq:intertw} defines a continuous bimodule
map.  Write $z=\sum_{\gamma}z_{\gamma}U_{\gamma}$ with $z_{\gamma}\in
E_P\oplus\C1$; the coefficient maps are continuous.  Take $y=c\in E_P$.  Since
$U_{\gamma}\beta_b(c)=\alpha_{\gamma}(\beta_b(c))U_{\gamma}$, comparing the
coefficients of $U_{\gamma}$ in \eqref{eq:intertw} gives
\[
\bigl(c(a^{-1}x)-c(b^{-1}\gamma^{-1}x)\bigr)\,z_{\gamma}(x)=0
\qquad(x\in D,\ c\in E_P).
\]
If $a\ne\gamma b$, then for $x\in\Omega$ the points $a^{-1}x$ and
$(\gamma b)^{-1}x$ of $\Omega$ are distinct, because every coordinate of $x$ is
nonzero; choosing $c$ which separates them gives $z_{\gamma}(x)=0$.  So
$z_{\gamma}$ vanishes on $\Omega$, hence on $D$ by continuity and density,
and therefore $z_{\gamma}=0$ by injectivity.
Thus $z=0$ unless $a=\gamma b$ for some $\gamma$, which is then unique.

Suppose $a=\gamma b$ and $z=hU_{\gamma}$.  Taking $y=U_{\delta}$ in
\eqref{eq:intertw}, and using that $\beta_a,\beta_b$ fix $U_{\delta}$ and that
$\Gamma_S$ is abelian, gives $\alpha_{\delta}(h)=h$ for every
$\delta\in\Gamma_S$.  Write $h=\lambda+e$ with $\lambda\in\C$ and $e\in E_P$;
then $e$ is invariant.  The seminorm $q(e):=\sup_{x\in D}e^{|t(x)|}|e(x)|$,
$t(x)=\log|x_{\infty}|$, is continuous on $E_P$: on elementary tensors it is at
most the product of the seminorm $\sup_te^{|t|}|f|$ of the real factor and
the sup norms of the finite factors, which are bounded by $|\lambda_p|+
v_{0,0}(f_p)$ for $\lambda_p\chi_p+f_p\in E_p$.  Fix $p\in S_f$.  Invariance
gives $e(x)=e(p^{n}x)$ for all $n$, and $t(p^{n}x)=t(x)+n\log p$, so
$|e(x)|\le q(e)e^{-|t(x)+n\log p|}\to0$.  Hence $e$ vanishes on $D$, so $e=0$, and $z=\lambda
U_{\gamma}$.  Conversely, if $a=\gamma b$ then
$U_{\gamma}\beta_b(y)U_{\gamma}^{-1}=\beta_a(y)$ for $y\in E_P$, since
$\alpha_{\gamma}\beta_b=\beta_{\gamma b}$, and for $y=U_{\delta}$ trivially; so
$z=U_{\gamma}$ satisfies \eqref{eq:intertw}.  The nonzero $\Hom$ space is one
dimensional, so the identification is topological.

For the last assertion recall that the objects of $G$ are built from the
generators $E_a$ with $a$ running over a fixed set of lifts of the classes of
$C_S$.  By
\eqref{eq:homclass}, $U_{\gamma}$ is an isomorphism $E_{\gamma a}\to E_a$, the
endomorphisms of $E_a$ are the scalars, and there are no nonzero maps between
generators of different classes.  So a bounded complex $K$ of finite sums of
generators is isomorphic to a finite direct sum $\bigoplus_{[a]}V_{[a]}^{\bullet}
\otimes E_a$, where each $V_{[a]}^{\bullet}$ is a bounded complex of
finite-dimensional vector spaces and the differentials are scalar matrices.
Each $V^{\bullet}$ is homotopy equivalent to its cohomology by finite
dimensional linear algebra, with scalar homotopies, and hence
$K\simeq\bigoplus_{[a]}\bigoplus_nH^{n}(V_{[a]})\otimes E_a[-n]$.
\end{proof}

The proposition is used in three places: Proposition~\ref{prop:cat}(iv), where it
gives infinitely many pairwise non-isomorphic generators; the vanishing
\eqref{eq:HomUN} and the sentence after Proposition~\ref{prop:obstruction},
where it is what removes the self-duality obstruction from the category of
Definition~\ref{def:G}; and Theorem~\ref{thm:daggerthm}, the involution.

\begin{definition}\label{def:G}
$G$ is the differential graded category of bounded complexes of finite direct
sums of the bimodules $E_a$, with $a$ running over a fixed set of lifts of the
classes in $C_S$ (the choice is immaterial up to isomorphism, by
Proposition~\ref{prop:hom}), with differentials given by matrices
of intertwiners and with the complete differential graded $\Hom$ complex.
Shifts, finite sums and cones of objects are objects.
\end{definition}

\begin{corollary}\label{cor:trivial}
The section $\sigma$ of \eqref{eq:section} satisfies
$\beta_{\sigma(u)}\beta_{\sigma(v)}=\beta_{\sigma(uv)}$ exactly, not only up to
a gauge unitary.
Consequently the multiplication cocycle of the associated Fell bundle over
$C_S$ is strictly trivialized, every fibre acquires a canonical vector
$1_{E_u}$, and the assignment
\begin{equation}\label{eq:smear}
F\longmapsto E_F:=\int_{C_S}F(u)\,E_{\sigma(u)}\dstar u
\end{equation}
carries convolution on $\DA$ to $\star$ and $\dagger$ to the involution of
Section~\ref{sec:dagger}.  The coefficient of $E_F$ at the label $u$ is
exactly $F(u)$, and there is no coefficient at any nontrivial gauge label.
\end{corollary}

\begin{proof}
$\sigma$ is a group homomorphism by Lemma~\ref{thm:split}, and
$\beta_a\beta_b=\beta_{ab}$ holds for external ideles because they act by
multiplication; so the cocycle is trivial in the trivialization given by
$\sigma$.  With $\sigma$ fixed, the Fell product becomes ordinary convolution
on $C_S$ and the Fell involution becomes inversion, whence the two
compatibilities.  The statement about coefficients is the definition of the
integral \eqref{eq:smear} in the trivialization: without a splitting there is
no canonical vector $1_{E_u}$ in the fibre, and changing the idele lift sends
$1$ to a gauge multiplier $U_{\gamma}$; with $\sigma$ fixed, no such ambiguity
remains.
\end{proof}

\begin{remark}\label{rem:types}
The smeared objects $E_F$ are \emph{not} objects of $G$: an object of $G$ is a
finite sum of single twists $E_a$, and an integral over a continuum of classes
is not such a sum.  We therefore work with two categories, $G$ (single twists;
Sections~\ref{sec:cat}--\ref{sec:shadow}) and the category $\Corrtr$ of
complexes presented over $\DA$ (Section~\ref{sec:presented}), related by
$F\mapsto E_F$, which is a homomorphism of monoids but not an equivalence.
No result proved for one is transported to the other except along this map.
\end{remark}

\begin{proposition}\label{prop:cat}
$G$ is stable and closed under $\star$; it is geometrically localized on
$Y_S$; its composition is the continuous extension of \eqref{eq:starmult}; and
it contains infinitely many pairwise
non-isomorphic objects.
\end{proposition}

\begin{proof}
\emph{(i) Stability.}  Shifts, finite sums and cones of objects of $G$ are
objects of $G$: the cone differential matrix
$\left(\begin{smallmatrix}d_L&f\\0&-d_K\end{smallmatrix}\right)$ squares to
zero when $f$ is a chain map, and the differential graded $\Hom$ complex is
complete with all products finite, the objects being finite sums in each
degree.  The homotopy category of a differential graded category with cones
and shifts is triangulated with the cone triangles.

\emph{(ii) Closure under $\star$.}  On generators this is \eqref{eq:starmult}.
On complexes one totalizes term by term with the Koszul differential; the mixed
terms of $d^{2}$ cancel, and tensoring continuous intertwiners over $U$ is
balanced and continuous.  Hence $K\star L\in G$.

\emph{(iii) Geometric localization.}  By \eqref{eq:graphbimod}, $E_a$ is the
structure kernel of the graph of $m_{[a]}$ inside $Y_S$, that is the
pushforward of the source along an actual closed subgroupoid of the square,
and the two projections of $Y_S$ are the two distinct anchor maps.  Objects of
$G$ are finite complexes of finite sums of such, hence localized.

\emph{(iv) Infinitely many objects.}  By \eqref{eq:homclass},
$\Hom(E_a,E_b)\ne0$ if and only if $a\in\Gamma_S b$, and then it is one
dimensional; so $E_a\cong E_b$ if and only if $[a]=[b]$ in $C_S$.  The group
$C_S$ is infinite, since $|\cdot|_S$ maps it onto $\R_{>0}$ by
Lemma~\ref{thm:split}.  Hence $\{E_a\}_{a\in C_S}$ is an infinite family of
pairwise non-isomorphic objects; in particular $G$ is not the span of two
rulings.
\end{proof}

\begin{remark}[How much structure $G$ carries]\label{rem:semisimple}
By \eqref{eq:homclass} the endomorphism algebra of every generator is $\C$
and there are no morphisms between non-isomorphic generators.  Hence $G$ is
semisimple: it is equivalent to the category of bounded complexes of
finite-dimensional $C_S$-graded complex vector spaces, the grading recording
the class $[a]$ of each summand $E_a$, with $\star$ given on generators by
multiplication of classes, $\Delta$ the unit class, and $\dagger$ given by
conjugate transposition together with inversion of the grading.  The formal
properties of Proposition~\ref{prop:cat} and Theorem~\ref{thm:daggerthm} are
correct, and they are of little depth: all of the numerical content of the
theory lies in the analytic trace of Section~\ref{sec:trace} and not in $G$.
\end{remark}

\subsection{Two-sided excision}

\subsubsection*{The localized ideals}

For $J\subseteq S_f$ let
\begin{equation}\label{eq:UJ}
U_J:=\Bigl\{x\in\R^{\times}\times\prod_{p\in S_f}\Q_p:\ x_p\ne0\ \text{for every }p\in J\Bigr\},\qquad
I_J:=\Bigl(E_{\R}\,\widehat\otimes\!\!\bigotimes_{p\in J}F_p\,\widehat\otimes\!
\!\bigotimes_{p\in S_f\setminus J}E_p\Bigr)\rtimes_{\mathrm{rapid}}\Gamma_S ,
\end{equation}
a closed two-sided ideal of $U$ supported on $U_J$; $I_{\emptyset}=N$ is the
coefficient ideal of \eqref{eq:source}, not the unit.  Every $\Gamma_S$-orbit
and every external dilation preserves each $U_J$, hence each $I_J$.

The finite splittings $E_p=\C\chi_p\oplus F_p$ give a decomposition of the
coefficient algebra as a \emph{finite} direct sum of closed subspaces indexed
by the subsets $L\subseteq S_f$: the summand $C_L$ is the one carrying $F_p$ at
the places $p\in L$ and $\C\chi_p$ at the places $p\notin L$, tensored with
$E_{\R}$ and with the rapid lattice factor.  Adjoining the scalar lattice
summand $\C1\rtimes\Gamma_S$ one gets
\begin{equation}\label{eq:Ldecomp}
U=\Bigl(\bigoplus_{L\subseteq S_f}C_L\Bigr)\oplus\bigl(\C1
\rtimes_{\mathrm{rapid}}\Gamma_S\bigr),\qquad
I_J=\bigoplus_{L\supseteq J}C_L ,
\end{equation}
and the projection onto any set of summands is continuous in all the seminorms
of $U$.  In particular $0\to I_J\to U\to U/I_J\to0$ is admissible.  This is a
decomposition of $U$ as a topological vector space, not as a product of
algebras: multiplication mixes the summands, and only the containments in
\eqref{eq:Ldecomp} are used below.

\subsubsection*{The adapted partition}

\begin{lemma}\label{lem:partition}
Fix $J\subseteq S_f$.  There are families $(e_i)_{i\in\mathcal I}$ and
$(\rho_i)_{i\in\mathcal I}$ of elements of $I_J$, all placed at the identity
lattice label, such that
\begin{equation}\label{eq:partition}
e_i\rho_i=\rho_i\ \ \text{for every }i,\qquad
\sum_i\rho_i=1\ \text{ pointwise on }U_J ,
\end{equation}
the sum being locally finite on $U_J$, and such that for every pair of
continuous seminorms $\alpha,\beta$ on $I_J$ there are a continuous seminorm
$\nu$ on $U$ and a constant $C$ with
\begin{equation}\label{eq:summable}
\sum_i\alpha(e_i)\,\beta(\rho_i a)\ \le\ C\,\nu(a)\qquad(a\in I_J).
\end{equation}
Moreover $\sum_i\rho_ia=a$ in every seminorm of $I_J$, for every $a\in I_J$.
\end{lemma}

\begin{proof}
\emph{Construction.}  On the real line of the coordinate $t$ choose a
smooth nonnegative partition of unity $(\psi_n)_{n\in\Z}$ with
$\sum_n\psi_n=1$, each $\psi_n$ supported in the translate by $n$ of a fixed
compact interval, and with derivatives bounded uniformly in $n$; choose smooth
compactly supported $\varepsilon_n$ with $\varepsilon_n\equiv1$ on
$\supp\psi_n$ and with the same uniform derivative bounds; we use these functions on each
of the two half-lines of $\R^{\times}$, with $t=\log|x_{\infty}|$, so that the
real index runs over $\Z\times\{\pm\}$.  At a place $p\in J$
use the valuation-shell idempotents $1_{v_p=j}$, $j\in\Z$, which partition
$\Q_p^{\times}$; at a place $p\in S_f\setminus J$ use instead $\chi_p=1_{v_p\ge
0}$ --- which carries the value at $x_p=0$ --- together with the shell
idempotents $1_{v_p=j}$ for $j<0$, which partition all of $\Q_p$.  All of these
lie in the respective factor $F_p$ or $E_p$ and have trivial unit Fourier
character, i.e.\ $c_p$-conductor $1$.  Let $\mathcal I$ be the index set of all
products of one such function at each place together with one $\psi_n$, and let
$\rho_i$ be that product and $e_i$ the corresponding product with $\psi_n$
replaced by $\varepsilon_n$ and each shell idempotent by itself; place both at
the identity lattice label.  Then $e_i\rho_i=\rho_i$, because
$\varepsilon_n\psi_n=\psi_n$ and the idempotents are idempotent; and
$\sum_i\rho_i=1$ on $U_J$, locally finitely, because at each place the chosen
family partitions the relevant set and the $\psi_n$ partition $\R$.  Each
$\rho_i$ and each $e_i$ lies in $I_J$: at $p\in J$ its $p$-component is a
shell idempotent, which lies in $F_p$.

The use of $\chi_p$ at the places outside $J$ is essential: a nonzero value at
$x_p=0$ must not be decomposed into infinitely many shells of positive
valuation.

\emph{The estimate.}  Fix $\alpha,\beta$.  It suffices to bound
$\alpha(e_i)\beta(\rho_ia)$ for a single index $i=(n,(j_p)_p)$ and to sum a
geometric series.  In the seminorms \eqref{eq:seminorms} and the archimedean
seminorms $\sup_te^{l|t|}|D^{m}f|$, a real translate of the fixed cutoff
$\varepsilon_n$ has $\alpha$-size $O(e^{l_{\alpha}|n|})$, and a shell
idempotent $1_{v_p=j}$ has $v_{l,N}$-size $p^{l|j|}$ and unit conductor; all
real derivatives of the $\varepsilon_n$ and $\psi_n$ are bounded uniformly in
$n$ before these weights, and the $\psi_n$ and the idempotents carry the
identity lattice label and trivial unit character, so they do not disturb the
lattice or conductor weights.  Therefore, for a fixed pair $(\alpha,\beta)$,
choose $\nu$ with a real weight $l_{\nu}>l_{\alpha}+l_{\beta}+1$ and shell
weights $l_{\nu,p}>l_{\alpha,p}+l_{\beta,p}+1$ at every $p\in S_f$, and with
the same conductor and lattice exponents as $\beta$.  Multiplying $a$ by
$\rho_i$ and applying $\beta$, the localization of $a$ to the window
$(n,(j_p))$ is bounded by $\nu(a)$ times $e^{-|n|}\prod_pp^{-|j_p|}$, by the
definition of the weights; the Leibniz rule contributes only a finite sum of
derivative terms, each with uniformly bounded factors from $\psi_n$.
Multiplying by $\alpha(e_i)$ and summing over $n\in\Z$ and over the
finitely many places with $j_p\in\Z$ gives a convergent geometric sum, which is
\eqref{eq:summable}.  This is exactly where the rapid decay of the coefficient
algebra at the boundaries $x_p\to0$ and $x_p\to\infty$ is used.

\emph{Convergence of $\sum_i\rho_ia$.}  The same weighted tail estimate shows
that the partial sums of $\sum_i\rho_ia$ converge in every seminorm of $I_J$,
and the limit is $a$ because $\sum_i\rho_i=1$ pointwise on $U_J$ and $a$ is
supported there.  No claim is made that $\sum_i\rho_i$ converges to an element
of the nonunital ideal $I_J$.
\end{proof}

\begin{lemma}\label{lem:projective}
Each $I_J$ is relatively projective as a right $U$-module and as a left
$U$-module, and $I_J\widehat\otimes_U(-)$ preserves admissible exact sequences.
Moreover the bar complex of $I_J$ admits a continuous linear contraction, so
$I_J$ is admissibly H-unital.
\end{lemma}

\begin{proof}
Put $h_J(a):=\sum_ie_i\otimes\rho_ia$.  By \eqref{eq:summable} the sum
converges absolutely in the completed projective tensor product
$I_J\widehat\otimes I_J$ and $h_J$ is continuous.  Composed with multiplication
it is the identity, by \eqref{eq:partition} and the last assertion of
Lemma~\ref{lem:partition}.  Regarded as a map into
$I_J\widehat\otimes U$, it is right $U$-linear: associativity gives
$\rho_i(ab)=(\rho_ia)b$.  Its retraction, multiplication, is also right
$U$-linear.  Hence $I_J$ is a continuous right-module retract of the free right
module $I_J^{\flat}\widehat\otimes U$, $I_J^{\flat}$ denoting the underlying
Fr\'echet space of $I_J$.  A free module lifts through any admissible surjection
of right modules, by the underlying continuous linear section; a retract of a
free module therefore lifts as well.  So $I_J$ is relatively projective on the
right.  For the left, use the mirrored partition
$h_J^{\mathrm{op}}(a):=\sum_i a\rho_i\otimes e_i$, which satisfies the same
estimate and exhibits $I_J$ as a retract of the free left module
$U\widehat\otimes I_J^{\flat}$; the same argument applies.

For the preservation of admissible exact sequences: for a free module it
follows by tensoring the explicit linear splitting; for a retract, apply the
retract to the resulting split exact complex --- if $h$ is the contraction of
the split complex and $\pi$ the commuting chain projection onto the retract,
then $\pi h\pi$ contracts it.  This gives continuous split exactness, not
merely algebraic injectivity.

Finally, inserting $h_J$ at the first factor of the nonunital multiplication
bar complex $B'(I_J)$ contracts it: the first boundary term of
$b'h_J(a)$ is $\sum_ie_i\rho_ia=a$, and the remaining terms cancel $h_Jb'$ by
right linearity.  The contraction is continuous by \eqref{eq:summable}.
\end{proof}

\begin{lemma}\label{lem:idealproduct}
For $J,K\subseteq S_f$, multiplication induces a topological isomorphism
$I_J\widehat\otimes_UI_K\cong I_{J\cup K}$, and consequently
\begin{equation}\label{eq:idealprod}
I_J\,\widehat\otimes^{\mathbf L}_U\,I_K\cong I_{J\cup K}[0].
\end{equation}
\end{lemma}

\begin{proof}
\emph{Injectivity.}  Tensor the admissible sequence $0\to I_K\to U\to U/I_K\to0$
on the left with $I_J$, which is relatively projective on the right
(Lemma~\ref{lem:projective}).  This embeds $I_J\widehat\otimes_UI_K$ as a
closed subspace of $I_J\widehat\otimes_UU=I_J$, the embedding being
multiplication.  Its image lies in $I_{J\cup K}$: the support of a crossed
product is contained in $U_J\cap U_K=U_{J\cup K}$, the lattice preserving each
$U_K$.

\emph{Surjectivity.}  Apply $h_{J\cup K}$ of Lemma~\ref{lem:projective} to
$I_{J\cup K}$: each first factor $e_i$ lies in $I_J$ and each second factor
$\rho_ia$ lies in $I_K$, since both $e_i$ and $\rho_i$ are supported in
$U_{J\cup K}\subseteq U_J\cap U_K$.  So $h_{J\cup K}$ followed by the balanced
quotient is a continuous right inverse to multiplication onto $I_{J\cup K}$.
Together with injectivity this proves the isomorphism, including continuity of
the inverse.

Since $I_J$ is relatively projective on the right, the relative derived tensor
product has no higher terms, which is \eqref{eq:idealprod}.
\end{proof}

\begin{lemma}\label{lem:cech}
Let $J_1,\dots,J_m\subseteq S_f$ be a finite nonempty list, put
$O=\bigcup_{a}U_{J_a}$, $I_O=\sum_aI_{J_a}$ and $B_O=U/I_O$; for
$\emptyset\ne T\subseteq\{1,\dots,m\}$ put $J_T=\bigcup_{a\in T}J_a$.  Then:
\begin{enumerate}
\item[(i)] $I_O$ is a closed ideal, $U\to B_O$ is admissible, and $I_O$ is the
direct sum of those summands $C_L$ of \eqref{eq:Ldecomp} for which
$J_a\subseteq L$ for some $a$;
\item[(ii)] the augmented complex
$P_O^{0}=U$, $P_O^{-n}=\bigoplus_{|T|=n}I_{J_T}$ $(1\le n\le m)$, with the
alternating inclusions and the augmentation $U\to B_O$, is admissibly exact,
and is therefore a finite relative projective resolution of $B_O$ on either
side;
\item[(iii)] $I_{J_T}\widehat\otimes_UB_O=0$ for every $T\ne\emptyset$;
\item[(iv)] $B_O\widehat\otimes^{\mathbf L}_UB_O\cong B_O$.
\end{enumerate}
\end{lemma}

\begin{proof}
(i)  By \eqref{eq:Ldecomp}, $I_{J_a}=\bigoplus_{L\supseteq J_a}C_L$, so the sum
$I_O$ is the direct sum of the $C_L$ with $J_a\subseteq L$ for some $a$, a
finite direct sum of closed summands; the projection onto it is continuous in
all the seminorms of $U$, so $I_O$ is closed and the quotient map is
continuously linearly split.  The scalar lattice summand is not in $I_O$.

(ii)  The differential squares to zero because deleting two indices in the two
possible orders gives opposite signs.  For exactness, work on one summand $C_L$
and put $\mathcal A_L:=\{a:J_a\subseteq L\}$.  By \eqref{eq:Ldecomp}, the
summand of $P_O^{-n}$ at $C_L$ is $\bigoplus\{C_L:|T|=n,\ J_T\subseteq L\}$,
i.e.\ the span of the $n$-element subsets of $\mathcal A_L$.  If
$\mathcal A_L\ne\emptyset$, the complex at $C_L$ is the augmented simplicial
chain complex of the full simplex on the finite vertex set $\mathcal A_L$,
which contracts by insertion of the least vertex with the usual sign and zero
on repeated vertices.  If $\mathcal A_L=\emptyset$, only the identity
$C_L\to C_L$ from $P^0_O$ to $B_O$ survives, and the same holds on the scalar
lattice summand.  These contractions are continuous, being finite sums on each
summand, and they are combined by the continuous finite projections of
\eqref{eq:Ldecomp}.  They are not $U$-linear, which is exactly what
admissibility permits.  Every negative-degree term is relatively projective on
either side by Lemma~\ref{lem:projective}, and $P_O^{0}=U$ is free.

(iii)  Choose $a\in T$; then $J_a\subseteq J_T$, so multiplication
$I_{J_T}\widehat\otimes_UI_{J_a}\to I_{J_T}$ is surjective by
Lemma~\ref{lem:idealproduct}, with a continuous splitting.  After tensoring
with $B_O$ every balanced elementary product vanishes, because $I_{J_a}$
annihilates $B_O$; the continuous splitting propagates this to the whole
completed tensor product.

(iv)  Tensor the resolution (ii) with the left module $B_O$.  By (iii) every
negative-degree term dies, so the resulting complex is $B_O$ in degree zero and
zero elsewhere; hence the multiplication counit is a relative derived
isomorphism.
\end{proof}

Two-sided induction to the stratum is
$\mathcal I(M):=B_O\widehat\otimes^{\mathbf L}M\widehat\otimes^{\mathbf L}B_O$.
We abbreviate the two statements of Lemma~\ref{lem:cech}(iii),(iv) as
\begin{equation}\label{eq:cechvanish}
I_{J_T}\,\widehat\otimes_U\,B_O=0\quad(T\neq\emptyset),\qquad
B_O\,\widehat\otimes^{\mathbf L}_U\,B_O\cong B_O .
\end{equation}

\begin{theorem}\label{thm:exctriangle}
For bounded-above $M,N$ there is a distinguished triangle
\begin{equation}\label{eq:exctriangle}
B_O\widehat\otimes^{\mathbf L}M\widehat\otimes^{\mathbf L}I_O
\widehat\otimes^{\mathbf L}N\widehat\otimes^{\mathbf L}B_O
\longrightarrow \mathcal I(M\star N)\longrightarrow
\mathcal I(M)\star_{B_O}\mathcal I(N)\longrightarrow[1],
\end{equation}
so that $\mathcal I$ is monoidal on the pair $(M,N)$ if and only if the left
hand term vanishes.
\end{theorem}

\begin{proof}
Apply the exact functor
$B_O\widehat\otimes^{\mathbf L}M\widehat\otimes^{\mathbf L}(-)
\widehat\otimes^{\mathbf L}N\widehat\otimes^{\mathbf L}B_O$ to the admissible
triangle $I_O\to U\to B_O\to[1]$, and use the second identity of
\eqref{eq:cechvanish} on the right hand term.  Exactness of the functor is
K-flatness of the outer factors.
\end{proof}

\begin{lemma}\label{lem:excvanish}
$I_{K\cup J_T}\widehat\otimes_U B_O=0$ for every $T\ne\emptyset$ and every
$K\subseteq S_f$.
\end{lemma}

\begin{proof}
Choose $a\in T$; then $J_a\subseteq K\cup J_T$, so multiplication
$I_{K\cup J_T}\widehat\otimes_U I_{J_a}\to I_{K\cup J_T}$ is an isomorphism
with continuous splitting, by Lemmas~\ref{lem:partition} and
\ref{lem:idealproduct}.
Tensoring with $B_O$ kills every balanced product by the first identity of
\eqref{eq:cechvanish}.
\end{proof}

\begin{theorem}[Two-sided excision]\label{thm:excision}
For $M,N$ in the sector generated by $G$ and the localized ideals,
$\mathcal I(M\star N)\cong\mathcal I(M)\star_{B_O}\mathcal I(N)$.  In
particular composition, the unit and the restriction functors to every closed
stratum commute on $G$, coherently for nested strata.
\end{theorem}

\begin{proof}
Both sides of the vanishing condition in Theorem~\ref{thm:exctriangle} are
exact in each variable and commute with sums, shifts, cones and retracts, so it
suffices to test on generators.  For $M=E_u$, $N=E_v$ and a \v Cech term
$I_{J_T}$ of the truncated resolution of $I_O$, an invertible twist preserves
every coordinate open, so
$E_u\widehat\otimes I_{J_T}\widehat\otimes E_v\cong I_{J_T,uv}$, an ideal of
the same support, killed by Lemma~\ref{lem:excvanish}.  For localized ideals
the corresponding term is $I_{K\cup J_T\cup L}$ up to a twist, also killed.  A
finite complex all of whose terms vanish vanishes.
\end{proof}

\begin{remark}
The vanishing is not asserted for the profile objects $[I_J\to U]$, which are
not objects of $G$.
\end{remark}

\section{The involution}\label{sec:dagger}

We begin with an obstruction which explains why the involution is defined on
the category $G$ and not on a larger category containing the coefficient
ideal.  We have
\begin{equation}\label{eq:HomUN}
\Hom_{U\text{-}U}(U,N)=0 .
\end{equation}
Indeed, a bimodule map $U\to U$ is right multiplication by a central element
$z$ of $U$, that is, by an element satisfying \eqref{eq:intertw} with $a=b=1$;
by Proposition~\ref{prop:hom}, $z$ is a scalar.  A nonzero scalar does not lie
in $N$, so a bimodule map $U\to N$ is zero.

The vanishing obstructs self-duality of ideals, not of the diagonal.

\begin{proposition}\label{prop:obstruction}
The vanishing $\Hom_{U\text{-}U}(U,N)=0$ forbids the simultaneous existence of
a degree zero fully faithful duality $D$ with $D(I)\cong I$ for some nonzero
ideal $I\subseteq N$ and $D(\Delta)\cong\Delta$.  It is compatible with
$D(\Delta)\cong\Delta$ whenever no nonzero subobject of $N$ in the category is
self-dual.
\end{proposition}

\begin{proof}
Suppose $D(I)\cong I$ and $D(\Delta)\cong\Delta$ with $D$ fully faithful.
Dualizing the inclusion $I\hookrightarrow U$ produces a nonzero map
$U\cong D(\Delta)\to D(I)\cong I\subseteq N$, i.e.\ a nonzero element of
$\Hom(U,N)=0$, a contradiction.  If $D(I)\not\cong I$ the dual of the
inclusion is a map $U\to D(I)$, which is constrained only if $D(I)$ embeds
into $N$.
\end{proof}

The ideal $N$ is not an object of $G$, so the
obstruction does not apply to the category of Definition~\ref{def:G}, where
what is required is only $\Delta^{\dagger}\cong\Delta$.

Put $h:=|\cdot|_S^{-1/2}$, so that $h^{2}=|\cdot|_S^{-1}$ is the module of the
idele class action.  By Lemma~\ref{lem:prodformula} the gauge dilations are
measure preserving, so the external dilations act unitarily on $L^{2}(A_S)$ by
the half-density formula $U_af(x)=|a|_S^{-1/2}f(a^{-1}x)$, and the realized
graph operators satisfy
\begin{equation}\label{eq:graphops}
\Pi(b,a)=\pi(b)U_a=|a|_S^{-1/2}T_{b,a},\quad
\Pi(b,a)\Pi(c,d)=\Pi(b\beta_a(c),ad),\quad
\Pi(b,a)^{*}=\Pi(\beta_a^{-1}(b^{*}),a^{-1}),
\end{equation}
as bounded operators, the realization being faithful.

The function $h$ is a square root of the dualizing character.  Unlike a theta
characteristic on a curve, it involves no choice: the product formula fixes
it.

\begin{definition}\label{def:dagger}
$E_a^{\dagger}:=E_{a^{-1}}$, realized through \eqref{eq:graphops}.  A
nonzero morphism $v=\lambda U_{\gamma}:E_a\to E_b$ exists only when
$a=\gamma b$ (Proposition~\ref{prop:hom}); then $b^{-1}=\gamma a^{-1}$, and we
set
\[
v^{\dagger}:=\bar\lambda\,U_{\gamma}:\ E_b^{\dagger}=E_{b^{-1}}\longrightarrow
E_{a^{-1}}=E_a^{\dagger},
\]
which is again an element of \eqref{eq:homclass}, so that $\dagger$ is
contravariant, antilinear and involutive on morphisms.  (In earlier versions
$v^{\dagger}$ was written $\bar\lambda U_{\gamma^{-1}}$; that is the Hilbert
adjoint of $v$, a map $E_b\to E_a$, not a map $E_b^{\dagger}\to E_a^{\dagger}$.)
The definition is extended to complexes by conjugate transposition of
intertwiner matrices with the standard $\Hom$-complex signs.
\end{definition}

\begin{theorem}\label{thm:daggerthm}
Then $\dagger$ is a continuous, antilinear,
contravariant involution on $G$ with $(K\star L)^{\dagger}\cong
L^{\dagger}\star K^{\dagger}$ and $\Delta^{\dagger}=\Delta$, defined from the
source alone.
\end{theorem}

\begin{proof}
On generators $\dagger$ is $a\mapsto a^{-1}$, and on morphisms it is the
rule of Definition~\ref{def:dagger}, which is contravariant and involutive.  By
\eqref{eq:homclass} the morphisms are finite scalar matrices times lift
unitaries, so conjugate transposition is well defined and continuous for the
finite-product topology on the $\Hom$ complex, and it commutes with the
differential up to the usual sign; it therefore passes to cones and shifts.
Antimultiplicativity is the second identity of \eqref{eq:graphops} together
with \eqref{eq:starmult}, and $\Delta=E_1$ is fixed because $1^{-1}=1$.
\end{proof}

Under the trivialization of Corollary~\ref{cor:trivial} the involution
$\dagger$ on smeared objects is $F\mapsto F^{\dagger}$ of
Definition~\ref{def:DA}.

\section{The Hochschild shadow of a graph bimodule}\label{sec:shadow}

In this section we compute the relative Hochschild shadow \eqref{eq:shadow} of
a graph bimodule.  The conclusion is that no additive invariant of these
complexes can serve as an intersection number; this is why the pairing of
Section~\ref{sec:trace} is analytic.

\subsection{Excision at the origin and the K\"unneth decomposition}

\begin{lemma}[H-unitality]\label{lem:hunital}
$E_P$, every $I_J$, and $N$ are admissibly H-unital: their bar complexes
$B'_n(I)=I^{\widehat\otimes n}$ with
$b'(a_1\otimes\cdots\otimes a_n)=\sum_j(-1)^{j-1}a_1\otimes\cdots\otimes
a_ja_{j+1}\otimes\cdots\otimes a_n$ admit a continuous linear contraction $h$
with $b'h+hb'=\id$ in degrees $n\ge1$.
\end{lemma}

\begin{proof}
For the ideals $I_J$, and in particular for $N=I_{\emptyset}$, this is the
last assertion of Lemma~\ref{lem:projective}, with the contraction
$h_J(a)=\sum_ie_i\otimes\rho_ia$ of Lemma~\ref{lem:partition}.  For $E_P$ one
runs the same construction on the coefficient algebra alone, with the
partition of Lemma~\ref{lem:partition} for $J=\emptyset$ (that is, $\chi_p$
together with the shells of negative valuation at every finite place) taken
without the lattice factor.  Its members lie in $E_P$ and sum to $1$ on all of
$\R^{\times}\times\prod_p\Q_p$, so $\sum_i\rho_ia=a$ for every $a\in E_P$;
and \eqref{eq:summable} holds verbatim, its proof never using the lattice
weights.
\end{proof}

\begin{theorem}[Split augmented excision]\label{thm:B2}
Let $R=\C\cdot1\oplus I$ be a unital augmented complete Fr\'echet algebra with
augmentation ideal $I$ admissibly H-unital, and let $\sigma$ be a continuous
unital endomorphism, $\sigma(1)=1$, with $\sigma(I)\subseteq I$ (an
automorphism with $\sigma(I)=I$ is the main case).  The twisted bimodule
$R_{\sigma}$ is $R$ with $a\cdot m\cdot b=\sigma(a)mb$; it is unital because
$\sigma(1)=1$, so the normalized complex computes its Hochschild homology.  Then there is a continuous
homotopy equivalence of complexes
\[
\bar C(R,R_{\sigma})\ \simeq\ \C[0]\ \oplus\ C(I,I_{\sigma}),
\]
where $\bar C$ is the normalized Hochschild complex and
$C(I,I_{\sigma})_n=I\widehat\otimes I^{\widehat\otimes n}$ with the twisted
Hochschild boundary.  In particular
$\HH(R,R_{\sigma})\cong\C[0]\oplus\HH(I,I_{\sigma})$ admissibly.
\end{theorem}

\begin{proof}
The normalized complex is
$\bar C_n=R_{\sigma}\widehat\otimes(R/\C1)^{\widehat\otimes n}
=R\widehat\otimes I^{\widehat\otimes n}
=(\C1\widehat\otimes I^{\widehat\otimes n})\oplus
(I\widehat\otimes I^{\widehat\otimes n})=:P_n\oplus Q_n$, a topological direct
sum because $R=\C1\oplus I$ is one.  The twisted boundary is
\[
b(a_0\otimes\cdots\otimes a_n)=\sum_{j<n}(-1)^{j}a_0\otimes\cdots\otimes
a_ja_{j+1}\otimes\cdots+(-1)^{n}\sigma(a_n)a_0\otimes a_1\otimes\cdots\otimes
a_{n-1}.
\]
On $Q$ all products lie in $I$ and $\sigma(I)\subseteq I$, so $b(Q_n)\subseteq Q_{n-1}$ and
$Q=C(I,I_{\sigma})$.  On $P$, applied to $1\otimes a_1\otimes\cdots\otimes
a_n$, the terms with $j\ge1$ are
$(-1)^{j}1\otimes\cdots\otimes a_ja_{j+1}\otimes\cdots\in P_{n-1}$ and
reproduce $-b'$ on $I^{\widehat\otimes n}$; the term $j=0$ is
$a_1\otimes\cdots\otimes a_n\in Q_{n-1}$; and the cyclic term is
$(-1)^{n}\sigma(a_n)\otimes a_1\otimes\cdots\otimes a_{n-1}\in Q_{n-1}$.  Hence
$P_{\bullet}$ is the bar complex $B'(I)$ with differential $-b'$, and there is
a connecting map $\phi:P_n\to Q_{n-1}$,
$\phi(a_1\otimes\cdots\otimes a_n)=a_1\otimes\cdots\otimes a_n+
(-1)^{n}\sigma(a_n)\otimes a_1\otimes\cdots\otimes a_{n-1}$, so that
$\bar C=\Cone(\phi)$ up to the standard sign convention.  In degree zero
$P_0=\C\cdot1$ carries $b=0$ and splits off as $\C[0]$, because
$b(1\otimes a_1)=a_1-\sigma(a_1)\in I$ lands in $Q_0$.  Finally $B'(I)$ is
contractible by $h$ (Lemma~\ref{lem:hunital}); for a cone $\Cone(\phi:X\to Y)$
with $X$ contractible via $h$, the inclusion $Y\to\Cone(\phi)$ is a homotopy
equivalence with inverse $(x,y)\mapsto y-\phi(hx)$ and homotopy built from $h$,
every map being a finite composite of $\phi$, $h$ and the differentials, hence
continuous.
\end{proof}

Applied to $R=E_P\oplus\C1$ with $I=E_P$, Theorem~\ref{thm:B2} gives
$\HH(R,R_{\sigma})\cong\C[0]\oplus\HH(E_P,(E_P)_{\sigma})$, the summand $\C[0]$
being the class of the adjoined unit.  Applied at a single finite place to the
unital algebra $A_1:=\C\chi_p\oplus F_p^{\ge0}$ of functions on $\Z_p$, with
augmentation the evaluation at the origin and ideal $F_p^{\ge0}$, it gives
$\HH(A_1,\sigma)\cong\C[0]\oplus\HH(F_p^{\ge0},\sigma)$, the summand $\C[0]$
being the origin class; it is untwisted because every dilation fixes $0$.

\begin{theorem}[Excision; cited for $\sigma=\id$]\label{thm:excision-input}
Let $A$ be a unital complete Fr\'echet algebra, $I\subseteq A$ a closed
two-sided ideal with a continuous linear complement, and suppose $I$ is
admissibly H-unital (its bar complex $B'(I)$ has a continuous linear
contraction).  Then the inclusion $C(I,I)\to\ker\bigl(C(A,A)\to
C(A/I,A/I)\bigr)$ is a quasi-isomorphism, so that
$\HH(I)\to\HH(A)\to\HH(A/I)$ is a long exact sequence.
\end{theorem}

This is Wodzicki's theorem \cite[Theorem~3.1]{Wodzicki} for algebras over a
field of characteristic zero, without topology; its version for Fr\'echet
algebras and continuous Hochschild homology, under exactly the hypothesis that
the extension has a continuous linear section and $B'(I)$ a continuous linear
contraction, is the main theorem of Brodzki and Lykova \cite{BrodzkiLykova}.  For bornological frameworks see Cuntz--Quillen
\cite{CuntzQuillen} and Meyer \cite{MeyerCyclic}.
Both sources treat the coefficients $A$, $I$, $A/I$, that is, $\sigma=\id$.
We need the same statement with the twisted coefficients $A_{\sigma}$,
$I_{\sigma}$, $(A/I)_{\sigma}$ for a continuous automorphism $\sigma$ of $A$
with $\sigma(I)=I$.  We have not found this twisted form stated in the
literature, and we do not prove it here; we record it as
\begin{equation}\label{eq:Excsigma}
(\mathrm{Exc}_{\sigma})\qquad
C(I,I_{\sigma})\xrightarrow{\ \simeq_{\mathrm{qis}}\ }
\ker\bigl(C(A,A_{\sigma})\to C(A/I,(A/I)_{\sigma})\bigr),
\end{equation}
and every statement below that depends on it says so.  The local computations
of Lemmas~\ref{lem:block}, \ref{lem:nonunit} and \ref{lem:unit} do not use it.

\begin{theorem}[K\"unneth over places; conditional on $(\mathrm{Exc}_{\sigma})$]\label{thm:kunneth}
Let $\sigma=\bigotimes_v\sigma_v$ act coordinatewise on $E_P$, and assume
$(\mathrm{Exc}_{\sigma})$ for the ideals named in Step~3 below.  Then
\begin{equation}\label{eq:kunneth}
\HH\bigl(E_P,(E_P)_{\sigma}\bigr)\ \cong\
H\Bigl(\mathrm{Tot}\,\widehat\bigotimes_{v\in S}C\bigl(E_v,(E_v)_{\sigma_v}\bigr)\Bigr),
\end{equation}
the isomorphism being induced by a zigzag of continuous chain maps which are
quasi-isomorphisms; here $E_{\infty}:=E_{\R}$.
\end{theorem}

\begin{proof}
\emph{Step 1 (unital Eilenberg--Zilber).}  For unital complete Fr\'echet
algebras $A,B$ with complete unital bimodules $M,N$, the Hochschild complex
$C(A\widehat\otimes B,M\widehat\otimes N)$ has $n$-th term
$(M\widehat\otimes N)\widehat\otimes(A\widehat\otimes B)^{\widehat\otimes n}
\cong(M\widehat\otimes A^{\widehat\otimes n})\widehat\otimes
(N\widehat\otimes B^{\widehat\otimes n})$, a continuous isomorphism by
associativity and commutativity of the completed projective tensor product,
with face maps $d_i\otimes d_i$: it is the diagonal of the bisimplicial module
$C_{p,q}=C_p(A,M)\widehat\otimes C_q(B,N)$, whose degeneracies exist because
$A$ and $B$ are unital.  The shuffle map $\nabla$ and the Alexander--Whitney
map $AW$ are chain maps with $AW\circ\nabla=\id$ and $\nabla\circ AW\simeq\id$
by the Eilenberg--Mac Lane homotopy; these are the classical combinatorial
identities among faces and degeneracies \cite[VIII.8]{MacLane}, valid in any
additive category.  Each face is a multiplication or a module action and each
degeneracy is a unit insertion, hence continuous; and every map above is in
each degree a finite signed sum of finite composites of them.  Therefore
$C(A\widehat\otimes B,M\widehat\otimes N)\simeq
\mathrm{Tot}(C(A,M)\widehat\otimes C(B,N))$ continuously and admissibly.  For
$\sigma=\sigma_A\otimes\sigma_B$ one applies this with $M=A_{\sigma_A}$ and
$N=B_{\sigma_B}$, and then induces on the number of factors.

\emph{Step 2 (support grading).}  Let $E_v^{+}=E_v\oplus\C1$ and
$\mathcal C^{\otimes}:=\widehat\bigotimes_{v\in S}E_v^{+}$.  As a topological
vector space $\mathcal C^{\otimes}=\bigoplus_{V\subseteq S}E_V$ with
$E_V=\widehat\bigotimes_{v\in V}E_v$ (units at the places outside $V$), and
$E_V\cdot E_W\subseteq E_{V\cup W}$.  Call $V$ the \emph{support} of an element
of $E_V$, and the \emph{support} of an elementary tensor of homogeneous letters
the union of the supports of its letters.  Every face of the Hochschild
complex multiplies two adjacent letters (or lets the last letter act on the
first), so it either kills a word or preserves its support; the twist
$\sigma=\bigotimes\sigma_v$ preserves each $E_V$; and the degeneracies insert
the unit, of support $\emptyset$.  Hence $C(\mathcal C^{\otimes},
\mathcal C^{\otimes}_{\sigma})=\bigoplus_{U\subseteq S}C^{(U)}$ as complexes, the
projections being continuous, and all maps of Step~1 (built from faces and
degeneracies) preserve the grading.  The same applies to each $E_v^{+}$, with
supports $\emptyset$ and $\{v\}$, and in the proof of Theorem~\ref{thm:B2} the
summand $\C[0]$ is the support-$\emptyset$ part while $Q=C(E_v,(E_v)_{\sigma_v})$
and the contractible complex $P_{\ge1}$ have support $\{v\}$; so
$C(E_v^{+})^{(\{v\})}\simeq C(E_v,(E_v)_{\sigma_v})$.  Taking the support-$S$
part of Step~1,
\begin{equation}\label{eq:topsupport}
C(\mathcal C^{\otimes},\mathcal C^{\otimes}_{\sigma})^{(S)}\ \simeq\
\mathrm{Tot}\,\widehat\bigotimes_{v\in S}C\bigl(E_v^{+}\bigr)^{(\{v\})}
\ \simeq\ \mathrm{Tot}\,\widehat\bigotimes_{v\in S}C\bigl(E_v,(E_v)_{\sigma_v}\bigr).
\end{equation}

\emph{Step 3 (the top support is $E_P$).}  It remains to identify
$C(E_P,(E_P)_{\sigma})$, which lies in support $S$, with the support-$S$ part
of $C(\mathcal C^{\otimes})$ up to quasi-isomorphism.  We argue by induction
on $|S|$; for $|S|=1$ this is the support-$\{v\}$ statement of Step~2.  For
$|S|\ge2$ fix $v\in S$, put $I:=E_v$, $J:=E_{S\setminus\{v\}}$, and let
$T:=I^{+}\widehat\otimes J^{+}=\C\oplus I\oplus J\oplus K$ with
$K:=I\widehat\otimes J=E_P$.  The grading by supports $\emptyset,\{1\},\{2\},
\{1,2\}$ of the letters $\C,I,J,K$ is preserved by the Hochschild boundary as
in Step~2.  Each of $I$, $J$ and $K$ is admissibly H-unital: the proof of
Lemma~\ref{lem:hunital} applies verbatim to $E_V$ for every $V$.

(i) $K\subseteq T$ is a closed ideal with continuous complement
$\C\oplus I\oplus J$, and $T/K=:Q'=\C\oplus I\oplus J$ with $IJ=JI=0$.  By
$(\mathrm{Exc}_{\sigma})$, $C(K,K_{\sigma})\to\ker(C(T,T_{\sigma})\to
C(Q',Q'_{\sigma}))$ is a quasi-isomorphism.  The quotient map preserves the
support grading, and for a word of support $\ne\{1,2\}$ no letter lies in $K$
and no product of letters reaches $K$, so it is an isomorphism in supports
$\ne\{1,2\}$.  Hence the kernel is concentrated in support $\{1,2\}$, and there
is a long exact sequence
\[
\cdots\to\HH(K,K_{\sigma})\to\HH(T,T_{\sigma})^{(12)}\to
\HH(Q',Q'_{\sigma})^{(12)}\to\cdots .
\]

(ii) $\HH(Q',Q'_{\sigma})^{(12)}=0$.  Indeed $I\subseteq Q'$ is a closed ideal,
admissibly H-unital, with quotient $Q'/I=J^{+}$.  By $(\mathrm{Exc}_{\sigma})$
again, and since inclusion and quotient preserve the support grading of
$C(Q')$ (letters $\C,I,J$ of supports $\emptyset,\{1\},\{2\}$; products of an
$I$-letter with a $J$-letter vanish), we get a long exact sequence in support
$\{1,2\}$ whose outer terms are $\HH(I,I_{\sigma})^{(12)}=0$ (every word of
$C(I)$ has support $\{1\}$) and $\HH(J^{+})^{(12)}=0$ (every word of $C(J^{+})$
has support $\subseteq\{2\}$).

(iii) Combining, $\HH(K,K_{\sigma})\cong\HH(T,T_{\sigma})^{(12)}$, induced by
the inclusion.  By Step~1 applied to $T=I^{+}\widehat\otimes J^{+}$ and the
grading, $C(T)^{(12)}\simeq\mathrm{Tot}(C(I^{+})^{(\{1\})}\widehat\otimes
C(J^{+})^{(\{2\})})$, and $C(I^{+})^{(\{1\})}\simeq C(I,I_{\sigma_v})$ and
$C(J^{+})^{(\{2\})}\simeq C(J,J_{\sigma'})$, $\sigma'=\bigotimes_{u\ne
v}\sigma_u$, both by Theorem~\ref{thm:B2} read in the support grading as in
Step~2.  By the induction hypothesis $C(J,J_{\sigma'})$ is quasi-isomorphic to
$\mathrm{Tot}\widehat\bigotimes_{u\ne v}C(E_u,(E_u)_{\sigma_u})$.  This proves
\eqref{eq:kunneth}.  The filtration used is the two-step chain of ideals
$K\subset I\oplus K\subset T$, read off the support grading; both steps are
ideals because $E_V\cdot E_W\subseteq E_{V\cup W}$.
\end{proof}

\begin{remark}\label{rem:kunnethstatus}
Only Step~3 uses $(\mathrm{Exc}_{\sigma})$, twice, for the ideals $K\subset T$
and $I\subset Q'$.  Steps~1--2 and Theorem~\ref{thm:B2} are proved.  The
consequences of Theorem~\ref{thm:kunneth} below (Theorem~\ref{thm:shadow}(b)
and the global form of Corollary~\ref{cor:blind}) inherit this condition; the
local statements, Theorem~\ref{thm:shadow}(c), do not.
\end{remark}

\subsection{The local twisted computations}

Fix a finite place $p$.  Write $F^{(j)}$ for the algebra of functions on the
shell $p^{j}\Z_p^{\times}$ with rapid conductor weights, $1_j$ for its unit,
$F_p^{\ge0}=\widehat\bigoplus_{j\ge0}F^{(j)}$, $F_p^{<0}=\widehat\bigoplus_{j<0}
F^{(j)}$, and $A_*:=\C\chi_p\oplus F_p^{\ge0}=\chi_pE_p$, the algebra of
functions on $\Z_p$, unital with unit $\chi_p$.  Then $E_p=A_*\oplus
\widehat\bigoplus_{j<0}F^{(j)}$ is a rapid orthogonal sum of unital algebras.
For an endomorphism $\rho$ we use the bimodule $A_{\rho}$ of
Theorem~\ref{thm:B2}, $a\cdot m\cdot b=\rho(a)mb$.

\begin{lemma}[Block lemma]\label{lem:block}
Let $A=\widehat\bigoplus_{a\in\mathcal I}A_a$ be a complete Fr\'echet algebra
which is a rapid orthogonal sum of unital closed ideals $A_a=e_aA$, with
orthogonal idempotents $e_a$ such that $x=\sum_ae_ax$ for all $x\in A$, and
such that for every continuous seminorm $\nu$ there is a continuous seminorm
$\nu'$ with $\sum_a\nu(e_a)\,\nu(e_ax)\le\nu'(x)$.  Let $M$ be a complete
$A$-bimodule with $m=\sum_{a,b}e_a\cdot m\cdot e_b$, satisfying the same
estimate, and put $M_{ab}:=e_a\cdot M\cdot e_b$.  Then
\[
C(A,M)\ =\ \widehat\bigoplus_{a}C(A_a,M_{aa})\ \oplus\ X
\]
as complexes, with $X$ admissibly contractible.  In particular, if
$M_{aa}=0$ for every $a$, then $C(A,M)$ is admissibly contractible.
\end{lemma}

\begin{proof}
By the hypotheses every chain is a rapidly convergent sum of
\emph{homogeneous} words $w=m\otimes a_1\otimes\cdots\otimes a_n$ with
$m\in M_{ab}$ and $a_r\in A_{k_r}$, and the projections onto multi-indices are
continuous.  Put $k_0:=b$ and $k_{n+1}:=a$.  Say that the edge $r$
($0\le r\le n$) of $w$ is \emph{matched} if $k_r=k_{r+1}$.  The face $d_r$
multiplies across edge $r$ (for $r=n$ it is the action $a_n\cdot m$), and it
vanishes on $w$ if the edge is unmatched, because then the product contains
$e_{k_r}e_{k_{r+1}}=0$.  A matched face merges two letters and leaves the
remaining edges, and whether they are matched, unchanged.  Therefore the
fully matched words, i.e.\ those with all indices equal to some $a$, span the
subcomplex $\widehat\bigoplus_aC(A_a,M_{aa})$, and the words with at least one
unmatched edge span a complementary subcomplex $X$.

On $X$ let $r(w)$ be the first unmatched edge and define
$s(w):=(-1)^{r}s_{r}(w)$, where $s_r$ inserts the unit $e_{k_r}$ of
$A_{k_r}$ immediately after the $r$-th slot (after $m$ if $r=0$).  Then
$d_rs_rw=w$, because $e_{k_r}$ is a unit for the $r$-th letter;
$d_{r+1}s_rw=0$, because $e_{k_r}$ meets the next letter (or, for $r=n$, acts
on $m$) across an unmatched edge; for $t<r$ one has $d_ts_r=s_{r-1}d_t$ and
$d_tw$ has first unmatched edge $r-1$, with the same unit to insert; and for
$t>r+1$ one has $d_ts_r=s_rd_{t-1}$, while $d_{t-1}w$ is either $0$ or has
first unmatched edge $r$.  Also $d_rw=0$.  Summing with the Hochschild signs,
all terms of $bs+sb$ cancel in pairs except $(-1)^{r}(-1)^{r}d_rs_rw=w$, so
$bs+sb=\id$ on $X$, in every degree including $0$ (there $w=m\in M_{ab}$ with
$a\ne b$, and $b(m\otimes e_b)=m\cdot e_b-e_b\cdot m=m$).  The estimate in the
hypothesis makes $s$ continuous, since inserting $e_{k_r}$ next to a letter of
$A_{k_r}$ costs at most the factor $\nu(e_{k_r})$ in a projective seminorm.
\end{proof}

For the shell decomposition of $E_p$ the estimate holds: with the seminorms
\eqref{eq:seminorms} one has $v_{l,N}(1_j)\,v_{l,N}(1_jx)\le
v_{2l,N}(1_jx)$ (the trivial character carrying the weight $1$), and the block
$A_*$ is a single summand.  Lemma~\ref{lem:block} replaces the orthogonal-sum
argument of earlier versions, which used a homotopy that does not satisfy
$bs+sb=\id$ on all mixed words.

\begin{lemma}[Non-unit twists]\label{lem:nonunit}
Let $\sigma$ be the dilation $f\mapsto f(w^{-1}\,\cdot\,)$ by
$w\in\Q_p^{\times}$ with $v:=v_p(w)\ne0$.  Then
\begin{equation}\label{eq:D1}
C\bigl(E_p,(E_p)_{\sigma}\bigr)\ \simeq\ \C[0]
\end{equation}
admissibly, the class being that of $\chi_p$, i.e.\ of the origin $x_p=0$, and
independent of $w$.  Moreover $C(F_p,(F_p)_{\sigma})$ is admissibly
contractible.
\end{lemma}

\begin{proof}
$\sigma$ maps the shell $j$ to the shell $j+v$, so $\sigma(1_j)=1_{j+v}$ and
$\sigma(\chi_p)=\mathbf 1_{w\Z_p}$; $\sigma$ preserves $E_p$ and $F_p$ but not
$A_*$.  Apply Lemma~\ref{lem:block} to $A=E_p$ with the blocks $A_*$ and
$F^{(j)}$, $j<0$, and $M=(E_p)_{\sigma}$, for which
$M_{ab}=\sigma(e_a)E_pe_b$.  For a shell block,
$M_{jj}=1_{j+v}1_jE_p=0$.  Hence $C(E_p,M)\simeq C(A_*,M_{**})$ with
$M_{**}=\mathbf 1_{w\Z_p}\chi_pE_p$.

If $v<0$, then $\Z_p\subset w\Z_p$ and $M_{**}=A_*$ with $a\cdot m\cdot
b=\sigma'(a)mb$, where $\sigma'(a):=\sigma(a)|_{\Z_p}$, i.e.\
$\sigma'(a)(x)=a(w^{-1}x)$ for $x\in\Z_p$; since $w^{-1}\Z_p\subseteq\Z_p$,
$\sigma'$ is a continuous unital endomorphism of $A_*$ fixing the value at the
origin, so $\sigma'(F_p^{\ge0})\subseteq F_p^{\ge0}$.  If $v>0$, then
$M_{**}=\sigma(A_*)$ consists of the functions on $w\Z_p$, and
$\psi:A_*\to M_{**}$, $\psi(c)=\sigma(c)$, is an isomorphism with
$a\cdot\psi(c)\cdot b=\psi\bigl(ac\,\rho(b)\bigr)$, $\rho(b)(y):=b(wy)$; so
$M_{**}$ is $A_*$ twisted on the right by the unital endomorphism $\rho$, which
again fixes the value at the origin.  Since $A_*$ is commutative, reversing
the order of the letters, $m\otimes a_1\otimes\cdots\otimes a_n\mapsto
\pm\,m\otimes a_n\otimes\cdots\otimes a_1$, identifies the complex of a
right-twisted bimodule with that of the left-twisted one.

In both cases Theorem~\ref{thm:B2} applies to $R=A_*=\C\chi_p\oplus
F_p^{\ge0}$ with augmentation the value at the origin, and gives
$C(A_*,M_{**})\simeq\C[0]\oplus C(F_p^{\ge0},(F_p^{\ge0})_{\rho'})$ with
$\rho'=\sigma'$ or $\rho$, the $\C[0]$ being the class of $\chi_p$.  The
endomorphism $\rho'$ sends $1_j$ to $1_{j-v}$ (or to $0$ when that shell is not
in $\Z_p$), never to $1_j$, so Lemma~\ref{lem:block} applied to the shells
$j\ge0$ shows that $C(F_p^{\ge0},(F_p^{\ge0})_{\rho'})$ is admissibly
contractible.  The last assertion is Lemma~\ref{lem:block} applied directly to
$F_p=\widehat\bigoplus_{j\in\Z}F^{(j)}$, where $M_{jj}=1_{j+v}1_jF_p=0$.
\end{proof}

\begin{lemma}[Unit twists]\label{lem:unit}
Let $\sigma_p$ be the dilation by a unit $u\in\Z_p^{\times}$, $u\ne1$, and let
$H_u:=\overline{\langle u\rangle}\subseteq\Z_p^{\times}$ be the closed subgroup
it generates.  Then for every shell $j$:
\begin{enumerate}
\item[(i)] the closure of $(1-\sigma)F^{(j)}$ is the closed span of the
characters $e_{\kappa}$ with $\kappa(u)\ne1$, and the \emph{reduced} group
$\HH_0^{\mathrm{red}}:=F^{(j)}/\overline{(1-\sigma)F^{(j)}}$ is the space of
functions in $F^{(j)}$ invariant under $H_u$, i.e.\ the closed span of the
characters trivial on $H_u$;
\item[(ii)] if $u$ has infinite order, $(1-\sigma)F^{(j)}$ is not closed, so
the unreduced $\HH_0$ is not Hausdorff; if $u$ has finite order it is closed
and $\HH_0=\HH_0^{\mathrm{red}}$;
\item[(iii)] $\dim\HH_0^{\mathrm{red}}=[\Z_p^{\times}:H_u]$, which equals
\[
N(u)=\begin{cases}(p-1)p^{k-1}, & u\in1+q_p\Z_p,\ k:=v_p(1-u),\\
2^{k'-1}, & p=2,\ u\in3+4\Z_2,\ u\ne-1,\ k':=v_2(1+u)\ge2,\\
\infty, & p=2,\ u=-1,
\end{cases}
\]
with $q_p=p$ for $p$ odd and $q_2=4$; for $p$ odd and $u\notin1+p\Z_p$ the
index $[\Z_p^{\times}:H_u]$ is again finite or infinite according as $u$ has
infinite or finite order, and it is not determined by $|1-u|_p$
(Remark~\ref{rem:closure}).
\end{enumerate}
Moreover $C(E_p,(E_p)_u)\simeq\C[0]\oplus\widehat\bigoplus_{j\in\Z}
C(F^{(j)},(F^{(j)})_u)$ admissibly, so that
\begin{equation}\label{eq:D2}
\HH_0^{\mathrm{red}}\bigl(E_p,(E_p)_{u}\bigr)\ \cong\ \C\ \oplus\
\widehat\bigoplus_{j\in\Z}\C^{N(u)},
\end{equation}
of infinite dimension in every case: the origin class together with the
$H_u$-invariant functions on each valuation shell.
\end{lemma}

\begin{proof}
$\sigma$ preserves every shell and acts on $F^{(j)}\cong
C^{\infty}_{\mathrm{rapid}}(\Z_p^{\times})$ by $x\mapsto u^{-1}x$; in the
character basis $\sigma(e_{\kappa})=\kappa(u)^{\pm1}e_{\kappa}$, and we write
$\kappa(u)$ for the eigenvalue.  Since $F^{(j)}$ is unital and commutative,
$\HH_0(F,F_{\sigma})=F/(1-\sigma)F$ with
$(1-\sigma)e_{\kappa}=(1-\kappa(u))e_{\kappa}$.

(i) If $\kappa(u)\ne1$ then $e_{\kappa}=(1-\sigma)e_{\kappa}/(1-\kappa(u))$ lies
in the image, so the closure of the image contains the closed span of these
characters; conversely the coefficient of $e_{\kappa}$ with $\kappa(u)=1$ is a
continuous functional vanishing on the image.  A character is trivial at $u$
if and only if it is trivial on $H_u$, by continuity.

(ii) If $u$ has finite order $e$, every $\kappa(u)\ne1$ is a nontrivial $e$-th
root of unity, so $|1-\kappa(u)|\ge|1-e^{2\pi i/e}|$ and division by
$1-\kappa(u)$ is continuous: the image is closed.  If $u$ has infinite order,
write $u=\zeta u_1$ with $\zeta$ in the torsion subgroup $\mu$ of
$\Z_p^{\times}$ and $1\ne u_1\in1+q_p\Z_p$, $k:=v_p(1-u_1)$, and consider the
characters $\kappa$ trivial on $\mu$.  For such $\kappa$ of conductor exponent
$m>k$ (and $m\ge3$ if $p=2$) the value $\kappa(u)=\kappa(u_1)$ is a primitive
root of unity of order $p^{m-k}$, by the logarithm isomorphism
$1+q_p\Z_p\cong q_p\Z_p$, so $|1-\kappa(u)|\asymp p^{-(m-k)}$.
The seminorms \eqref{eq:seminorms} weight the conductor only polynomially,
by $c_p(\kappa)^{N}$; choosing one such $\kappa_m$ for each $m$ and
$g:=\sum_m e^{-\sqrt m}e_{\kappa_m}\in\overline{(1-\sigma)F}$, any preimage
would have coefficients of size $e^{-\sqrt m}|1-\kappa_m(u)|^{-1}$, which grow
exponentially, so $g$ is not in the image.

(iii) The $H_u$-invariant functions are the functions on the finite or
profinite group $\Z_p^{\times}/H_u$, so the dimension is the index.  For
$u\in1+q_p\Z_p$ the $p$-adic logarithm identifies $1+q_p\Z_p$ with
$q_p\Z_p$ and $u^{\Z_p}$ with $p^{k}\Z_p$, so $H_u=1+p^{k}\Z_p$ and the index is
$|\Z_p^{\times}/(1+p^{k}\Z_p)|=(p-1)p^{k-1}$.  For $p=2$ and $u\in3+4\Z_2$,
$u^{2}\in1+8\Z_2$ has $v_2(1-u^{2})=1+k'$, so
$\overline{\langle u^{2}\rangle}=1+2^{k'+1}\Z_2$ and $H_u$ is the union of this
group and its translate by $u$, of index $2^{k'}/2=2^{k'-1}$; if $u=-1$ then
$H_u=\{\pm1\}$, of infinite index.

For the last assertion, $\sigma$ fixes $\chi_p$ and every $1_j$, so it preserves
every block of $E_p=A_*\oplus\widehat\bigoplus_{j<0}F^{(j)}$ and of
$F_p^{\ge0}=\widehat\bigoplus_{j\ge0}F^{(j)}$.  Lemma~\ref{lem:block} for the
first decomposition, Theorem~\ref{thm:B2} for $A_*=\C\chi_p\oplus F_p^{\ge0}$,
and Lemma~\ref{lem:block} for the second give the stated splitting, and
\eqref{eq:D2} follows from (i)--(iii).
\end{proof}

\begin{remark}\label{rem:higherdegrees}
For a unit twist we compute only $\HH_0$.  Whether
$\HH_n(F^{(j)},(F^{(j)})_{\sigma})$ vanishes for $n\ge1$ is not determined
here: the algebraic vanishing is standard, $F^{(j)}$ being the union of the
finite-dimensional separable algebras of functions on
$\Z_p^{\times}/(1+p^{N}\Z_p)$, but the separability idempotent at level $N$ is
a separability idempotent for that level only, and its size in the seminorms
\eqref{eq:seminorms} is not bounded uniformly in $N$, so the resulting
contractions do not pass to an admissible contraction of the limit.  Nothing
below uses positive-degree information: the statements that are used are that
the twisted complex of $E_p$ is homotopy equivalent to $\C[0]$ in the non-unit
case (Lemma~\ref{lem:nonunit}) and that its reduced $\HH_0$ is
infinite-dimensional in the unit case (Lemma~\ref{lem:unit}).
\end{remark}

\begin{remark}\label{rem:closure}
The count $N(u)=(p-1)p^{k-1}=(1-p^{-1})|1-u|_p^{-1}$ holds exactly on
$1+q_p\Z_p\setminus\{1\}$, with $q_p=p$ for odd $p$ and $q_2=4$.  Outside this
set the index $[\Z_p^{\times}:H_u]$ is unrelated to $|1-u|_p$.  For $p=7$ and
$u=-1$ the closure is $\{\pm1\}$, of index $3$, while $|1-u|_7=1$.  For $p=2$,
where $1+p\Z_2$ is all of $\Z_2^{\times}$, the unit $u=3$ has $|1-u|_2^{-1}=2$
and $(p-1)p^{k-1}=1$, but $H_3=\{\pm1\}\cdot(1+8\Z_2)$ has index $2$; and
$u=-1$ has $H_{-1}=\{\pm1\}$, of infinite index.  An earlier version of
Lemma~\ref{lem:unit} asserted $H_u=1+p^{k}\Z_p$ for all $u\in1+p\Z_p$, which
fails at $p=2$, $k=1$; these are the counterexamples.
\end{remark}

\subsection{The shadow of a graph}

The decomposition of the Hochschild complex of a crossed product over the
group elements is classical in the algebraic setting; see Nistor
\cite[\S\S3--4]{Nistor}, Getzler--Jones \cite[\S3]{GetzlerJones} and
Brylinski \cite{Brylinski}.  We need a continuous version for the rapid
crossed product, and we prove it for free abelian groups by means of a
Koszul resolution of the crossed product as a bimodule.

Throughout the next proposition, $p_1,\dots,p_d$ are the primes of $S_f$,
$\Gamma=\Gamma_S\cong\Z^{d}$ with basis $p_1,\dots,p_d$, and $H$ is the height
of Section~\ref{sec:carrier}.  For a complete Fr\'echet space $V$ we write
$\ell^{1}_{H}(X;V)$, $X=\Gamma^{m}$, for the space of families
$(v_x)_{x\in X}$ with $\sum_xH(x_1)^{k}\cdots H(x_m)^{k}\nu(v_x)<\infty$ for all
$k$ and all continuous seminorms $\nu$; by Grothendieck's description of
$\ell^{1}\widehat\otimes_{\pi}$, one has
$\ell^{1}_{H}(X;V)=\ell^{1}_{H}(X)\widehat\otimes V$ and
$\ell^{1}_{H}(X)\widehat\otimes\ell^{1}_{H}(Y)=\ell^{1}_{H}(X\times Y)$.

\begin{proposition}\label{prop:crossed}
Let $A$ be a unital complete Fr\'echet algebra with an action $\alpha$ of
$\Gamma$ which is tempered with respect to $H$, let
$B=A\rtimes_{\mathrm{rapid}}\Gamma$, and let $\beta$ be a continuous
automorphism of $B$ with $\beta(A)=A$ and $\beta(U_{\gamma})=U_{\gamma}$ for
every $\gamma$.  For $\gamma\in\Gamma$ let $\alpha$ act diagonally on the
Hochschild complex $C(A,A_{\alpha_{\gamma}\beta})$, and let
$\mathcal K_{\Gamma}(\alpha)$ denote the Koszul complex
$\Lambda^{\bullet}(\C^{d})$ with the differential
$e_I\otimes v\mapsto\sum_{r}(-1)^{r-1}e_{I\setminus i_r}\otimes
(\alpha_{p_{i_r}}-1)v$, $I=\{i_1<\dots<i_m\}$.  Then the label-$\gamma$
subcomplex of the twisted Hochschild complex $C(B,B_{\beta})$ is
admissibly homotopy equivalent to
$\mathrm{Tot}\bigl(\mathcal K_{\Gamma}(\alpha)\widehat\otimes
C(A,A_{\alpha_{\gamma}\beta})\bigr)$.
\end{proposition}

Here the \emph{label} of an elementary tensor $a_0U_{g_0}\otimes\cdots\otimes
a_nU_{g_n}$ is $g_0\cdots g_n$; as in the proof of Theorem~\ref{thm:shadow}(a)
below, the twisted Hochschild boundary preserves labels, and the label
projections are continuous.  The action $\alpha$ commutes with the
Hochschild boundary of $C(A,A_{\alpha_{\gamma}\beta})$ because $\beta$
commutes with every $\alpha_g$: $\beta(\alpha_g(a))=\beta(U_gaU_g^{-1})=
U_g\beta(a)U_g^{-1}=\alpha_g(\beta(a))$.

\begin{proof}
\emph{Step 0: freeness of $B$ over $A$.}  The maps
$\ell^{1}_{H}(\Gamma)\widehat\otimes A\to B$, $\delta_g\otimes a\mapsto
aU_g$ and $\delta_g\otimes a\mapsto U_ga$, are topological isomorphisms of
left, respectively right, $A$-modules: the first is the definition of $B$,
and the second differs from it by $a\mapsto\alpha_g(a)$, since
$U_ga=\alpha_g(a)U_g$; this map and its inverse cost at most $CH(g)^{M}$ by
temperedness and are therefore absorbed by the rapid weights.  Consequently,
for a complete $A$-bimodule $X$, the map $U_g\otimes x\otimes U_h\mapsto
\delta_g\otimes x\otimes\delta_h$ identifies
$B\widehat\otimes_AX\widehat\otimes_AB$ with $\ell^{1}_{H}(\Gamma)\widehat
\otimes X\widehat\otimes\ell^{1}_{H}(\Gamma)$ as a topological vector space.
Under this identification the map induced by an $A$-bimodule map $\phi$ is
$1\otimes\phi\otimes1$; in particular, if a complex of $A$-bimodules has a
continuous linear contraction $s$ (not necessarily $A$-linear), then
$1\otimes s\otimes1$ is a continuous linear contraction of the induced
complex, and it preserves the labels $g,h$.

\emph{Step 1: a Koszul resolution of $B$.}  Put
$K_m:=B\widehat\otimes_AB\otimes\Lambda^{m}(\C^{d})$, a $B$-bimodule through
the outer factors, and for $i=1,\dots,d$ let
\[
D_i(x\otimes y):=xU_{p_i}\otimes U_{p_i}^{-1}y-x\otimes y .
\]
This is well defined on the balanced tensor product, because
$xaU_{p_i}\otimes U_{p_i}^{-1}y=xU_{p_i}\alpha_{p_i}^{-1}(a)\otimes
U_{p_i}^{-1}y=xU_{p_i}\otimes U_{p_i}^{-1}ay$; it is a bimodule map, the
$D_i$ commute, and multiplication $\mu:K_0\to B$ satisfies $\mu D_i=0$.  We
give $K_{\bullet}$ the differential $\partial(\xi\otimes e_I):=
\sum_{r}(-1)^{r-1}D_{i_r}\xi\otimes e_{I\setminus i_r}$,
$I=\{i_1<\dots<i_m\}$, with the same signs as in $\mathcal K_{\Gamma}(\alpha)$;
it squares to zero because the $D_i$ commute, and $K_{\bullet}$ is augmented
to $B$ by $\mu$.  We claim that it is admissibly exact.  By Step 0, $B\widehat\otimes_AB
\cong\ell^{1}_{H}(\Gamma\times\Gamma)\widehat\otimes A$ through
$U_g\otimes U_ha\leftrightarrow\delta_{(g,h)}\otimes a$, and in these
coordinates $D_i$ replaces $\delta_{(g,h)}$ by
$\delta_{(gp_i,hp_i^{-1})}-\delta_{(g,h)}$ while $\mu$ sends
$\delta_{(g,h)}\otimes a$ to $U_{gh}a$.  Change variables to $(g,s)$ with
$s=gh$; the weight families $H(g)^kH(h)^k$ and $H(g)^kH(s)^k$ are cofinal in
each other, since $H(s)\le H(g)H(h)$ and $H(h)\le H(g)H(s)$.  So the augmented
complex is the tensor product of $\ell^{1}_{H}(\Gamma_s)\widehat\otimes A$,
with zero differential, and the augmented Koszul complex
\begin{equation}\label{eq:koszulell}
0\to\ell^{1}_{H}(\Gamma)\otimes\Lambda^{d}\to\cdots\to
\ell^{1}_{H}(\Gamma)\otimes\Lambda^{0}\xrightarrow{\ \Sigma\ }\C\to0
\end{equation}
of the translations $\tau_i:\delta_g\mapsto\delta_{gp_i}$, augmented by
$\Sigma(\phi)=\sum_g\phi(g)$.  It suffices to give \eqref{eq:koszulell} a
continuous contraction.  Writing $\gamma=a/b$ in lowest terms, the height
satisfies $H(\gamma)\le\prod_ip_i^{|v_{p_i}(\gamma)|}=ab\le H(\gamma)^{2}$, so
$\ell^{1}_{H}(\Gamma)=\widehat\bigotimes_i\ell^{1}_{H_i}(\Z)$ with
$H_i(n)=p_i^{|n|}$, and \eqref{eq:koszulell} is the tensor product of the $d$
augmented complexes $0\to\ell^{1}_{H_i}(\Z)\xrightarrow{\tau-1}
\ell^{1}_{H_i}(\Z)\xrightarrow{\Sigma}\C\to0$.  Each of these is contracted
by $s_{-1}(1)=\delta_0$ and
\[
s_0(\delta_n)=\sum_{0\le j<n}\delta_j\quad(n>0),\qquad
s_0(\delta_0)=0,\qquad
s_0(\delta_n)=-\sum_{n\le j<0}\delta_j\quad(n<0);
\]
indeed $(\tau-1)s_0(\delta_n)=\delta_n-\delta_0=\delta_n-s_{-1}\Sigma(\delta_n)$
and $s_0((\tau-1)\delta_n)=s_0(\delta_{n+1})-s_0(\delta_n)=\delta_n$ for
every $n$.  Since $H_i(j)\le H_i(n)$ for $j$ between $0$ and $n$,
$\lVert s_0\delta_n\rVert_k\le|n|p_i^{k|n|}\le p_i^{(k+1)|n|}$, so $s_0$ is
continuous from the weight $k+1$ to the weight $k$.  If $X\to\C$ and
$X'\to\C$ are augmented complexes with augmentations $\varepsilon,\varepsilon'$,
sections $\iota,\iota'$ of them and contractions $s,s'$ with
$\partial s+s\partial=1-\iota\varepsilon$, then
$h:=s\otimes1+(-1)^{\deg}\iota\varepsilon\otimes s'$ satisfies
$\partial h+h\partial=1-\iota\varepsilon\otimes\iota'\varepsilon'$ on
$X\otimes X'$.  Applying this $d-1$ times gives a continuous contraction of
\eqref{eq:koszulell}, which proves the claim.

\emph{Step 2: a free resolution.}  Let $\mathrm{Bar}(A)\to A$ be the bar
resolution, $\mathrm{Bar}_n(A)=A\widehat\otimes A^{\widehat\otimes
n}\widehat\otimes A$, with the continuous contraction
$\xi\mapsto1\otimes\xi$.  By Step 0 the complexes
$P_{m,\bullet}:=B\widehat\otimes_A\mathrm{Bar}_{\bullet}(A)\widehat\otimes_A
B\otimes\Lambda^{m}$ are admissibly exact resolutions of $K_m$, and
$P_{m,n}\cong B\widehat\otimes A^{\widehat\otimes n}\widehat\otimes
B\otimes\Lambda^{m}$ is a free $B$-bimodule.  The maps
\[
\widetilde D_i(x\otimes a_1\otimes\cdots\otimes a_n\otimes y):=
xU_{p_i}\otimes\alpha_{p_i}^{-1}(a_1)\otimes\cdots\otimes
\alpha_{p_i}^{-1}(a_n)\otimes U_{p_i}^{-1}y-x\otimes a_1\otimes\cdots\otimes
a_n\otimes y
\]
are continuous bimodule maps lifting $D_i$; they commute with the bar
differential, because $xU_{p_i}\alpha_{p_i}^{-1}(a_1)=xa_1U_{p_i}$ and
$\alpha_{p_i}^{-1}(a_n)U_{p_i}^{-1}y=U_{p_i}^{-1}a_ny$, and they commute with
one another.  Let $P$ be the total complex of $P_{\bullet,\bullet}$, with
differential $\sum_r(-1)^{r-1}\widetilde D_{i_r}$ on the factor $e_I$ plus
$(-1)^{|I|}$ times the bar differential; it squares to zero because
$\widetilde D_i$ commutes with the bar differential.  The augmentations of the
columns give a label-preserving map of complexes of bimodules
$P\to K_{\bullet}$.  Each column $P_{m,\bullet}\to K_m$ has a continuous
linear contraction, by Step 0 applied to the bar contraction, so the cone of
$P\to K_{\bullet}$ is filtered by the $d+1$ columns with linearly contractible
graded pieces; by the perturbation lemma, the perturbation being nilpotent
because there are finitely many columns, the cone has a continuous linear
contraction, built from finitely many composites of the column contractions
and the differentials.  Composing with Step 1, the augmented complex $P\to B$
has a continuous linear contraction: $P$ is a resolution of $B$ by free
$B$-bimodules, and all maps are label-preserving.

\emph{Step 3: comparison with the bar resolution.}  The bar resolution
$\mathrm{Bar}(B)\to B$ is also a free resolution with a continuous
contraction.  By the comparison theorem of relative homological algebra, built
inductively from the contractions, there are continuous chain maps of
bimodules $\mathrm{Bar}(B)\rightleftarrows P$ over $\mathrm{id}_B$ which are
inverse to each other up to continuous homotopy.  Every space above is graded
by the total label (with $U_g$ of label $g$ and $A$ of label $1$), all
differentials preserve this grading, and so do the contractions: the bar
contraction $\xi\mapsto1\otimes\xi$ trivially, and the contraction of Step 1
because it changes $g$ at fixed $s=gh$.  Hence the comparison maps and
homotopies can be chosen label-preserving.

\emph{Step 4: the Hochschild complexes.}  Let $M=B_{\beta}$.  Then
$C(B,M)=M\widehat\otimes_{B^{e}}\mathrm{Bar}(B)$, and for a complete space
$V$ one has $M\widehat\otimes_{B^{e}}(B\widehat\otimes V\widehat\otimes B)\cong
M\widehat\otimes V$, by $m\otimes(x\otimes v\otimes y)\mapsto y\,m\,x\otimes
v$.  Under this identification $M\widehat\otimes_{B^{e}}P_{m,\bullet}$ is
$C(A,M)\otimes\Lambda^{m}$, the Hochschild complex of $A$ with coefficients
in $M$ regarded as an $A$-bimodule, and $\widetilde D_i$ becomes
\[
m\otimes a_1\otimes\cdots\otimes a_n\mapsto
U_{p_i}^{-1}\cdot m\cdot U_{p_i}\otimes\alpha_{p_i}^{-1}(a_1)\otimes\cdots
\otimes\alpha_{p_i}^{-1}(a_n)-m\otimes a_1\otimes\cdots\otimes a_n .
\]
Applying $M\widehat\otimes_{B^{e}}(-)$ to Step 3 therefore gives a continuous,
label-preserving homotopy equivalence between $C(B,M)$ and the total complex
of $C(A,M)\otimes\Lambda^{\bullet}$ with the Koszul differential built from
these operators.

\emph{Step 5: the label-$\gamma$ piece.}  As an $A$-bimodule,
$M=\widehat\bigoplus_{\gamma}AU_{\gamma}$, and $a\cdot(cU_{\gamma})\cdot a'=
acU_{\gamma}\beta(a')=ac\,\alpha_{\gamma}(\beta(a'))U_{\gamma}$; so
$AU_{\gamma}\cong A_{\alpha_{\gamma}\beta}$ by $cU_{\gamma}\mapsto c$.  The
operator $m\mapsto U_{p_i}^{-1}\cdot m\cdot U_{p_i}$ sends $cU_{\gamma}$ to
$U_{p_i}^{-1}cU_{\gamma}U_{p_i}=\alpha_{p_i}^{-1}(c)U_{\gamma}$, since
$\beta(U_{p_i})=U_{p_i}$ and $\Gamma$ is abelian.  So the label-$\gamma$ part
of the complex of Step 4 is $\mathrm{Tot}(\Lambda^{\bullet}\otimes
C(A,A_{\alpha_{\gamma}\beta}))$ with the Koszul differential of the diagonal
operators $\alpha_{p_i}^{-1}-1$.  Finally, $\alpha_{p_i}^{-1}-1=
-\alpha_{p_i}^{-1}(\alpha_{p_i}-1)$, and the $\alpha_{p_i}$ are invertible,
commute with one another and with the Hochschild boundary, so $e_I\otimes v\mapsto e_I\otimes
\prod_{i\in I}(-\alpha_{p_i})v$ is an isomorphism from
$\mathcal K_{\Gamma}(\alpha)\widehat\otimes C(A,A_{\alpha_{\gamma}\beta})$ onto
this complex.  Restricting the label-preserving homotopy equivalence of Step
4 to label $\gamma$ proves the proposition.
\end{proof}

\begin{theorem}\label{thm:shadow}
Let $w\in C_S$ be an external class.  In $D^{-}_{\mathrm{rel}}$:
\begin{enumerate}
\item[(a)] \emph{(label splitting)}
\begin{equation}\label{eq:labelsplit}
\Sh(E_w)\ \cong\ \widehat\bigoplus_{\gamma\in\Gamma_S}
\mathrm{Tot}\Bigl(\mathcal K_{\Gamma}(\alpha)\ \widehat\otimes\
C\bigl(E_P^{u},(E_P^{u})_{\gamma w}\bigr)\Bigr),
\end{equation}
where $E_P^{u}=E_P\oplus\C1$, $C(A,A_{\sigma})$ is the Hochschild
complex with $\sigma$-twisted coefficients, and $\mathcal K_{\Gamma}(\alpha)$
is the Koszul complex of $\Gamma_S\cong\Z^{|S_f|}$ of
Proposition~\ref{prop:crossed}.
\item[(b)] \emph{(K\"unneth; conditional on $(\mathrm{Exc}_{\sigma})$)}  The
coefficient factor decomposes over the places as in \eqref{eq:kunneth}.
\item[(c)] \emph{(finite places)}  $\HH_0(E_p,E_p)\cong E_p$.  If
$v_p(w_p)\ne0$ then $\HH(E_p,(E_p)_{w_p})\simeq\C[0]$, the class of the fixed
point $x_p=0$, concentrated in degree zero, of dimension one, and
\emph{independent of $w_p$}.  If $w_p=u\in\Z_p^{\times}$ with $u\ne1$, then
the reduced group $\HH_0^{\mathrm{red}}(E_p,(E_p)_u)$ is \eqref{eq:D2}, of
infinite dimension.
\end{enumerate}
\end{theorem}

\begin{proof}
(a)  An elementary tensor of
$E_w\widehat\otimes U^{\widehat\otimes n}$ carries
a total lattice label, the product of the labels of its factors.  Since
$\Gamma_S$ is abelian and $\beta_w$ fixes every gauge unitary $U_{\gamma}$, the
Hochschild boundary, including the cyclic term which passes through
$\beta_w$, preserves the total label; so the complex is the rapid direct sum of
its label pieces, the summands being closed and the projections continuous.
The identification of the label pieces is then
Proposition~\ref{prop:crossed} applied with $A=E_P^{u}$, $\Gamma=\Gamma_S$ and
$\beta=\beta_w$: the hypotheses are met, $E_P^{u}$ being unital complete
Fr\'echet with a tempered $\Gamma_S$-action by \eqref{eq:coeff}, $\Gamma_S$
free abelian of rank $|S_f|$, and $\beta_w$ continuous, preserving $E_P^{u}$
and fixing every $U_{\gamma}$.  The equivalences for the individual labels are
the restrictions of the single label-preserving homotopy equivalence of Step 4
of that proof, and the isomorphism of Step 5 does not depend on $\gamma$; so
they assemble into \eqref{eq:labelsplit}, where $\widehat\bigoplus_{\gamma}$
denotes the $H$-rapid sum, the topology being inherited from
$E_w=\widehat\bigoplus_{\gamma}E_P^{u}U_{\gamma}$.

(b) is Theorem~\ref{thm:kunneth} applied to $\sigma=\gamma w$ acting
coordinatewise, together with Theorem~\ref{thm:B2} for the adjoined unit.

(c)  Since $E_p$ is commutative, the untwisted statement is immediate from
$\HH_0(E_p,E_p)=E_p/[E_p,E_p]$.  The two twisted statements are respectively
Lemma~\ref{lem:nonunit}, where $v_p(w_p)\ne0$, and Lemma~\ref{lem:unit}, where
$w_p$ is a unit.  In the first of them the conclusion is a homotopy
equivalence of complexes, not merely a statement about homology in degree
zero.  Parts (a) and (c) do not use $(\mathrm{Exc}_{\sigma})$.
\end{proof}

\begin{remark}\label{rem:notorsion}
The fixed point $x_p=0$ of a nontrivial dilation contributes $\C[0]$: rank one,
no differential.  In the smooth polynomial model $\R[z]$ the same fixed point
contributes the Koszul complex
$[\,\R[z]\xrightarrow{(\sigma-1)z}\R[z]\,]$, acyclic over $\C$ but with torsion
$|\sigma-1|$.  The totally disconnected coefficient algebra has no tangent
direction at the origin, so the Koszul differential has nowhere to live, and
its determinant survives only in the \emph{normalization} of the point class.
That normalization is what Section~\ref{sec:trace} measures.
\end{remark}

\begin{corollary}\label{cor:blind}
Let $w$ be an external class.  At every finite place where $v_p(w_p)\ne0$ the
twisted complex $C(E_p,(E_p)_{w_p})$ is homotopy equivalent to the single
origin class $\C[0]$, independent of $w_p$; at every finite place where $w_p$
is a unit $\ne1$ its reduced $\HH_0$ is infinite dimensional.  If
$(\mathrm{Exc}_{\sigma})$ holds, these complexes are the local factors of
$\Sh(E_w)$ by Theorem~\ref{thm:shadow}(b), and then any pairing on $G$ defined through ranks, Euler characteristics
or determinant lines of the shadow is either \emph{constant} on every family
all of whose finite components have nonzero valuation, or \emph{undefined} on
unit components.  In particular no boundary Euler-characteristic pairing and
no ``determinant of the shadow'' separates the generators $E_a$.
\end{corollary}

\begin{proof}
Immediate from Theorem~\ref{thm:shadow}(c).  A complex concentrated in degree
zero with one-dimensional homology has determinant line canonically $\C$ and
no torsion, and its rank and Euler characteristic are $1$ whatever $w_p$ is; a
complex with infinite-dimensional $\HH_0$ has no finite rank, no Euler
characteristic and no determinant line, whatever its higher homology.
\end{proof}

By contrast, the analytic trace constructed in the next section \emph{does}
separate; see Theorem~\ref{thm:separate}.  At a unit twist
$u\in1+q_p\Z_p\setminus\{1\}$ the finite quantity attached to one shell is the
orbit count $N(u)=(p-1)p^{k-1}=c_p|1-u|_p^{-1}$, $c_p=1-p^{-1}$; its relation to
the density $\tau_p$ is recorded in Proposition~\ref{prop:haarrank}.

\section{The trace}\label{sec:trace}

\subsection{The local densities}

\begin{definition}\label{def:tau}
For a place $v$ of $\Q$ and $x\in\Q_v^{\times}$ with $x\ne1$ put
\begin{equation}\label{eq:tau}
\tau_v(x):=\frac{|x|_v^{1/2}}{|1-x|_v}.
\end{equation}
\end{definition}

\begin{lemma}\label{lem:tau}
Let $p$ be finite and $x=p^{k}u$ with $u\in\Z_p^{\times}$.
\begin{enumerate}
\item[(i)] If $k\ne0$ then $\tau_p(x)=p^{-|k|/2}$.
\item[(ii)] If $k=0$ and $x\ne1$ then $\tau_p(x)=|1-x|_p^{-1}=p^{\,v_p(1-x)}$,
a finite number.
\end{enumerate}
Moreover the finite-rank cutoff traces of the dilation by $x$, which equal
$|x|_p^{1/2}\min\bigl(p^{L},|1-x|_p^{-1}\bigr)$ at cutoff level $L$, take the
exact value $\tau_p(x)$ for every $L\ge v_p(1-x)$; in case (i) this holds for
every $L\ge0$.
\end{lemma}

\begin{proof}
(i) $|x|_p^{1/2}=p^{-k/2}$.  If $k>0$ then $|x|_p<1$, so $|1-x|_p=1$ and
$\tau_p=p^{-k/2}$.  If $k<0$ then $|x|_p>1$, so $|1-x|_p=|x|_p=p^{-k}$ and
$\tau_p=p^{-k/2}p^{k}=p^{k/2}=p^{-|k|/2}$.  In both cases $|1-x|_p\ge1$, so
$\min(p^{L},|1-x|_p^{-1})=|1-x|_p^{-1}$ for every $L\ge0$.

(ii) For a fixed unit $x\ne1$ the integer $j:=v_p(1-x)\ge0$ is finite and
$|1-x|_p^{-1}=p^{j}$; hence $\min(p^{L},p^{j})=p^{j}$ for every $L\ge j$.
\end{proof}

\begin{remark}\label{rem:unitdiv}
The pointwise finiteness of (ii) is compatible with the divergence of the
integral over the unit group,
$\int_{\Z_p^{\times}}\min(p^{L},|1-u|_p^{-1})\,du=L+1$, because
$p^{v_p(1-u)}$ is unbounded as $u\to1$.  This is why the regularization below
touches only a neighbourhood of the identity class.
\end{remark}

\begin{lemma}\label{lem:inv}
$\tau_v(x^{-1})=\tau_v(x)$ for every place $v$ and every $x\ne1$.
\end{lemma}

\begin{proof}
$|x^{-1}|_v^{1/2}=|x|_v^{-1/2}$ and
$|1-x^{-1}|_v=|x|_v^{-1}|x-1|_v$; dividing, the factors $|x|_v^{-1/2}$ and
$|x|_v^{-1}$ combine to $|x|_v^{1/2}$.
\end{proof}

At the archimedean place, in the multiplicative coordinate $x=\pm e^{t}$,
\begin{equation}\label{eq:tauinf}
\tau_{\infty}(+e^{t})=\frac{1}{2|\sinh(t/2)|},\qquad
\tau_{\infty}(-e^{t})=\frac{1}{2\cosh(t/2)} .
\end{equation}

\subsection{The trace}

For $v\in S$ let $j_v:\Q_v^{\times}\to C_S$ be the map sending $x$ to the class
of the idele having $x$ at $v$ and $1$ at every other place of $S$.  By
Lemma~\ref{thm:split}, the canonical representative of that class in
$C_S\cong\R^{\times}\times\prod_{q\in S_f}\Z_q^{\times}$ is obtained by
dividing by $\varphi$ of the idele, and therefore
\begin{equation}\label{eq:slices}
\begin{aligned}
j_p(p^{k}w)&=\bigl(p^{-k}\text{ at }\infty,\ w\text{ at }p,\
p^{-k}\text{ at every other place of }S\bigr),\qquad p\in S_f,\\
j_{\infty}(x)&=\bigl(x\text{ at }\infty,\ 1\text{ at every finite place of }S
\bigr).
\end{aligned}
\end{equation}

\begin{definition}\label{def:trace}
Fix a real constant $c$.  The \emph{trace} is the linear functional on $\DA$
defined by its local decomposition
\begin{equation}\label{eq:tracedef}
T(F):=\sum_{v\in S}W_v(F)+c\,F(1),\qquad F\in\DA ,
\end{equation}
where, at a finite place $p$, with $a:=F\circ j_p$ and with $\dstar x$ giving
every shell $p^{k}\Z_p^{\times}$ the volume $\log p$,
\begin{equation}\label{eq:Wp}
W_p(F)=\int_{v_p(x)\neq0}a(x)\,\tau_p(x)\dstar x
\ +\ (\log p)\int_{\Z_p^{\times}}\frac{a(u)-a(1)}{|1-u|_p}\,du,
\end{equation}
the second integral against the probability Haar measure of
$\Z_p^{\times}$; and at the archimedean place, with the two profiles
\[
b_{+}(t):=F\bigl(j_{\infty}(e^{t})\bigr),\qquad
b_{-}(t):=F\bigl(j_{\infty}(-e^{t})\bigr),\qquad t\in\R,
\]
of $F$ on the two real half-lines (so that $b_{+}(0)=F(1)$),
\begin{equation}\label{eq:Winf}
\begin{aligned}
W_{\infty}(F)={}&(\gamma+\log4\pi)\,b_{+}(0)
+\int_0^{\infty}\Bigl[K_{s}(t)\bigl(b_{+}(t)+b_{+}(-t)\bigr)
-2K(t)\,b_{+}(0)\,e^{-t/2}\Bigr]dt\\
&+\int_0^{\infty}K_{c}(t)\bigl(b_{-}(t)+b_{-}(-t)\bigr)dt,
\end{aligned}
\end{equation}
with the kernels
\[
K_{s}(t)=\frac{1}{4\sinh(t/2)},\qquad K_{c}(t)=\frac{1}{4\cosh(t/2)},\qquad
K(t)=K_{s}(t)+K_{c}(t)=\frac{e^{-t/2}}{1-e^{-2t}} .
\]
\end{definition}

\begin{convention}[The archimedean term for all test functions]\label{conv:Winf}
Formula \eqref{eq:Winf} is the definition of $W_{\infty}$ on all of $\DA$,
with no parity assumption: the positive half-line enters only through $b_{+}$
and the kernel $K_{s}$, the negative half-line only through $b_{-}$ and the
kernel $K_{c}$, and the regularization (the constant $\gamma+\log4\pi$ and the
counterterm $2K(t)b_{+}(0)e^{-t/2}$) involves only the value $b_{+}(0)=F(1)$ at
the identity, which is the only point where $\tau_{\infty}$ is singular.  Both
integrals in \eqref{eq:Winf} converge absolutely: near $t=0$ one has
$2K_{s}(t)=t^{-1}+O(t)$ and $2K(t)e^{-t/2}=t^{-1}+O(1)$, so the first bracket is
$b_{+}(t)+b_{+}(-t)-2b_{+}(0)$ times $O(t^{-1})$ plus a bounded function times
$b_{+}(0)$, hence bounded; at infinity $K_{s},K_{c}=O(e^{-t/2})$ and
$K(t)e^{-t/2}=O(e^{-t})$, while $b_{\pm}$ have compact support.  Writing
$F=F^{\mathrm{ev}}+F^{\mathrm{odd}}$ for the decomposition under translation by
$j_{\infty}(-1)$, the even part has $b_{+}=b_{-}=:b$ and \eqref{eq:Winf}
becomes
\begin{equation}\label{eq:Winfeven}
W_{\infty}(F^{\mathrm{ev}})=(\gamma+\log4\pi)\,b(0)
+\int_0^{\infty}K(t)\bigl[b(t)+b(-t)-2b(0)e^{-t/2}\bigr]dt,
\end{equation}
while for the odd part $b_{-}=-b_{+}$ and the $K_{c}$-integral enters with the
opposite sign.  If $F$ vanishes on a neighbourhood of $1$, then $b_{+}(0)=0$,
the counterterm drops, and since $\tau_{\infty}(\pm e^{t})$ equals
$2K_{s}(|t|)$, respectively $2K_{c}(|t|)$, by \eqref{eq:tauinf}, while
$\dstar x=dt/2$ on each half-line by \eqref{eq:archmeasure},
\begin{equation}\label{eq:Winfint}
\begin{aligned}
W_{\infty}(F)&=\int_{\R^{\times}}F\bigl(j_{\infty}(x)\bigr)\,\tau_{\infty}(x)\,
\frac{dx}{2|x|}\\
&=\int_0^{\infty}K_{s}(t)\bigl(b_{+}(t)+b_{+}(-t)\bigr)dt
+\int_0^{\infty}K_{c}(t)\bigl(b_{-}(t)+b_{-}(-t)\bigr)dt .
\end{aligned}
\end{equation}
Thus $W_{\infty}$ is the regularized integral of $F$ against $\tau_{\infty}$
over $\R^{\times}$, with the measure \eqref{eq:archmeasure}, regularized at the
identity only; this is the meaning of $\int'_{\R^{\times}}$ wherever it
appears below.  The factor $\tfrac12$ in \eqref{eq:archmeasure} is the folding
factor of the two half-lines.
\end{convention}

The two local terms are the regularized integrals of $F$ against the densities
$\tau_v$ of \eqref{eq:tau}: at a finite place this is manifest from
Lemma~\ref{lem:tau}, and at the archimedean place it is \eqref{eq:Winfint}.

$T$ is linear, and it is continuous on each fixed-support Fr\'echet step of
$\DA$: each $W_v$ is integration of $F$ against a fixed distribution, and the
counterterms $a(1)$, $b_{+}(0)$ are evaluations.

\begin{theorem}[Absolute convergence off the identity class]\label{thm:absconv}
For $F\in\DA^{\mathrm{off}}$,
\begin{equation}\label{eq:absconv}
T(F)=\sum_{v\in S}\int_{\Q_v^{\times}}F(j_v(x))\,\tau_v(x)\dstar x ,
\end{equation}
an absolutely convergent sum of integrals, with no constant term, no
counterterm and no limiting operation.
\end{theorem}

\begin{proof}
For $F\in\DA^{\mathrm{off}}$ one has $F(1)=0$, so the term $cF(1)$ drops
whatever $c$ is; and $a=F\circ j_v$ vanishes on a neighbourhood of $1$ in
$\Q_v^{\times}$, because $j_v$ is continuous with $j_v(1)=1$.  Hence the
counterterms $(\log p)a(1)$, $(\gamma+\log4\pi)b_{+}(0)$ and $2K(t)b_{+}(0)e^{-t/2}$ all
vanish.  What remains at a finite place is
$\int_{v_p\ne0}a\tau_p\dstar x+(\log p)\int_{\Z_p^{\times}}a(u)|1-u|_p^{-1}du$.
The first is a finite sum over the finitely many shells met by the compact
support of $a$, each term bounded by $\sup|a|\cdot p^{-|k|/2}\log p$; the
second is the integral of a compactly supported locally constant function
against a function bounded on $\supp a$, by Lemma~\ref{lem:tau}(ii), since
$|1-u|_p^{-1}\le p^{j_0}$ with $j_0$ determined by the distance of $\supp a$
from $1$.  At the archimedean place \eqref{eq:Winfint} applies: $K_{s}$ is continuous
on the compact support of $b_{+}$, which avoids $t=0$, and $K_{c}$ is
continuous everywhere.
\end{proof}

\begin{theorem}[Uniqueness off the identity class]\label{thm:unique}
Let $T'$ be a continuous linear functional on $\DA$ agreeing with
\eqref{eq:absconv} on $\DA^{\mathrm{off}}$, and assume that its archimedean
part is a distribution of finite order.  Then $T-T'$ is supported at the
identity class; its finite-place components are multiples of $\delta_1$ and its
archimedean component is a finite combination of derivatives of $\delta_1$.  In
particular the one-parameter family \eqref{eq:tracedef}, $c\in\R$, is the only
ambiguity of order zero.
\end{theorem}

\begin{proof}
$T-T'$ annihilates every test vanishing near $1$, hence is supported at
$\{1\}$.  On the totally disconnected factors, $\DA$ consists of locally
constant functions and the subgroups $1+p^{m}\Z_p$ form a neighbourhood basis
of $1$ by compact open subgroups; a continuous functional vanishing on all
tests supported off $1$ is determined by its values on the characteristic
functions of those subgroups, and these stabilize by continuity, so the
functional is a multiple of $\delta_1$.  On the real factor, a distribution of
finite order supported at a point is a finite combination of derivatives of
$\delta_1$.
\end{proof}

\begin{proposition}[Comparison with an analytic
regularization]\label{prop:comparison}
Suppose that for compactly supported tests the cutoff operators
$A_{\Lambda,S}(F)=R_{\Lambda}\theta(F)$ are trace class and satisfy
\begin{equation}\label{eq:cutoff}
\Tr\bigl(A_{\Lambda,S}(F)\bigr)=2F(1)\log\Lambda+T_{\mathrm{an}}(F)+o(1)
\qquad(\Lambda\to\infty)
\end{equation}
for a continuous limiting functional $T_{\mathrm{an}}$ of finite archimedean
order, whose off-identity part is given by the local integrals with densities
$\tau_v$.  Then there is a real constant $c_{\mathrm{cut}}$ with
$T_{\mathrm{an}}=\sum_{v\in S}W_v+(c+c_{\mathrm{cut}})F(1)$; that is,
$T_{\mathrm{an}}$ is a member of the family \eqref{eq:tracedef}.
\end{proposition}

\begin{proof}
Off the identity class both functionals are given by the same absolutely
convergent local integrals, by hypothesis and by
Theorem~\ref{thm:absconv}.  Apply Theorem~\ref{thm:unique}: the residue is of
order zero at the finite places by the argument given there, and at the real
place the counterterm in \eqref{eq:cutoff} is $2F(1)\log\Lambda$, of order zero
in $F$, so the residue is a multiple of $F(1)$.
\end{proof}

Proposition~\ref{prop:comparison} is the only place where an operator
regularization appears; we do not determine $c_{\mathrm{cut}}$, and it is not
needed later, since the trace is defined by \eqref{eq:tracedef}.

\begin{remark}\label{rem:diagonal}
The diagonal $\Delta=E_1$ corresponds to the Dirac section at the identity
class, which is a multiplier of the section algebra and not an element of it.
So $T(\Delta)$ is not defined.  This is to be expected: the identity operator
on an infinite-dimensional space is not trace class, and the counterterms in
\eqref{eq:Wp}--\eqref{eq:Winf} play the role of the diagonal.
\end{remark}

\begin{convention}\label{conv:c}
From here on $c$ denotes a fixed real constant, the same for every set $S$.
All results are proved for arbitrary $c$ unless the contrary is stated; the
value $c=0$ is adopted where a specific normalization is needed, namely in
Theorem~\ref{thm:RR}(f) and the auxiliary
Theorem~\ref{thm:mertens} of Appendix~\ref{app:mertens}.  The choice $c=0$ is the statement that no global
constant is added on top of the intrinsic constants of the local symbols
computed in Theorem~\ref{thm:ML}.
\end{convention}

\subsection{Separation}

The following computation involves no regularization.  It shows that the
trace distinguishes classes that the additive invariants of
Corollary~\ref{cor:blind} cannot.

\begin{theorem}\label{thm:separate}
Fix $p\in S_f$ and units $u,u'\in1+p\Z_p$ with
$j:=v_p(1-u)\ne j':=v_p(1-u')$, and fix $m>\max(j,j')$.  Let
$\theta:=\min_{q\in S_f}q^{1/2}>1$ and let $h\in C_c^{\infty}(\R^{\times})$
satisfy $h\ge0$, $\supp h\subset(\theta^{-1},\theta)$ and $h(1)>0$.  In the
coordinates of Lemma~\ref{thm:split} define
\begin{equation}\label{eq:sepsections}
F\bigl(x_{\infty},(w_q)\bigr):=h(x_{\infty})\cdot
\mathbf 1_{u(1+p^{m}\Z_p)}(w_p)\cdot\prod_{q\ne p}\mathbf 1_{\Z_q^{\times}}(w_q),
\end{equation}
and $F'$ by the same formula with $u$ replaced by $u'$.  Then $F,F'\in
\DA^{\mathrm{off}}$ and, with $\nu:=h(1)\vol(1+p^{m}\Z_p)>0$,
\begin{equation}\label{eq:sepvalues}
T(F)=\nu(\log p)p^{\,j},\qquad T(F')=\nu(\log p)p^{\,j'},\qquad
T(F)-T(F')=\nu(\log p)\bigl(p^{\,j}-p^{\,j'}\bigr)\neq0,
\end{equation}
while $\lambda_-(F)=\lambda_-(F')$ and $\lambda_+(F)=\lambda_+(F')$ for the two
functionals of \eqref{eq:lambdadef} below.  In particular $T$ separates two
sections with identical pushforward to the modulus group.
\end{theorem}

\begin{proof}
\emph{Membership.}  $F$ and $F'$ are finite products of a compactly supported
smooth function of $x_{\infty}$ with characteristic functions of compact open
subsets of the compact groups $\Z_q^{\times}$, hence lie in $\DA$.  The
identity class is $1=(1,(1)_q)$; since $u\in1+p\Z_p$ with $v_p(1-u)=j<m$, we
have $1\notin u(1+p^{m}\Z_p)$, and that set is open and closed, so $F$ vanishes
on a neighbourhood of $1$.  By Theorem~\ref{thm:absconv} the value $T(F)$ is
the absolutely convergent sum \eqref{eq:absconv}, with no counterterm and no
dependence on $c$.

\emph{The archimedean term vanishes.}  By \eqref{eq:slices} the archimedean
slice is $j_{\infty}(x)=(x,1,\dots,1)$, whose $p$-coordinate is
$1\notin u(1+p^{m}\Z_p)$; hence $F\circ j_{\infty}\equiv0$ and $W_{\infty}(F)=0$.

\emph{The other finite terms vanish.}  Let $q\in S_f$, $q\ne p$.  By
\eqref{eq:slices} the $p$-coordinate of $j_q(q^{k}w)$ is $q^{-k}$.  For $k=0$
this is $1\notin u(1+p^{m}\Z_p)$, so $F\circ j_q$ vanishes on the units and
both the unit packet and the value $a(1)$ vanish.  For $k\ne0$ the
archimedean coordinate of $j_q(q^kw)$ is $q^{-k}$, which is not in $\supp h$
because $|k|\ge1$ and $\supp h\subset(\theta^{-1},\theta)$ with
$\theta\le q^{1/2}<q$.  So $F\circ j_q\equiv0$ and $W_q(F)=0$.

\emph{The term at $p$.}  By \eqref{eq:slices},
$a(p^{k}w)=h(p^{-k})\mathbf 1_{u(1+p^{m}\Z_p)}(w)$.  For $k\ne0$ we have
$p^{-k}\notin\supp h$, so the shell integral vanishes.  For $k=0$,
$a(w)=h(1)\mathbf 1_{u(1+p^{m}\Z_p)}(w)$ and $a(1)=0$, so by \eqref{eq:Wp}
\[
W_p(F)=(\log p)\,h(1)\int_{u(1+p^{m}\Z_p)}|1-w|_p^{-1}\,dw .
\]
On the coset write $w=u(1+p^{m}y)$ with $y\in\Z_p$; then $1-w=(1-u)-up^{m}y$,
and since $v_p(1-u)=j<m$ we get $v_p(1-w)=j$ exactly, so $|1-w|_p^{-1}=p^{j}$
is constant on the coset.  Therefore
$W_p(F)=(\log p)h(1)p^{j}\vol(1+p^{m}\Z_p)=\nu(\log p)p^{j}$, with $\nu>0$
because $h(1)>0$.  The same computation with $u'$, and subtraction, gives
\eqref{eq:sepvalues}; the difference is nonzero because $j\ne j'$.

\emph{The Tate functionals agree.}  By Lemma~\ref{thm:split}, $|w|_S=|w_\infty|$
on $C_S$, so
\[
\lambda_{\mp}(F)=\Bigl(\int_{\R^{\times}}h(x)|x|^{\mp1/2}\dstar x\Bigr)
\cdot\vol\bigl(u(1+p^{m}\Z_p)\bigr)\cdot\prod_{q\ne p}\vol(\Z_q^{\times}),
\]
and $\vol(u(1+p^{m}\Z_p))=\vol(1+p^{m}\Z_p)=\vol(u'(1+p^{m}\Z_p))$ by
translation invariance of Haar measure on $\Z_p^{\times}$.
\end{proof}

\begin{proposition}[The local weight and the shadow]\label{prop:haarrank}
Let $p\in S_f$, $x\in\Q_p^{\times}$, $x\ne1$, and $c_p:=1-p^{-1}$.
\begin{enumerate}
\item[(a)] If $v_p(x)\ne0$, the local factor of the shadow at the dilation $x$
is $\C[0]$, of rank $1$ (Lemma~\ref{lem:nonunit}), and
$\tau_p(x)=|x|_p^{1/2}|1-x|_p^{-1}=p^{-|v_p(x)|/2}$.
\item[(b)] If $x\in1+q_p\Z_p$ ($q_p=p$ for odd $p$, $q_2=4$), the reduced
$\HH_0$ of each shell has dimension $N(x)=(p-1)p^{k-1}$, $k=v_p(1-x)$
(Lemma~\ref{lem:unit}), and
\[
\tau_p(x)=|1-x|_p^{-1}=N(x)/c_p .
\]
\item[(c)] If $x$ is a unit outside $1+q_p\Z_p$, then $\tau_p(x)=|1-x|_p^{-1}$
while the reduced $\HH_0$ of a shell has dimension $[\Z_p^{\times}:H_x]$, and
the two are unrelated: for $p$ odd and $x\notin1+p\Z_p$ one has
$\tau_p(x)=1$, and for $p=2$, $x=3$ one has $\tau_2(3)=2$ while
$[\Z_2^{\times}:H_3]/c_2=4$; for $x=-1$ the index is infinite.
\end{enumerate}
\end{proposition}

\begin{proof}
(a) is Lemma~\ref{lem:tau}(i) together with Lemma~\ref{lem:nonunit}.  (b) is
$(p-1)p^{k-1}=(1-p^{-1})p^{k}$ and $|1-x|_p^{-1}=p^{k}$, together with
Lemma~\ref{lem:unit}(iii).  (c) is Lemma~\ref{lem:unit}(iii) and
Remark~\ref{rem:closure}.
\end{proof}

Proposition~\ref{prop:haarrank} is a comparison, not a derivation: the trace
of Definition~\ref{def:trace} is defined by the densities $\tau_v$ directly, and
it is not obtained from the shadow.  On the shells and on the units of
$1+q_p\Z_p$ the density agrees with the orbit count of the shadow divided by
$c_p$; on the remaining units, which carry part of the unit packet of
\eqref{eq:Wp}, it does not.  Together with Corollary~\ref{cor:blind} this is
why the numerical theory below is analytic.

\section{The Tate characters and the symbol of the trace}\label{sec:symbol}

\subsection{The two Tate characters}

\begin{theorem}\label{thm:tate}
The map
\begin{equation}\label{eq:rhoT}
\rho_T(u):=\bigl(|u|_S^{-1/2},\,|u|_S^{+1/2}\bigr)
\end{equation}
is well defined on the class group $C_S$, independent of the lift, and
satisfies $\rho_T(u^{-1})=\rho_T(u)^{\sharp}$ for the swap $\sharp$ of the two
coordinates.  Integrating, the two \emph{Tate functionals}
\begin{equation}\label{eq:lambdadef}
\lambda_{\mp}(F):=\int_{C_S}F(u)\,|u|_S^{\mp1/2}\dstar u,\qquad F\in\DA,
\end{equation}
are continuous --- they are the local zeta functionals of Tate's thesis
\cite{Tate} evaluated at the two exponents $\mp\tfrac12$ --- and satisfy
$\lambda_{\mp}(F\star G)=\lambda_{\mp}(F)\lambda_{\mp}(G)$ and
$\lambda_{\mp}(F^{\dagger})=\overline{\lambda_{\pm}(F)}$; and
$M:=(\lambda_-,\lambda_+):\DA\to\BT:=\C_-\oplus\C_+$ is a surjective
continuous algebra homomorphism.  With $e_-=(1,0)$ and $e_+=(0,1)$, the
Hermitian form
$\beta_T(z,w)=z_-\overline{w_+}+z_+\overline{w_-}$ has matrix
$\left(\begin{smallmatrix}0&1\\1&0\end{smallmatrix}\right)$.
\end{theorem}

\begin{proof}
$|\gamma|_S=1$ for $\gamma\in\Gamma_S$ by Lemma~\ref{lem:prodformula}, so
$|\cdot|_S^{\pm1/2}$ is constant on $\Gamma_S$-orbits and \eqref{eq:rhoT}
descends to $C_S$.  Multiplicativity of $\lambda_{\mp}$ is the Fourier
transform of convolution at a fixed character; continuity is continuity of
integration against a locally bounded character on each fixed-support step.
For the involution, substituting $v=u^{-1}$ and using inversion invariance,
\[
\lambda_{\mp}(F^{\dagger})=\int\overline{F(u^{-1})}|u|_S^{\mp1/2}\dstar u
=\overline{\int F(v)|v|_S^{\pm1/2}\dstar v}=\overline{\lambda_{\pm}(F)}.
\]
Surjectivity: $|\cdot|_S:C_S\to\R_{>0}$ is surjective, so $\lambda_{\mp}$ are
the Mellin transforms of the radial pushforward of $F$ at the two distinct
points $s=\mp1/2$; and a compactly supported smooth function on $\R_{>0}$ may
have prescribed Mellin values at those two points, because the two functionals
$f\mapsto\int fe^{t/2}$ and $f\mapsto\int fe^{-t/2}$ on $C_c^{\infty}(\R)$ are
linearly independent, the functions $e^{t/2}$ and $e^{-t/2}$ being linearly
independent on $\R$.  (The same argument, with two further functionals, is
written out in the proof of Theorem~\ref{thm:basechange}(c).)
\end{proof}

\subsection{The dual group}

By Lemma~\ref{thm:split}, $C_S\cong\R^{\times}\times\prod_{q\in
S_f}\Z_q^{\times}$, so its Pontryagin dual is
\begin{equation}\label{eq:dual}
\Ch=\bigl\{\chi=(\varepsilon,\tau;(\kappa_q))\bigr\},\qquad
\chi(w)=\sgn(w_{\infty})^{\varepsilon}|w_{\infty}|^{i\tau}
\prod_{q\in S_f}\kappa_q(w_q),
\end{equation}
with $\varepsilon\in\{0,1\}$, $\tau\in\R$ and $\kappa_q$ a finite-order
character of $\Z_q^{\times}$ of conductor exponent $c(\kappa_q)\ge0$.  As a
topological space $\Ch$ is a disjoint union of copies of $\{0,1\}\times\R$,
one for each of the countably many systems $(\kappa_q)$; we call these its
\emph{components}, and we write $\mathfrak c_0$ for the \emph{principal
component}, $\varepsilon=0$ and $\kappa_q=1$ for all $q$.  Set
\begin{equation}\label{eq:FT}
\widehat F(\chi):=\int_{C_S}F(w)\overline{\chi(w)}\dstar w .
\end{equation}
If $F\in\DA$ is invariant under $\prod_q(1+q^{N}\Z_q)$ then $\widehat F$ is
supported on the finitely many components with $c(\kappa_q)\le N$ for all $q$,
and on each of those it is a Schwartz function of $\tau$, the radial profile of
$F$ being $C_c^{\infty}$.  Indeed, by Paley--Wiener, each $\widehat
F(\mathfrak c,\cdot)$ extends to an entire function of $\tau$ of exponential
type.  In particular the two Tate functionals are values of the analytic
continuation of $\widehat F$ off the unitary axis:
\begin{equation}\label{eq:lambdaFT}
\lambda_{\mp}(F)=\widehat F\bigl(\varepsilon=0,\ \tau=\mp i/2;\ (1)\bigr).
\end{equation}
They are therefore \emph{not} evaluations at points of $\Ch$.  We normalize
$d\chi|_{\mathfrak c}=d\tau/2\pi$ on each component, the Plancherel measure
dual to $\dstar w$.

\subsection{The symbol}

\begin{theorem}\label{thm:symbol}
For $\chi\in\Ch$ put
\begin{equation}\label{eq:symboldef}
\Th(\chi):=\Wh_{\infty}(\chi)+\sum_{p\in S_f}\Wh_p(\chi),\qquad
\Wh_v(\chi):=\int'_{\Q_v^{\times}}\chi(j_v(x))\,\tau_v(x)\dstar x,
\end{equation}
with the regularization and the measures of Definition~\ref{def:trace}.  Then
each integral converges and:
\begin{enumerate}
\item[(a)] \emph{(finite places)}  Put
$\eta_p(\chi):=p^{-i\tau}\prod_{q\in S_f,\,q\ne p}\kappa_q(p)^{-1}$, of modulus
one, and $\zeta_p(\chi):=p^{-1/2}\overline{\eta_p(\chi)}$.  Then
\begin{equation}\label{eq:Wphat}
\Wh_p(\chi)=
\begin{cases}
2(\log p)\,\Rel\bigl(\zeta_p/(1-\zeta_p)\bigr), & \kappa_p=1,\\[2pt]
-c(\kappa_p)\log p, & \kappa_p\neq1 .
\end{cases}
\end{equation}
In both cases $\Wh_p$ is real, real-analytic in $\tau$, and bounded on each
component by $2(\log p)p^{-1/2}(1-p^{-1/2})^{-1}$, respectively by
$c(\kappa_p)\log p$.
\item[(b)] \emph{(the real place)}  $\Wh_{\infty}$ is real, real-analytic in
$\tau$ on each component, bounded above, and
\begin{equation}\label{eq:Winfhat}
\Wh_{\infty}(\varepsilon,\tau)=-\log|\tau|+O(1)\qquad(|\tau|\to\infty),
\end{equation}
uniformly in $\varepsilon$.  For $\varepsilon=0$ one has the closed form
\begin{equation}\label{eq:Winfclosed}
\Wh_{\infty}(0,\tau)=\log\pi-\Rel\,\psi\bigl(\tfrac14+\tfrac{i\tau}{2}\bigr),
\end{equation}
$\psi$ the digamma function.  For $\varepsilon=1$, with $W_{\infty}$ as in
Convention~\ref{conv:Winf}, $\Wh_{\infty}(1,\tau)=\log\pi-\Rel\,\psi(\tfrac34+
\tfrac{i\tau}2)=\Wh_{\infty}(0,\tau)-\pi/\cosh(\pi\tau)$, see
\eqref{eq:Winfodd}.
\item[(c)] $\Th$ is real-valued, real-analytic on each component, bounded above
by a finite constant $m:=\sup_{\Ch}\Th<\infty$, and
$\Th(\varepsilon,\tau;(\kappa_q))\to-\infty$ as $|\tau|\to\infty$.  Hence for
every $c\in\R$ the function $\Th+c$ vanishes identically on no component, and
its zero set is a discrete subset of each component.
\item[(d)] \emph{(Plancherel form)}  For every $H\in\DA$,
\begin{equation}\label{eq:plancherel}
T(H)=\int_{\Ch}\widehat H(\chi)\bigl(\Th(\chi)+c\bigr)d\chi ,
\end{equation}
the integral being absolutely convergent; and the interchange is justified
termwise, so that for every $v\in S$,
\begin{equation}\label{eq:plancherelv}
W_v(H)=\int_{\Ch}\widehat H(\chi)\,\Wh_v(\chi)\,d\chi ,
\end{equation}
also absolutely convergent.
\end{enumerate}
\end{theorem}

\begin{proof}
(a)  By \eqref{eq:slices}, $\sgn(p^{-k})=+1$ and
\[
\chi\bigl(j_p(p^{k}w)\bigr)=p^{-ik\tau}\kappa_p(w)\prod_{q\ne p}\kappa_q(p)^{-k}
=\eta_p(\chi)^{k}\kappa_p(w).
\]
Insert this into \eqref{eq:Wp} with $a=\chi\circ j_p$, which has $a(1)=1$.

\emph{Shell part.}  By Lemma~\ref{lem:tau}(i) and the shell volume $\log p$,
\[
\int_{p^{k}\Z_p^{\times}}a\,\tau_p\dstar x
=(\log p)\,p^{-|k|/2}\,\eta_p^{k}\int_{\Z_p^{\times}}\kappa_p(w)\,dw ,
\]
and the unit integral is $1$ if $\kappa_p=1$ and $0$ otherwise, by
orthogonality.  For $\kappa_p=1$ the sum over
$k\ne0$ is $\sum_{k\ge1}p^{-k/2}(\eta_p^{k}+\eta_p^{-k})$; with
$\zeta:=p^{-1/2}\overline{\eta_p}$ and $\eta_p^{-1}=\overline{\eta_p}$ this is
$\sum_{k\ge1}(\bar\zeta^{k}+\zeta^{k})=2\Rel(\zeta/(1-\zeta))$, absolutely
convergent since $|\zeta|=p^{-1/2}<1$, and manifestly real.

\emph{Unit packet.}  For $k=0$, $j_p(u)=(1,u,1,\dots)$, so the packet is
$(\log p)\int_{\Z_p^{\times}}(\kappa_p(u)-1)|1-u|_p^{-1}du$ against the
probability Haar measure.  We claim
\begin{equation}\label{eq:conductor}
\int_{\Z_p^{\times}}\frac{\kappa_p(u)-1}{|1-u|_p}\,du=-c(\kappa_p).
\end{equation}
Indeed, we may assume $c:=c(\kappa_p)\ge1$, the case $\kappa_p=1$ being
trivial.  Decompose $\Z_p^{\times}$ into the sets
$V_n:=\{u\in\Z_p^{\times}:v_p(1-u)=n\}$, $n\ge0$, on which $|1-u|_p^{-1}=p^{n}$;
one has $V_0=\Z_p^{\times}\setminus(1+p\Z_p)$ and
$V_n=(1+p^{n}\Z_p)\setminus(1+p^{n+1}\Z_p)$ for $n\ge1$.  Write
$\mu_n:=\vol(1+p^{n}\Z_p)=\bigl((p-1)p^{n-1}\bigr)^{-1}$ for $n\ge1$ and
$\mu_0:=\vol(\Z_p^{\times})=1$, and put
\[
A_n:=\int_{1+p^{n}\Z_p}\bigl(\kappa_p(u)-1\bigr)du\quad(n\ge1),\qquad
A_0:=\int_{\Z_p^{\times}}\bigl(\kappa_p(u)-1\bigr)du,
\]
so that $\int_{V_n}(\kappa_p-1)\,du=A_n-A_{n+1}$ for every $n\ge0$.  Since
$\kappa_p$ is trivial on $1+p^{n}\Z_p$ exactly for $n\ge c$, orthogonality of
characters on the group $1+p^{n}\Z_p$ (respectively on $\Z_p^{\times}$ when
$n=0$, using $c\ge1$) gives $\int_{1+p^n\Z_p}\kappa_p\,du=\mu_n$ for $n\ge c$
and $0$ for $n<c$.  Hence $A_n=0$ for $n\ge c$ and $A_n=-\mu_n$ for $n<c$.
Therefore
\[
\int_{\Z_p^{\times}}\frac{\kappa_p(u)-1}{|1-u|_p}\,du
=\sum_{n\ge0}p^{n}\bigl(A_n-A_{n+1}\bigr)
=-\sum_{n=0}^{c-1}p^{n}\mu_n+\sum_{n=0}^{c-2}p^{n}\mu_{n+1}
=-\mu_0+\sum_{m=1}^{c-1}\bigl(p^{m-1}-p^{m}\bigr)\mu_m .
\]
For $m\ge1$ one has $(p^{m-1}-p^{m})\mu_m=p^{m-1}(1-p)/\bigl((p-1)p^{m-1}\bigr)
=-1$, so the right hand side is $-1-(c-1)=-c$, which is \eqref{eq:conductor}.

For $\kappa_p=1$ the packet vanishes identically, since the integrand is
zero.  Combining the two parts gives \eqref{eq:Wphat}, the two cases being exclusive because the shell
part carries the factor $\delta_{\kappa_p,1}$ and the packet the factor
$c(\kappa_p)$, which vanishes when $\kappa_p=1$.  Finally
$\zeta_p=p^{-1/2}e^{i(\tau\log p+\theta_p)}$ with
$\theta_p:=\arg\prod_{q\ne p}\kappa_q(p)$ a constant, so
$2\Rel(\zeta/(1-\zeta))$ is a real-analytic $2\pi/\log p$-periodic function of
$\tau$; the ramified case is constant.

(b)  By \eqref{eq:slices}, $\chi(j_{\infty}(x))=\sgn(x)^{\varepsilon}
|x|^{i\tau}$, so the two profiles of Definition~\ref{def:trace} are
$b_{+}(t)=e^{i\tau t}$ and $b_{-}(t)=(-1)^{\varepsilon}e^{i\tau t}$.  Put
these into \eqref{eq:Winf} (Convention~\ref{conv:Winf}).  The $K_{c}$-integral
is $(-1)^{\varepsilon}\int_0^{\infty}K_c(t)\,2\cos(\tau t)\,dt$, and the
remaining terms form the regularized $K_{s}$-integral, which is the same for
both parities; since $2K_{s}(|t|)=1/(2|\sinh(t/2)|)$ and
$2K_{c}(|t|)=1/(2\cosh(t/2))$, this reads
\begin{equation}\label{eq:Winfhat2}
\Wh_{\infty}(\varepsilon,\tau)=\frac12\int'_{\R}
\frac{e^{i\tau t}\,dt}{2|\sinh(t/2)|}
+(-1)^{\varepsilon}\cdot\frac12\int_{\R}\frac{e^{i\tau t}\,dt}{2\cosh(t/2)} .
\end{equation}
where, by definition,
\[
\frac12\int'_{\R}\frac{e^{i\tau t}\,dt}{2|\sinh(t/2)|}
:=(\gamma+\log4\pi)+\int_0^{\infty}\bigl[2K_{s}(t)\cos\tau t-2K(t)e^{-t/2}\bigr]dt,
\]
\[
\frac12\int_{\R}\frac{e^{i\tau t}\,dt}{2\cosh(t/2)}=\int_0^{\infty}2K_{c}(t)\cos\tau t\,dt .
\]
Both kernels are even in $t$, so both integrals are real, the odd part of
$e^{i\tau t}$ integrating to zero; the regularization at $t=0$, which subtracts
the value $1$ of $\chi\circ j_{\infty}$ at the identity, preserves evenness.
The second integral converges absolutely and equals $\pi/\cosh(\pi\tau)$, hence
is bounded by $\pi$ and tends to $0$; it contributes $O(1)$ uniformly in
$\varepsilon$, and it is the only place where the parity enters.  For the first,
write $1/(2|\sinh(t/2)|)=1/|t|+r(t)$ with $r$ even, bounded on $|t|\le1$ and
exponentially decaying; then
\[
\int'_{\R}\frac{e^{i\tau t}dt}{2|\sinh(t/2)|}
=\int_{|t|\le1}\frac{e^{i\tau t}-1}{|t|}dt+O(1)
=2\int_0^{1}\frac{\cos\tau t-1}{t}dt+O(1)
=-2\operatorname{Cin}(|\tau|)+O(1),
\]
and $\operatorname{Cin}(y)=\int_0^{y}(1-\cos v)v^{-1}dv=\log y+\gamma-
\operatorname{Ci}(y)=\log y+O(1)$; the $O(1)$'s are uniform in $\tau$, by
absolute convergence away from $t=0$ and by the Riemann--Lebesgue lemma for the
tail.  Multiplying by the factor $\tfrac12$ of \eqref{eq:archmeasure} gives
\eqref{eq:Winfhat}, uniformly in $\varepsilon$; the parity never enters the
divergence because the singularity of $\tau_{\infty}$ sits at $x=+1$, where
$\sgn(+1)^{\varepsilon}=1$ for both parities.

\emph{The closed form.}  We derive \eqref{eq:Winfclosed} from
Definition~\ref{def:trace}.  For $\varepsilon=0$ the two profiles coincide,
$b_{+}=b_{-}=b$ with $b(t)=e^{i\tau t}$, and $b(0)=1$; so the even form
\eqref{eq:Winfeven} gives
\[
\Wh_{\infty}(0,\tau)=(\gamma+\log4\pi)+\int_0^{\infty}K(t)
\bigl[2\cos\tau t-2e^{-t/2}\bigr]dt,\qquad
K(t)=\frac{e^{-t/2}}{1-e^{-2t}} .
\]
Substituting $u=2t$ turns the integral into
\[
\mathcal J(\tau):=\int_0^{\infty}
\frac{e^{-u/4}\cos(\tau u/2)-e^{-u/2}}{1-e^{-u}}\,du
=\int_0^{\infty}\frac{\tfrac12\bigl(e^{-\alpha u}+e^{-\bar\alpha u}\bigr)
-e^{-u/2}}{1-e^{-u}}\,du,\qquad \alpha:=\tfrac14+\tfrac{i\tau}{2},
\]
each integrand being continuous at $u=0$ and exponentially small at infinity,
so that the integral converges absolutely.  Gauss's integral representation
of the digamma function \cite[\S5.9]{DLMF},
\[
\psi(s)=-\gamma+\int_0^{\infty}\frac{e^{-u}-e^{-su}}{1-e^{-u}}\,du
\qquad(\Rel s>0),
\]
gives, for $\Rel s>0$,
$\int_0^{\infty}\bigl(e^{-su}-e^{-u/2}\bigr)(1-e^{-u})^{-1}du
=\psi(\tfrac12)-\psi(s)$, by subtracting the two representations.  Applying
this with $s=\alpha$ and with $s=\bar\alpha$, and halving,
\[
\mathcal J(\tau)=\psi(\tfrac12)-\Rel\,\psi\bigl(\tfrac14+\tfrac{i\tau}{2}\bigr),
\qquad \psi(\tfrac12)=-\gamma-2\log2 .
\]
Hence $\Wh_{\infty}(0,\tau)=\gamma+\log4\pi-\gamma-2\log2-
\Rel\,\psi(\tfrac14+\tfrac{i\tau}{2})=\log\pi-
\Rel\,\psi(\tfrac14+\tfrac{i\tau}{2})$, which is \eqref{eq:Winfclosed}.  This
is the archimedean term of Weil's explicit formula in the present
normalization \cite{Weil1952,CCWeilArch}; the derivation above fixes the
normalization, and no comparison of conventions is required.  Since
$\Rel\,\psi(\tfrac14+\tfrac{i\tau}{2})=\log|\tau/2|+O(\tau^{-2})$, the closed
form agrees with \eqref{eq:Winfhat}, including the factor $\tfrac12$.

\emph{The odd component.}  Subtracting the two cases of \eqref{eq:Winfhat2},
$\Wh_{\infty}(1,\tau)=\Wh_{\infty}(0,\tau)-\int_{\R}e^{i\tau t}(2\cosh(t/2))^{-1}dt
=\Wh_{\infty}(0,\tau)-\pi/\cosh(\pi\tau)$.  The reflection formula
$\psi(1-z)-\psi(z)=\pi\cot\pi z$ at $z=\tfrac14+\tfrac{i\tau}2$, together with
$\Rel\cot(\tfrac\pi4+\tfrac{i\pi\tau}2)=\sin(\tfrac\pi2)/(\cosh(\pi\tau)-\cos\tfrac\pi2)
=1/\cosh(\pi\tau)$ and $\Rel\psi(\bar w)=\Rel\psi(w)$, gives
$\Rel\psi(\tfrac34+\tfrac{i\tau}2)-\Rel\psi(\tfrac14+\tfrac{i\tau}2)=\pi/\cosh(\pi\tau)$.
Hence
\begin{equation}\label{eq:Winfodd}
\Wh_{\infty}(1,\tau)=\log\pi-\Rel\,\psi\bigl(\tfrac34+\tfrac{i\tau}{2}\bigr),
\end{equation}
the classical archimedean factor of odd parity, now derived from
Convention~\ref{conv:Winf}.

\emph{Monotonicity and the maximum.}  By the Gauss series
$\psi(z)=-\gamma+\sum_{n\ge0}\bigl[\tfrac1{n+1}-\tfrac1{n+z}\bigr]$, locally
uniformly convergent off $-\N_0$,
\[
\Rel\,\psi\bigl(\tfrac14+\tfrac{i\tau}{2}\bigr)
=-\gamma+\sum_{n\ge0}\Bigl[\frac1{n+1}
-\frac{n+\tfrac14}{(n+\tfrac14)^{2}+\tau^{2}/4}\Bigr].
\]
Each subtracted term is a positive decreasing function of $|\tau|$, so the sum
is increasing in $|\tau|$; hence $\Rel\,\psi(\tfrac14+\tfrac{i\tau}{2})$ is
increasing in $|\tau|$ and $\Wh_{\infty}(0,\cdot)$ attains its maximum at
$\tau=0$:
\begin{equation}\label{eq:maxarch}
\max_{\tau}\Wh_{\infty}(0,\tau)=\log\pi-\psi(\tfrac14)
=\log\pi+\gamma+3\log2+\tfrac{\pi}{2}\approx5.372 ,
\end{equation}
the value $\psi(\tfrac14)=-\gamma-3\log2-\tfrac{\pi}{2}$ being Gauss's digamma
theorem \cite[\S5.4(iii)]{DLMF}.  For $\varepsilon=1$ the same two displays
bound $\Wh_{\infty}(1,\cdot)$ above by $\max_{\tau}\Wh_{\infty}(0,\tau)+\pi/2$,
the parity term $(-1)^{\varepsilon}\tfrac12\pi/\cosh(\pi\tau)$ of
\eqref{eq:Winfhat2} being bounded by $\pi/2$ in absolute value.
Either way $\Wh_{\infty}$ is bounded above.

(c)  $S_f$ is finite and each $\Wh_p$ is real and bounded by (a), while
$\Wh_{\infty}$ is real, bounded above and tends to $-\infty$ by (b); so $\Th$
is real, bounded above by a finite $m$, and tends to $-\infty$ along every
component.  Adding a real constant preserves all of this, so $\Th+c$ is not the
zero function on any component; and a real-analytic function on $\R$ which is
not identically zero has isolated zeros.

(d)  Insert the Fourier inversion $H(w)=\int\widehat H(\chi)\chi(w)d\chi$ into
\eqref{eq:tracedef} and interchange; the term $cH(1)$ becomes
$c\int\widehat H\,d\chi$, since $\delta_1$ has symbol identically $1$.  The
interchange is justified term by term: $\widehat H$ is supported on finitely many
components and Schwartz in $\tau$ there; the shell sums are dominated by
$\sum_kp^{-|k|/2}$ uniformly in $\chi$; the unit packets are dominated by
$c(\kappa_p)\log p$, finite on the finitely many active components by
\eqref{eq:conductor}; and $\Wh_{\infty}$ grows only logarithmically by (b).
The product of a Schwartz function with a logarithmically growing one is
integrable.
\end{proof}

\begin{remark}\label{rem:symbolshape}
Expanding \eqref{eq:Wphat} at an unramified place gives
\begin{equation}\label{eq:symbolexp}
\Wh_p(\chi)=\sum_{k\ge1}2(\log p)\,p^{-k/2}
\cos\bigl(k\tau\log p+k\theta_p\bigr).
\end{equation}
Both ingredients are derived rather than assumed: the weights $p^{-|k|/2}$ are
the orbital integrals $\tau_p(p^{k})$ of Lemma~\ref{lem:tau}, and $\log p$ is
the shell volume fixed in Definition~\ref{def:trace}.  Theorem~\ref{thm:symbol}
is a statement about the \emph{size} of a symbol: real-analyticity, boundedness
of the finite-place terms, and divergence of the archimedean term.  The
arguments of Section~\ref{sec:numerical} use nothing else, and would apply
verbatim to any symbol with a bounded finite part and a dominant logarithmic
archimedean factor.
\end{remark}

\subsection{Mittag--Leffler form and the Tate poles}

\begin{theorem}\label{thm:ML}
Fix a component $\mathfrak c=(\varepsilon,(\kappa_q))$ of $\Ch$ and write
$\tau$ for its real coordinate.
\begin{enumerate}
\item[(a)] Let $p\in S_f$ with $\kappa_p=1$, and put $L:=\log p$ and
$\theta_p:=\arg\prod_{q\ne p}\kappa_q(p)$.  Then $\Wh_p(\mathfrak c,\cdot)$
extends from $\R$ to a meromorphic function on $\C$, $2\pi/L$-periodic in
$\tau$, given by
\begin{equation}\label{eq:MLp}
\Wh_p(\tau)=-\log p-i\sum_{k\in\Z}^{\mathrm{sym}}
\Bigl[\frac{1}{\tau-\tau_k^{+}}-\frac{1}{\tau-\tau_k^{-}}\Bigr],
\qquad \tau_k^{\pm}:=\pm\frac{i}{2}+\frac{2\pi k-\theta_p}{\log p},
\end{equation}
the sum taken symmetrically and converging locally uniformly off the poles.
All poles are simple; the residue at $\tau_k^{+}$ is $-i$ and at $\tau_k^{-}$
is $+i$, for every $k$ and every $p$.  For $\kappa_p\ne1$ the symbol is the
constant $-c(\kappa_p)\log p$, hence entire.
\item[(b)] With $\omega_n:=i(2n+\tfrac12)$, $n\ge0$,
\begin{equation}\label{eq:MLinf}
\Wh_{\infty}(0,\tau)=\log\pi+\gamma+\sum_{n\ge0}
\Bigl[\frac{-i}{\tau-\omega_n}+\frac{i}{\tau+\omega_n}-\frac{1}{n+1}\Bigr],
\end{equation}
locally uniformly convergent off the poles; $\Wh_{\infty}(0,\cdot)$ is
meromorphic on $\C$ with simple poles at $\pm\omega_n$ and residues $\mp i$.
\item[(c)] On the principal component $\mathfrak c_0$ one has $\theta_p=0$ for
every $p$, so $\tau_0^{\pm}=\pm i/2=\pm\omega_0$: every place of $S$
contributes a simple pole at $\tau=\pm i/2$ with residue $\mp i$, and
\begin{equation}\label{eq:residues}
\Res_{\tau=+i/2}\Th(\mathfrak c_0,\cdot)=-i\bigl(|S_f|+1\bigr),\qquad
\Res_{\tau=-i/2}\Th(\mathfrak c_0,\cdot)=+i\bigl(|S_f|+1\bigr).
\end{equation}
\item[(d)] Let $F\in\DA$ be supported on $\mathfrak c_0$ with radial profile
$f\in C_c^{\infty}(\R)$ in $t=\log|w|_S$, so
$\widehat F(\tau)=\int f(t)e^{-i\tau t}dt$.  Then for every $h>\tfrac12$
\begin{equation}\label{eq:contourp}
W_p(F)=\sum_{k\in\Z}\widehat F\Bigl(\frac{i}{2}+\frac{2\pi k}{\log p}\Bigr)
+\frac{1}{2\pi}\int_{\Im\tau=h}\widehat F(\tau)\Wh_p(\tau)\,d\tau ,
\end{equation}
both terms absolutely convergent, and the sum is the Poisson periodization
\begin{equation}\label{eq:poisson}
\sum_{k\in\Z}\widehat F\Bigl(\frac i2+\frac{2\pi k}{\log p}\Bigr)
=(\log p)\sum_{m\in\Z}f(m\log p)\,p^{m/2}.
\end{equation}
\item[(e)] For the same $F$ and any $\tfrac12<h<\tfrac52$,
\begin{equation}\label{eq:contourinf}
W_{\infty}(F)=\widehat F(i/2)
+\frac{1}{2\pi}\int_{\Im\tau=h}\widehat F(\tau)\Wh_{\infty}(0,\tau)\,d\tau
=\lambda_+(F)+\frac{1}{2\pi}\int_{\Im\tau=h}\widehat F\,\Wh_{\infty}\,d\tau .
\end{equation}
\end{enumerate}
\end{theorem}

\begin{proof}
(a)  By \eqref{eq:Wphat} with
$\zeta=\zeta_p(\chi)=e^{-L/2+i(\tau L+\theta_p)}$ one has, for real $\tau$,
\begin{equation}\label{eq:twoexp}
\Wh_p=2L\,\Rel\frac{\zeta}{1-\zeta}
=L\Bigl[\frac{1}{1-e^{-w}}+\frac{1}{1-e^{-\tilde w}}-2\Bigr],\qquad
w:=\tfrac L2-i(\tau L+\theta_p),\quad \tilde w:=\tfrac L2+i(\tau L+\theta_p),
\end{equation}
since $\zeta/(1-\zeta)=1/(1-\zeta)-1$ and $\bar\zeta=e^{-\tilde w}$ for real
$\tau$.  The right hand side is meromorphic in $\tau\in\C$ and
$2\pi/L$-periodic, since $\tau\mapsto\tau+2\pi/L$ sends $w\mapsto w-2\pi i$ and
$\tilde w\mapsto\tilde w+2\pi i$.  Apply the Mittag--Leffler expansion
\begin{equation}\label{eq:coth}
\frac{1}{1-e^{-z}}=\frac12+\frac12\coth\frac z2
=\frac12+\sum_{k\in\Z}^{\mathrm{sym}}\frac{1}{z-2\pi ik},
\end{equation}
valid for $z\notin2\pi i\Z$ with locally uniform symmetric convergence, the
poles of $\tfrac12\coth(z/2)$ being simple at $2\pi ik$ with residue $1$ and
the pairs $\pm k$ combining to $2z/(z^{2}+4\pi^{2}k^{2})$.  Substituting
$z=w$ and $z=\tilde w$,
\[
\Wh_p=L\Bigl[1+\sum_k^{\mathrm{sym}}\Bigl(\frac1{w-2\pi ik}
+\frac1{\tilde w-2\pi ik}\Bigr)-2\Bigr]
=-L+L\sum_k^{\mathrm{sym}}\Bigl(\frac1{w-2\pi ik}+\frac1{\tilde w-2\pi ik}
\Bigr).
\]
Now $\tilde w-2\pi ik=iL(\tau-\tau_k^{+})$, because
$iL\bigl(\tau-\tfrac i2-\tfrac{2\pi k-\theta_p}{L}\bigr)
=iL\tau+\tfrac L2+i\theta_p-2\pi ik=\tilde w-2\pi ik$; and
$w-2\pi ik=-iL(\tau-\tau_{-k}^{-})$, because
$-iL\bigl(\tau+\tfrac i2+\tfrac{2\pi k+\theta_p}{L}\bigr)
=-iL\tau+\tfrac L2-i\theta_p-2\pi ik=w-2\pi ik$.  Hence
$L/(\tilde w-2\pi ik)=-i/(\tau-\tau_k^{+})$ and
$L/(w-2\pi ik)=+i/(\tau-\tau_{-k}^{-})$; re-indexing $k\mapsto-k$ in the second
family, which is legitimate inside a symmetric sum, gives \eqref{eq:MLp}, and
the residues are read off the displayed simple poles.

(b)  By \eqref{eq:Winfclosed}, $\Wh_{\infty}(0,\tau)=\log\pi-\Rel\,
\psi(\tfrac14+\tfrac{i\tau}2)$ for real $\tau$.  Use the Gauss series
$\psi(z)=-\gamma+\sum_{n\ge0}\bigl[\tfrac1{n+1}-\tfrac1{n+z}\bigr]$, locally
uniformly convergent off $-\N_0$.  For real $\tau$,
$\Rel\,\psi(\tfrac14+\tfrac{i\tau}2)=\tfrac12\bigl[\psi(\tfrac14+\tfrac{i\tau}2)
+\psi(\tfrac14-\tfrac{i\tau}2)\bigr]$, and both terms are meromorphic in
$\tau$, so the right hand side of \eqref{eq:MLinf} is the analytic
continuation.  Since $n+\tfrac14+\tfrac{i\tau}2=\tfrac i2(\tau-\omega_n)$ and
$n+\tfrac14-\tfrac{i\tau}2=-\tfrac i2(\tau+\omega_n)$, one has
$\bigl(n+\tfrac14+\tfrac{i\tau}2\bigr)^{-1}=-2i(\tau-\omega_n)^{-1}$ and
$\bigl(n+\tfrac14-\tfrac{i\tau}2\bigr)^{-1}=2i(\tau+\omega_n)^{-1}$, whence
\[
\tfrac12\bigl[\psi(\tfrac14+\tfrac{i\tau}2)+\psi(\tfrac14-\tfrac{i\tau}2)\bigr]
=-\gamma+\sum_{n\ge0}\Bigl[\frac1{n+1}+\frac{i}{\tau-\omega_n}
-\frac{i}{\tau+\omega_n}\Bigr],
\]
and subtracting from $\log\pi$ gives \eqref{eq:MLinf}.  The bracketed terms are
$O(n^{-2})$ for fixed $\tau$, so the series converges absolutely and locally
uniformly off the poles.

(c)  $\tau_0^{\pm}=\pm i/2-\theta_p/L$ equals $\pm i/2$ if and only if
$\theta_p\equiv0\pmod{2\pi}$, and on $\mathfrak c_0$ every $\kappa_q$ is
trivial, so $\theta_p=0$; and $\omega_0=i/2$.  Summing the residues of the
$|S_f|$ finite terms and the archimedean one gives \eqref{eq:residues}.

(d)  Restricting \eqref{eq:plancherelv} to the single component $\mathfrak c_0$
gives $W_p(F)=\tfrac1{2\pi}\int_{\R}\widehat F\Wh_p\,d\tau$, the other
components not meeting $\supp\widehat F$.  By Paley--Wiener, $\widehat F$ is
entire of exponential type with $|\widehat F(\sigma+iy)|\le
C_Ne^{R|y|}(1+|\sigma|)^{-N}$ for every $N$, $R$ being the radius of $\supp f$.
Fix $h>\tfrac12$.  In the closed strip $0\le\Im\tau\le h$ the poles of
$\Wh_p$ are exactly the points $i/2+2\pi k/L$, $k\in\Z$, all at height
$\tfrac12$ by \eqref{eq:MLp} with $\theta_p=0$, the family $\tau_k^{-}$ lying
at height $-\tfrac12$.  Put $R_j:=2\pi(j+\tfrac12)/L$ and integrate
$\widehat F\Wh_p$ over the positively oriented boundary of the rectangle
$[-R_j,R_j]\times[0,h]$.  By $2\pi/L$-periodicity in $\Re\tau$, the modulus of
$\Wh_p$ is bounded on the two vertical sides by a constant independent of $j$,
each side being a translate by an integer multiple of $2\pi/L$ of the fixed
compact segment $\{\Re\tau=\pi/L,\ 0\le\Im\tau\le h\}$, which avoids all poles;
with the rapid decay of $\widehat F$ in $\Re\tau$, uniform for
$0\le\Im\tau\le h$, the two side contributions tend to $0$.  On the two
horizontal sides $\Wh_p$ is bounded: for $\Im\tau=y$ one has
$|e^{-w}|=e^{-L/2-yL}$ and $|e^{-\tilde w}|=e^{-L/2+yL}$ by \eqref{eq:twoexp},
so at $y=0$ both moduli are $e^{-L/2}<1$, and at $y=h>\tfrac12$ the first is
$<e^{-L}<1$ while the second is $>e^{L(h-1/2)}>1$; in both cases $|1-e^{-w}|$
and $|1-e^{-\tilde w}|$ are bounded below along the whole line.  Letting
$j\to\infty$, the residue theorem gives
\[
\int_{\Im\tau=0}\widehat F\Wh_p\,d\tau-\int_{\Im\tau=h}\widehat F\Wh_p\,d\tau
=2\pi i\sum_{k\in\Z}\Res_{\tau=i/2+2\pi k/L}\bigl(\widehat F\Wh_p\bigr)
=2\pi i\sum_{k\in\Z}(-i)\widehat F\Bigl(\frac i2+\frac{2\pi k}{L}\Bigr),
\]
the sum converging absolutely by the Paley--Wiener bound at height $\tfrac12$.
Dividing by $2\pi$ gives \eqref{eq:contourp}.  For \eqref{eq:poisson} set
$g(t):=f(t)e^{t/2}\in C_c^{\infty}(\R)$; then
$\widehat F(i/2+2\pi k/L)=\int f(t)e^{t/2}e^{-2\pi ikt/L}dt=\widehat g(2\pi k/L)$,
and Poisson summation for $g\in C_c^{\infty}$ gives
$\sum_k\widehat g(2\pi k/L)=L\sum_mg(mL)$.

(e)  The same contour argument with $\Wh_{\infty}$ in place of $\Wh_p$.  In the
strip $0\le\Im\tau\le h<\tfrac52$ the only pole is $\omega_0=i/2$, of residue
$-i$, so the residue term is $\tfrac1{2\pi}\cdot2\pi i\cdot(-i)\widehat
F(i/2)=\widehat F(i/2)$, which is $\lambda_+(F)$ by \eqref{eq:lambdaFT}.
Boundedness of $\Wh_{\infty}$ on the two horizontal lines and decay on the
vertical sides follow from $\psi(z)=\log z+O(|z|^{-1})$ uniformly for
$|\arg z|\le\pi-\delta$, so that $\Wh_{\infty}$ grows only logarithmically in
$|\Re\tau|$ on the strip away from the poles, which is dominated by the rapid
decay of $\widehat F$.
\end{proof}

\begin{remark}\label{rem:MLconseq}
Three consequences are used below.  First, the two functionals
$\lambda_{\mp}(F)=\widehat F(\mp i/2)$ are exactly the residue contributions of
the symbol at the poles $\pm i/2$; \eqref{eq:contourinf} exhibits $\lambda_+$
as a summand of the archimedean local trace.  Second, at a finite place the
residue term is not $\lambda_+$ but its $2\pi/\log p$-periodization, which by
\eqref{eq:poisson} is the $p$-adic periodization of the Tate-twisted profile.
Third, by \eqref{eq:residues} the total residue at $\pm i/2$ is
$\mp i(|S_f|+1)$: each place of $S$ carries its own copy of the Tate pole, so
there is no \emph{chart-independent} normalization of the scalar form in which
the pair $(\lambda_-,\lambda_+)$ appears with hyperbolic coefficient $1$.  This
is the analytic origin of Theorem~\ref{thm:scalartrivial}.  Parts (b), (d) and
(e) are proved for the even component only, which is all that is used: on the
odd component the archimedean symbol involves $\psi(\tfrac34+\tfrac{i\tau}2)$,
whose poles lie at $i(2n+\tfrac32)$, so there is no Tate pole at $\pm i/2$
there --- consistently with \eqref{eq:lambdaFT}, which places $\lambda_{\mp}$
on $\mathfrak c_0$.
\end{remark}

\section{The numerical theory}\label{sec:numerical}

\subsection{The pairing and its radical}

\begin{definition}\label{def:pairing}
For $F,G\in\DA$ put
\begin{equation}\label{eq:IA}
I_A(F,G):=T\bigl(F\star G^{\dagger}\bigr),
\end{equation}
and let
\begin{equation}\label{eq:radical}
\Nrad:=\bigl\{F\in\DA:\ I_A(F,G)=0\ \text{for every }G\in\DA\bigr\},
\qquad N_A:=\DA/\Nrad,
\end{equation}
with the quotient topology.
\end{definition}

\begin{proposition}\label{prop:bilin}
$I_A$ is linear in the first variable, conjugate-linear in the second,
Hermitian, i.e.\ $\overline{I_A(F,G)}=I_A(G,F)$, and jointly continuous on
$\DA\times\DA$.  Moreover
\begin{equation}\label{eq:daggerreal}
T(H^{\dagger})=\overline{T(H)}\qquad(H\in\DA).
\end{equation}
\end{proposition}

\begin{proof}
Linearity and conjugate-linearity are those of $(F,G)\mapsto F\star G^{\dagger}$
together with linearity of $T$.  For \eqref{eq:daggerreal}: by
Theorem~\ref{thm:absconv} applied after subtracting the identity counterterms,
\[
T(H^{\dagger})=\sum_v\int'\overline{H(j_v(x^{-1}))}\,\tau_v(x)\dstar x
+\bigl(\text{real multiple of }H^{\dagger}(1)\bigr),
\]
and the substitution $x\mapsto x^{-1}$, which preserves multiplicative Haar
measure, together with $\tau_v(x^{-1})=\tau_v(x)$ (Lemma~\ref{lem:inv}), turns
the sum into $\overline{\sum_v\int'H(j_v(x))\tau_v(x)\dstar x}$; the identity
counterterms are real multiples of $H(1)$ and satisfy
$H^{\dagger}(1)=\overline{H(1)}$, so they transform the same way.  Hermitian
symmetry: $\DA$ is commutative, so
$F\star G^{\dagger}=\bigl(G\star F^{\dagger}\bigr)^{\dagger}$, whence
$I_A(F,G)=\overline{I_A(G,F)}$ by \eqref{eq:daggerreal}.  Joint continuity:
$\star$ and $\dagger$ are continuous on each fixed-support step of $\DA$, and
$T$ is continuous on each step.
\end{proof}

\begin{proposition}\label{prop:ideal}
$\Nrad$ is a closed two-sided $\star$-ideal, stable under $\dagger$, and $N_A$
is a Hausdorff locally convex space carrying continuous descended operations
$\star$, $\dagger$ and a nondegenerate Hermitian pairing.  Moreover $T$
vanishes on $\Nrad$.
\end{proposition}

\begin{proof}
\emph{Closed.}  $\Nrad=\bigcap_G\ker I_A(\cdot,G)$, an intersection of kernels
of continuous functionals.

\emph{Ideal.}  For $F\in\Nrad$ and $G,H\in\DA$, using commutativity and
$(H\star G^{\dagger})^{\dagger}=G\star H^{\dagger}$,
\[
I_A(F\star H,G)=T\bigl(F\star H\star G^{\dagger}\bigr)
=T\bigl(F\star(G\star H^{\dagger})^{\dagger}\bigr)=I_A(F,G\star H^{\dagger})=0 .
\]

\emph{Stability under $\dagger$.}  For $F\in\Nrad$ and any $G$, using
$(a\star b)^{\dagger}=b^{\dagger}\star a^{\dagger}$, commutativity and
\eqref{eq:daggerreal},
\[
I_A(F^{\dagger},G)=T\bigl(F^{\dagger}\star G^{\dagger}\bigr)
=T\bigl((F\star G)^{\dagger}\bigr)=\overline{T(F\star G)}
=\overline{I_A(F,G^{\dagger})}=0 .
\]

\emph{$T$ vanishes on $\Nrad$.}  Let $F\in\Nrad$ and let $(e_{\varepsilon})$ be
the approximate identity of Remark~\ref{rem:approxid}.  Then
$F\star e_{\varepsilon}^{\dagger}\to F$ in the fixed-support Fr\'echet step
determined by a fixed compact neighbourhood of $\supp F$, on which $T$ is
continuous; hence $T(F)=\lim_{\varepsilon}I_A(F,e_{\varepsilon})=0$.

\emph{Quotient.}  The quotient of a locally convex space by a closed subspace
is Hausdorff; the descended pairing is nondegenerate by the definition of
$\Nrad$; and the descended operations are continuous because $\Nrad$ is a
closed $\star$- and $\dagger$-stable subspace.
\end{proof}

\begin{theorem}\label{thm:radical}
$\Nrad=\{0\}$, for every value of the constant $c$.  Hence $N_A=\DA$, the form
$I_A$ is nondegenerate on the whole test algebra, and $\lambda_{\pm}$
annihilate $\Nrad$ trivially.
\end{theorem}

\begin{proof}
Let $F\in\Nrad$ and $G\in\DA$.  First,
$(G^{\dagger})^{\wedge}(\chi)=\int\overline{G(w^{-1})}\,\overline{\chi(w)}
\dstar w=\overline{\int G(v)\overline{\chi(v)}\dstar v}=\overline{\widehat
G(\chi)}$, and convolution goes to the product of transforms; so
$(F\star G^{\dagger})^{\wedge}=\widehat F\,\overline{\widehat G}$.  By
\eqref{eq:plancherel},
\begin{equation}\label{eq:IAfourier}
I_A(F,G)=\int_{\Ch}\widehat F(\chi)\,\overline{\widehat G(\chi)}\,
\bigl(\Th(\chi)+c\bigr)\,d\chi ,
\end{equation}
and the hypothesis says that this vanishes for every $G$.

Fix a component $\mathfrak c$, given by a parity $\varepsilon$ and a system
$(\kappa_q)$.  Choose $G=\chi_{\mathfrak c}\cdot g$, where
$\chi_{\mathfrak c}=\prod_q\kappa_q$ is regarded as a function on the compact
part, multiplied by the characteristic function of $\prod_q\Z_q^{\times}$, and
$g\in C_c^{\infty}(\R^{\times})$ has parity $\varepsilon$.  Then $G\in\DA$,
$\widehat G$ is supported on $\mathfrak c$, and there it equals
$\widehat g(\tau)$, the Fourier transform of a $C_c^{\infty}$ function in the
logarithmic coordinate.  Thus \eqref{eq:IAfourier} becomes
\[
\int_{\R}\widehat F(\mathfrak c,\tau)\,\overline{\widehat g(\tau)}\,
\bigl(\Th(\mathfrak c,\tau)+c\bigr)\,d\tau=0
\qquad\text{for every } g\in C_c^{\infty}.
\]
The function $m(\tau):=\widehat F(\mathfrak c,\tau)(\Th(\mathfrak c,\tau)+c)$
is locally integrable and is a Schwartz function times a function of at most
logarithmic growth, hence a tempered distribution; the Fourier transforms of
$C_c^{\infty}$ functions are dense in the Schwartz space; therefore $m=0$
almost everywhere.  By Theorem~\ref{thm:symbol}(c) the zeros of
$\Th(\mathfrak c,\cdot)+c$ are isolated, so $\widehat F(\mathfrak c,\tau)=0$ off
a discrete set; and $\widehat F(\mathfrak c,\cdot)$ is continuous, so
$\widehat F\equiv0$ on $\mathfrak c$.  As $\mathfrak c$ was arbitrary,
$\widehat F\equiv0$, hence $F=0$.
\end{proof}

\begin{remark}
The proof uses exactly three properties of the symbol: it is real-analytic on
each component, its finite-place terms are bounded, and its archimedean term
diverges to $-\infty$.  The open constant $c$ is immaterial, since $\Th+c$
inherits all three.
\end{remark}

\subsection{The Tate boundary sequence}

\begin{definition}\label{def:prim}
$\DA^{\mathrm{prim}}:=\ker\lambda_-\cap\ker\lambda_+$, the \emph{primitive}
section space.  Put
$I_{\mathrm{Tate}}(F,G):=\beta_T(MF,MG)=\lambda_-(F)\overline{\lambda_+(G)}
+\lambda_+(F)\overline{\lambda_-(G)}$,
$I_{\mathrm{vec}}:=(I_A,I_{\mathrm{Tate}})$, and
$\Nrad_W:=\{F:I_{\mathrm{vec}}(F,G)=0\ \forall G\}$.
\end{definition}

\begin{proposition}\label{prop:NW}
$\Nrad_W=\Nrad\cap\DA^{\mathrm{prim}}$, it is a closed two-sided
$\star$-ideal stable under $\dagger$, and it is zero.
\end{proposition}

\begin{proof}
If $F\in\Nrad_W$ then the first component gives $F\in\Nrad$, and the second
says $\beta_T(MF,MG)=0$ for every $G$; since $M$ is surjective and $\beta_T$
nondegenerate (Theorem~\ref{thm:tate}), $MF=0$, i.e.\
$F\in\DA^{\mathrm{prim}}$.  The converse is immediate.  $\Nrad$ has the stated
properties by Proposition~\ref{prop:ideal}; $\DA^{\mathrm{prim}}$ is closed
because $\lambda_{\pm}$ are continuous, is an ideal because $\lambda_{\mp}$ are
multiplicative, and is $\dagger$-stable because
$\lambda_{\mp}(F^{\dagger})=\overline{\lambda_{\pm}(F)}$, so that $MF=0$ if and
only if $MF^{\dagger}=0$.  The intersection inherits all three, and it vanishes
by Theorem~\ref{thm:radical}.
\end{proof}

\begin{theorem}\label{thm:tateseq}
There is a canonical exact sequence of topological $\DA$-modules with
continuous $\star$ and $\dagger$,
\begin{equation}\label{eq:tateseq}
0\longrightarrow \DA^{\mathrm{prim}}\longrightarrow N_A
\stackrel{M}{\longrightarrow}\BT\longrightarrow 0,\qquad N_A=\DA,\quad
\BT=\C^{2}\ \text{with}\
\beta_T=\begin{pmatrix}0&1\\1&0\end{pmatrix},
\end{equation}
in which $\dagger$ exchanges the two basis directions $e_{\mp}$ of $\BT$, the
external classes act on $\BT$ by the two characters $|u|_S^{\mp1/2}$ and on the
kernel by translation.  The kernel $\DA^{\mathrm{prim}}$ is nonzero, and $I_A$
is nondegenerate on it.
\end{theorem}

\begin{proof}
Exactness and surjectivity are Theorem~\ref{thm:tate} together with
Definition~\ref{def:prim} and $N_A=\DA$ (Theorem~\ref{thm:radical}); the
exchange of $e_{\mp}$ is $\lambda_{\mp}(F^{\dagger})=\overline{\lambda_{\pm}(F)}$;
the module structure and the form are \eqref{eq:rhoT} and
Theorem~\ref{thm:tate}.

\emph{The kernel is nonzero.}  Let $F,F'$ be the sections of
Theorem~\ref{thm:separate} and put $\Phi:=F-F'$.  Then $\lambda_{\mp}(\Phi)=0$,
so $M\Phi=0$ and $\Phi\in\DA^{\mathrm{prim}}$; and
$T(\Phi)=\nu(\log p)(p^{j}-p^{j'})\ne0$ while $T$ vanishes on $\Nrad$
(Proposition~\ref{prop:ideal}), so $\Phi\ne0$ in $N_A$.  Thus $[\Phi]$ is a
nonzero primitive class.

\emph{Nondegeneracy on the kernel.}  Let $\Phi\in\DA^{\mathrm{prim}}$ satisfy
$I_A(\Phi,G)=0$ for every $G\in\DA^{\mathrm{prim}}$.  The conjugate-linear
functional $G\mapsto I_A(\Phi,G)$ then vanishes on the closed subspace
$\DA^{\mathrm{prim}}$, which has codimension two, hence equals
$a\,\overline{\lambda_-}+b\,\overline{\lambda_+}$ for scalars $a,b$.  Test on
translates: for $w\in C_S$ put $G_w(v):=G(w^{-1}v)$, so
$\widehat{G_w}(\chi)=\overline{\chi(w)}\widehat G(\chi)$ and
$\lambda_{\mp}(G_w)=|w|_S^{\mp1/2}\lambda_{\mp}(G)$.  By
\eqref{eq:IAfourier},
\[
I_A(\Phi,G_w)=\int_{\Ch}\widehat\Phi(\chi)\overline{\widehat G(\chi)}\,
\chi(w)\,\bigl(\Th(\chi)+c\bigr)d\chi ,
\]
which is the inverse Fourier transform at $w$ of the integrable measure
$\widehat\Phi\overline{\widehat G}(\Th+c)\,d\chi$ --- a Schwartz function
times a logarithmically growing one on finitely many components --- and is
therefore \emph{bounded} in $w$.  The right hand side is
$a|w|_S^{-1/2}\overline{\lambda_-(G)}+b|w|_S^{1/2}\overline{\lambda_+(G)}$,
unbounded as $|w|_S\to0$ unless $a\overline{\lambda_-(G)}=0$, and unbounded as
$|w|_S\to\infty$ unless $b\overline{\lambda_+(G)}=0$.  Choosing $G\in\DA$ with
$G\ge0$ and $G\ne0$, so that $\lambda_{\pm}(G)>0$, forces $a=b=0$.  Then
$I_A(\Phi,\cdot)\equiv0$ on all of $\DA$, so $\Phi\in\Nrad=\{0\}$.
\end{proof}

\begin{theorem}\label{thm:nosplit}
\begin{enumerate}
\item[(a)] There is no $\DA$-module map $s:\BT\to N_A$ with $M\circ s=\id$,
where $\DA$ acts on $\BT$ through $M$.
\item[(b)] The $I_A$-orthogonal complement of $\DA^{\mathrm{prim}}$ in $\DA$ is
zero, and the $I_{\mathrm{vec}}$-orthogonal complement coincides with it.
\end{enumerate}
\end{theorem}

\begin{proof}
(a)  Suppose $s$ exists and put $\Phi_{\mp}:=s(e_{\mp})\in\DA$; then
$M(\Phi_{\mp})=e_{\mp}\ne0$, so $\Phi_{\mp}\ne0$, and $\DA$-linearity reads
\begin{equation}\label{eq:eigen}
a\star\Phi_{\mp}=\lambda_{\mp}(a)\,\Phi_{\mp}\qquad\text{for every }a\in\DA .
\end{equation}
Taking Fourier transforms, $\widehat a(\chi)\widehat\Phi_{\mp}(\chi)=
\lambda_{\mp}(a)\widehat\Phi_{\mp}(\chi)$ for every $\chi\in\Ch$ and every $a$.
Suppose $\widehat\Phi_{\mp}(\chi_0)\ne0$ for some $\chi_0$; then
$\widehat a(\chi_0)=\lambda_{\mp}(a)$ for every $a\in\DA$.  But
$\lambda_{\mp}(a)=\widehat a(\mp i/2)$ by \eqref{eq:lambdaFT}, a value of the
analytic continuation of $\widehat a$ \emph{off} the unitary axis, whereas
$\widehat a(\chi_0)$ is a value \emph{on} it; and the two functionals are
independent.  Indeed, for the translate $a_w(v)=a(w^{-1}v)$ one has
$|\widehat{a_w}(\chi_0)|=|\widehat a(\chi_0)|$, constant in $w$, while
$|\lambda_{\mp}(a_w)|=|w|_S^{\mp1/2}|\lambda_{\mp}(a)|$, which is unbounded in
$w$ whenever $\lambda_{\mp}(a)\ne0$.  Taking $a\ge0$, $a\ne0$, so that
$\lambda_{\mp}(a)>0$, and translating until
$|\lambda_{\mp}(a_w)|>|\widehat a(\chi_0)|$, we reach a contradiction.  Hence
$\widehat\Phi_{\mp}\equiv0$, so $\Phi_{\mp}=0$, contradicting
$M(\Phi_{\mp})=e_{\mp}$.

(b)  Let $F$ be $I_A$-orthogonal to $\DA^{\mathrm{prim}}$.  Exactly as in the
last paragraph of the proof of Theorem~\ref{thm:tateseq}, the functional
$G\mapsto I_A(F,G)$ equals $a\overline{\lambda_-}+b\overline{\lambda_+}$, is
bounded on translates, and the right hand side is unbounded unless $a=b=0$;
so $I_A(F,\cdot)\equiv0$ and $F\in\Nrad=\{0\}$.  For $I_{\mathrm{vec}}$, note
that $I_{\mathrm{Tate}}(F,G)=\beta_T(MF,MG)=0$ automatically when
$G\in\DA^{\mathrm{prim}}$, so the two orthogonal complements agree.
\end{proof}

\begin{remark}\label{rem:weakhyp}
By Theorem~\ref{thm:symbol}(c) and Theorem~\ref{thm:signature} below, for
$c>-m$ the Hermitian form $I_A$ is indefinite, so isotropic vectors exist and
one may choose $F,G$ with $I_A(F,F)=I_A(G,G)=0$ and $I_A(F,G)=1$: an abstract
hyperbolic plane inside $N_A$.  This is \emph{not} the ruling plane.  It is not
canonical, depending on an arbitrary choice of isotropic vectors, of which an
indefinite form has a continuum; it carries no arithmetic, the construction
using only a sign change of $\Th+c$; and it is not compatible with the module
structure and does not map onto $\BT$, since by (a) the eigenvalue equation
\eqref{eq:eigen} has no solutions.  The ruling plane is characterized by being
the image of $M$, and Theorem~\ref{thm:nosplit}(a) says precisely that this
characterization is realizable only on a quotient.
\end{remark}

\subsection{The separated completion}

\begin{theorem}\label{thm:completion}
Let $\DA'$ be the strong dual of the nuclear LF space $\DA$ and put
\begin{equation}\label{eq:J}
J:N_A\longrightarrow\DA'\oplus\BT,\qquad
J(F):=\bigl(G\mapsto I_A(F,G),\ M(F)\bigr).
\end{equation}
Then $J$ is injective; $\NAd:=\overline{J(N_A)}$ is complete and nuclear; the
vector pairing $I_{\mathrm{vec}}\bigl((\ell,m),G\bigr)=
\bigl(\ell(G),\beta_T(m,MG)\bigr)$ is separating on $\NAd$; and $\DA$ acts
continuously on $\NAd$ by $(a\cdot\ell)(G)=\ell(a^{\dagger}\star G)$ and
$M(a\star F)=M(a)M(F)$, with $\dagger$ extended by duality.
\end{theorem}

\begin{proof}
Injectivity: $J(F)=0$ gives $F\in\Nrad_W=\{0\}$ by
Proposition~\ref{prop:NW}.  The space $\DA$ is a countable strict inductive
limit of nuclear Fr\'echet spaces; it is bornological, so its strong dual is
complete \cite[IV.6.1, Cor.~1]{Schaefer}, and that dual is nuclear, being the
projective limit of the strong duals of the steps, which are nuclear
(DF)-spaces \cite[III.7.4 and IV.9.6]{Schaefer}, \cite[Prop.~50.6]{Treves}.
Hence $\DA'\oplus\BT$ is complete and nuclear, and a closed subspace of a
complete nuclear space is complete and nuclear.  Separation: if $I_{\mathrm{vec}}((\ell,m),G)=0$ for all
$G$ then $\ell=0$, and surjectivity of $M$ with nondegeneracy of $\beta_T$
gives $m=0$.  The displayed module formulas are continuous for the strong dual
topology.
\end{proof}

\begin{lemma}\label{lem:Nseminorm}
For $G\in\DA$ put
\begin{equation}\label{eq:Nsemi}
N(G):=\int_{\Ch}|\widehat G(\chi)|\bigl(|\Th(\chi)|+|c|\bigr)d\chi .
\end{equation}
Then $N(G)<\infty$ for every $G$; $|I_A(F,G)|\le
\bigl(\sup_{\chi}|\widehat F(\chi)|\bigr)N(G)$ for all $F,G\in\DA$; and $N$ is
bounded on every bounded subset of $\DA$.
\end{lemma}

\begin{proof}
A bounded set $B\subseteq\DA$ lies in one Fr\'echet step: there are a compact
$K\subseteq C_S$ and an integer $N_0\ge0$ such that every $G\in B$ is supported
in $K$, invariant under $U_{N_0}=\prod_q(1+q^{N_0}\Z_q)$, and has all step
seminorms bounded by $C_B$.  Invariance forces $\widehat G(\varepsilon,\tau;
(\kappa_q))=0$ unless every $\kappa_q$ is trivial on $1+q^{N_0}\Z_q$, and
there are only finitely many such systems; with the two parities this leaves
finitely many components, say $d$ of them.  On each,
$\widehat G(\mathfrak c,\tau)=\int G_{\mathfrak c}(t)e^{-i\tau t}dt$ with
$G_{\mathfrak c}$ supported in the fixed compact logarithmic window of $K$ and
with $|G_{\mathfrak c}|,|G_{\mathfrak c}''|\le C'_B$; integrating by parts twice
gives $|\widehat G(\mathfrak c,\tau)|\le C''_B\min(1,\tau^{-2})$ uniformly in
$\mathfrak c$ and in $G\in B$.  By Theorem~\ref{thm:symbol},
\begin{equation}\label{eq:logbound}
|\Th(\mathfrak c,\tau)|\le A_S+\sum_{p\in S_f}c(\kappa_p)\log p
+B_S\log(2+|\tau|)=:A_{\mathfrak c}+B_S\log(2+|\tau|),
\end{equation}
where $A_S$ and $B_S$ depend only on $S$ (the shell parts, the archimedean
parity term and the archimedean $O(1)$ of \eqref{eq:Winfhat}), while the
ramified unit packets $-c(\kappa_p)\log p$ of \eqref{eq:Wphat} make
$A_{\mathfrak c}$ depend on the component through its conductors.  The bound is
therefore uniform on any finite set of components, and not over all of $\Ch$:
$\sup_{\mathfrak c}A_{\mathfrak c}=\infty$.  Put
$A_B:=\max A_{\mathfrak c}$ over the $d$ components met by $B$.  Hence
\[
N(G)\le d\,C''_B\int_{\R}\min(1,\tau^{-2})
\bigl(A_B+|c|+B_S\log(2+|\tau|)\bigr)\frac{d\tau}{2\pi}<\infty ,
\]
a bound independent of $G\in B$; the case $B=\{G\}$ gives finiteness for a
single $G$.  The inequality is \eqref{eq:IAfourier} with $|\widehat F|$ replaced
by its supremum.
\end{proof}

\subsection{The rulings}

\begin{theorem}\label{thm:rulings}
Fix a real $\vartheta\in C_c^{\infty}((0,\infty))$ with $\vartheta\ge0$,
$\int_{\R}\vartheta(s)\,ds=1$ and $\supp\vartheta\subseteq[a,b]$, $0<a<b<\infty$.  For
$n\ge1$ define $F_n,G_n\in\DA$ by
\begin{equation}\label{eq:spread}
F_n(w):=n^{-1}|w|_S^{-1/2}\vartheta\Bigl(\frac{\log|w|_S}{n}\Bigr),\qquad
G_n(w):=n^{-1}|w|_S^{+1/2}\vartheta\Bigl(-\frac{\log|w|_S}{n}\Bigr),
\end{equation}
functions of $|w|_S$ alone.  Then:
\begin{enumerate}
\item[(a)] $\lambda_+(F_n)=1$ exactly and $0<\lambda_-(F_n)\le\lVert\vartheta
\rVert_{L^1}e^{-na}$; $\lambda_-(G_n)=1$ exactly and
$0<\lambda_+(G_n)\le\lVert\vartheta\rVert_{L^1}e^{-na}$.  Hence $M(F_n)\to e_+$ and
$M(G_n)\to e_-$ in $\BT$.
\item[(b)] $\widehat F_n$ and $\widehat G_n$ are supported on the principal
component $\mathfrak c_0$, and
\begin{equation}\label{eq:supFhat}
\sup_{\chi}|\widehat F_n(\chi)|\le\lVert\vartheta\rVert_{L^1}e^{-na/2},\qquad
\sup_{\chi}|\widehat G_n(\chi)|\le\lVert\vartheta\rVert_{L^1}e^{-na/2}.
\end{equation}
\item[(c)] $I_A(F_n,\cdot)\to0$ and $I_A(G_n,\cdot)\to0$ in $\DA'$, that is,
uniformly on every bounded subset of $\DA$.
\item[(d)] Consequently $J(F_n)\to(0,e_+)=:e_+^{\mathrm{dist}}$ and
$J(G_n)\to(0,e_-)=:e_-^{\mathrm{dist}}$ in $\DA'\oplus\BT$, and both limits
lie in $\NAd$.
\item[(e)] $\DA$ acts on them by the two Tate characters:
$a\cdot e_{\pm}^{\mathrm{dist}}=\lambda_{\pm}(a)e_{\pm}^{\mathrm{dist}}$ for
every $a\in\DA$, and $\dagger$ exchanges them.
\end{enumerate}
\end{theorem}

\begin{proof}
Write $t=\log|w|_S$.  By Lemma~\ref{thm:split} and the normalization
\eqref{eq:archmeasure}, $\dstar w=dt\otimes\mu_{\mathrm{cpt}}$ with
$\mu_{\mathrm{cpt}}$ a probability measure on the compact part; so for a
function of $|w|_S$ alone,
$\int_{C_S}\Phi(|w|_S)\dstar w=\int_{\R}\Phi(e^{t})dt$.

(a)  $\lambda_{\mp}(F_n)=\int_{\R}n^{-1}e^{-t/2}\vartheta(t/n)e^{\mp t/2}dt$.  For
the upper sign, substituting $t=ns$ gives
$\lambda_-(F_n)=\int\vartheta(s)e^{-ns}ds$, which is positive since $\vartheta\ge0$ and
$\vartheta\ne0$, and bounded by $e^{-na}\lVert\vartheta\rVert_{L^1}$ since $s\ge a$ on
$\supp\vartheta$.  For the lower sign,
$\lambda_+(F_n)=\int n^{-1}\vartheta(t/n)dt=\int\vartheta(s)ds=1$.  The statements for
$G_n$ follow after $t\mapsto-t$.

(b)  $F_n$ is constant on the compact coordinates and even in the sign, so by
orthogonality $\widehat F_n$ vanishes on every component but $\mathfrak c_0$,
where
$\widehat F_n(\tau)=\int n^{-1}e^{-t/2}\vartheta(t/n)e^{-i\tau t}dt=
\int\vartheta(s)e^{-ns(\frac12+i\tau)}ds$; hence $|\widehat F_n(\tau)|\le
\int\vartheta(s)e^{-ns/2}ds\le e^{-na/2}\lVert\vartheta\rVert_{L^1}$, uniformly in
$\tau$.  Likewise $\widehat G_n(\tau)=\int\vartheta(s)e^{-ns(\frac12-i\tau)}ds$.

(c)  By Lemma~\ref{lem:Nseminorm} and \eqref{eq:supFhat}, for every bounded
subset $B\subseteq\DA$ one has
\[
\sup_{G\in B}|I_A(F_n,G)|\le\bigl(\sup_{\chi}|\widehat F_n|\bigr)
\sup_{G\in B}N(G)\le\lVert\vartheta\rVert_{L^1}e^{-na/2}C_B\longrightarrow0 .
\]

(d)  $J(F_n)=(I_A(F_n,\cdot),M(F_n))$; the first component tends to $0$ in
$\DA'$ by (c) and the second to $e_+$ by (a).  Each $J(F_n)$ lies in $J(N_A)$,
so the limit lies in its closure.

(e)  By Theorem~\ref{thm:completion} the action is
$a\cdot(\ell,m)=(a\cdot\ell,M(a)m)$.  On $(0,e_+)$ the first component stays
$0$ and the second is $M(a)e_+=\lambda_+(a)e_+$, the product in $\BT$ being
componentwise.  The involution acts on $\BT$ by $\sharp$, which exchanges $e_-$
and $e_+$.
\end{proof}

\begin{theorem}[Canonical splitting of the completion]\label{thm:splitting}
Let $\bar M:\DA'\oplus\BT\to\BT$ be the second-component projection, put
$\iota(m):=(0,m)$ and
$N_A^{\mathrm{dist,prim}}:=\NAd\cap(\DA'\oplus0)=\ker(\bar M|_{\NAd})$.  Then $\iota$ is
a continuous injective $\DA$-module map landing in $\NAd$, $\bar M\circ
\iota=\id_{\BT}$, and
\begin{equation}\label{eq:split}
\NAd=N_A^{\mathrm{dist,prim}}\ \oplus\ \iota(\BT)
\end{equation}
as topological $\DA$-modules; the decomposition involves no choices.  The involution
preserves the decomposition, acting by $\sharp$ on the second summand.
\end{theorem}

\begin{proof}
$\iota$ lands in $\NAd$ because $\iota(e_{\pm})=e_{\pm}^{\mathrm{dist}}\in\NAd$ by
Theorem~\ref{thm:rulings}(d) and $\NAd$ is a linear subspace; it is clearly
continuous and injective with $\bar M\iota=\id$.  It is a module map:
$a\cdot \iota(m)=(0,M(a)m)=\iota(M(a)m)$ by Theorem~\ref{thm:rulings}(e).  Given
$x=(\ell,m)\in\NAd$ we have $\iota(m)\in\NAd$, hence $x-\iota(m)=(\ell,0)\in\NAd$ and
lies in $N_A^{\mathrm{dist,prim}}$; the two summands meet in $0$ since $\bar M$
kills the first and is injective on the second.  Both are closed, the first as
the kernel of a continuous map and the second being finite dimensional, and
the two projections $x\mapsto \iota(\bar Mx)$ and $x\mapsto x-\iota(\bar Mx)$ are
continuous, so the sum is topological.  Both summands are submodules by the
displayed action, and $\dagger$ acts by $(\ell,m)^{\dagger}=(\ell^{\dagger},
m^{\sharp})$, preserving each.
\end{proof}

\begin{theorem}\label{thm:rulnull}
Put
\begin{multline}\label{eq:Rboundary}
R_{\partial}:=\bigl\{x\in\NAd:\ \text{the first component of }
I_{\mathrm{vec}}(x,G)\text{ vanishes for every }G\in\DA\bigr\}
\\
=\ \bigl\{(\ell,m)\in\NAd:\ \ell=0\bigr\}.
\end{multline}
Then:
\begin{enumerate}
\item[(a)] $R_{\partial}=\iota(\BT)$, a two-dimensional plane on which the Tate
component of $I_{\mathrm{vec}}$ is the nondegenerate hyperbolic form
$\beta_T(m,MG)$, with matrix
$\left(\begin{smallmatrix}0&1\\1&0\end{smallmatrix}\right)$ in the basis
$e_{\mp}^{\mathrm{dist}}$;
\item[(b)] $\lim_{n,m\to\infty}I_A(F_n,G_m)=0$, uniformly in one index while
the other tends to infinity, whereas
$\beta_T(MF_n,MG_m)\to\beta_T(e_+,e_-)=1$;
\item[(c)] $R_{\partial}\cap J(N_A)=0$.
\end{enumerate}
\end{theorem}

\begin{proof}
(a)  The two descriptions in \eqref{eq:Rboundary} agree because
$I_{\mathrm{vec}}((\ell,m),G)=(\ell(G),\beta_T(m,MG))$, and $\ell(G)=0$ for all
$G$ if and only if $\ell=0$; the identification with $\iota(\BT)$ is
Theorem~\ref{thm:splitting}, and the form is Theorem~\ref{thm:tate}.

(b)  The quantity $N(G_m)$ is bounded in $m$: by the computation in
Theorem~\ref{thm:rulings}(b), $\widehat G_m$ lives on $\mathfrak c_0$ and
$\widehat G_m(\tau)=\int\vartheta(s)e^{-ms(\frac12-i\tau)}ds$, so two integrations
by parts in $s$ give $|\widehat G_m(\tau)|\le C_{\vartheta}\min(1,(m|\tau|)^{-2})$
with $C_{\vartheta}$ depending only on $\vartheta$; hence
\[
N(G_m)\le C_{\vartheta}\int_{\R}\min\bigl(1,(m|\tau|)^{-2}\bigr)
\bigl(A_{\mathfrak c_0}+|c|+B_S\log(2+|\tau|)\bigr)\frac{d\tau}{2\pi}\le C_{\vartheta}C_S ,
\]
uniformly in $m\ge1$; the integral is in fact $O(m^{-1}\log m)$.  Therefore
$|I_A(F_n,G_m)|\le(\sup|\widehat F_n|)N(G_m)\le\lVert\vartheta\rVert_{L^1}
C_{\vartheta}C_Se^{-na/2}$, which tends to $0$ as $n\to\infty$ uniformly in $m$; by
symmetry the same holds with the roles exchanged.  Meanwhile
$\beta_T(MF_n,MG_m)=\lambda_-(F_n)\overline{\lambda_+(G_m)}+
\lambda_+(F_n)\overline{\lambda_-(G_m)}\to0\cdot0+1\cdot1=1$ by
Theorem~\ref{thm:rulings}(a).

(c)  If $J(F)\in R_{\partial}$ then $I_A(F,G)=0$ for every $G$, i.e.\
$F\in\Nrad=\{0\}$ by Theorem~\ref{thm:radical}, so $J(F)=0$.
\end{proof}

\begin{definition}\label{def:admissible}
A net $(F_{\alpha})$ in $\DA$ is an \emph{admissible approximation of $e_+$} if
\begin{equation}\label{eq:admissible}
\text{(i) } M(F_{\alpha})\to e_+\ \text{in }\BT,\qquad
\text{(ii) } \lVert F_{\alpha}\rVert_{L^1(\dstar w)}\to0 ;
\end{equation}
similarly for $e_-$.  A net $(G_{\beta})$ in $\DA$ is \emph{$N$-bounded} if
$\sup_{\beta}N(G_{\beta})<\infty$, with $N$ as in \eqref{eq:Nsemi}.
\end{definition}

Condition (ii) is what makes the approximation \emph{spread}: by
the inequality between the arithmetic and geometric means, for $F\ge0$ with
$M(F)$ near $e_+$ one has only $\lVert F\rVert_{L^1}\le\tfrac12(\lambda_++
\lambda_-)\approx\tfrac12$; so (ii) says that the mass of $F_{\alpha}$ escapes
to $|w|_S\to\infty$ at exactly the rate that makes the Tate weight
$|w|_S^{1/2}$ carry all of it.

\begin{theorem}\label{thm:rulnullclass}
\begin{enumerate}
\item[(a)] If $(F_{\alpha})$ is an admissible approximation of $e_{\pm}$ then
$J(F_{\alpha})\to \iota(e_{\pm})$ in $\DA'\oplus\BT$; that is, $I_A(F_{\alpha},
\cdot)\to0$ uniformly on bounded subsets of $\DA$.
\item[(b)] If moreover $(G_{\beta})$ is any $N$-bounded net, then
$|I_A(F_{\alpha},G_{\beta})|\le\lVert F_{\alpha}\rVert_{L^1}\sup_{\beta}
N(G_{\beta})\to0$ uniformly in $\beta$; in particular for any admissible
$N$-bounded approximation $(G_{\beta})$ of $e_{\mp}$,
\begin{equation}\label{eq:doublelimit}
\lim_{\alpha}\lim_{\beta}I_A(F_{\alpha},G_{\beta})
=\lim_{\beta}\lim_{\alpha}I_A(F_{\alpha},G_{\beta})=0,\qquad\text{while}\quad
\beta_T(MF_{\alpha},MG_{\beta})\to\beta_T(e_{\pm},e_{\mp})=1 .
\end{equation}
\item[(c)] The family \eqref{eq:spread} is admissible and $N$-bounded:
$(F_n)$ for $e_+$ and $(G_m)$ for $e_-$.
\end{enumerate}
\end{theorem}

\begin{proof}
(a)  $\sup_{\chi}|\widehat F_{\alpha}(\chi)|\le\lVert F_{\alpha}\rVert_{L^1}$
by the Riemann--Lebesgue estimate, so for $G$ in a bounded set $B$,
$|I_A(F_{\alpha},G)|\le\lVert F_{\alpha}\rVert_{L^1}\sup_BN\to0$ by
Lemma~\ref{lem:Nseminorm}; the second component is \eqref{eq:admissible}(i).
(b) is the same inequality with $\sup_{\beta}N(G_{\beta})$ in place of
$\sup_BN$, and $\beta_T$ is continuous.  (c)
$\lVert F_n\rVert_{L^1}=\int\vartheta(s)e^{-ns/2}ds\le\lVert\vartheta\rVert_{L^1}
e^{-na/2}\to0$, and $M(F_n)\to e_+$ by Theorem~\ref{thm:rulings}(a); and
$N(F_n)=O(n^{-1}\log n)$ by the estimate in the proof of
Theorem~\ref{thm:rulnull}(b), the two families being mirror images under
$t\mapsto-t$.
\end{proof}

\begin{remark}\label{rem:noextension}
Theorem~\ref{thm:rulnullclass} does not extend $I_A$ continuously to
$\NAd\times\NAd$, and we do not construct such an extension.  The completion
$\NAd$ is a module, not an algebra: the product of two arbitrary elements is
not defined.  The theorem is a statement about limits along the class of
approximations of Definition~\ref{def:admissible}.
\end{remark}

\begin{remark}\label{rem:nocontradiction}
There is no conflict with Theorem~\ref{thm:nosplit}(a).  That theorem forbids
an eigenvector $\Phi_{\mp}\in\DA$ with $a\star\Phi_{\mp}=\lambda_{\mp}(a)
\Phi_{\mp}$; its proof compares a value $\widehat a(\chi_0)$ \emph{on} the
unitary axis with $\widehat a(\mp i/2)$ \emph{off} it.  The eigenvectors
$e_{\pm}^{\mathrm{dist}}$ of Theorem~\ref{thm:rulings}(e) have zero
distribution component: there is no $\chi_0$ at which to evaluate, and the
eigenvalue equation holds in the $\BT$-component alone, where the action is by
the characters $\lambda_{\mp}$ by the very definition \eqref{eq:J} of the
module structure.  Theorem~\ref{thm:rulnull}(c) says precisely that these
eigenvectors lie outside the dense image $J(N_A)$, which is what
Theorem~\ref{thm:nosplit}(a) is a statement about.  The two theorems together
say: the ruling plane is a quotient of the dense test sector and a subobject of
its completion, and the two facts are compatible because the completion is
strictly larger.  Formula \eqref{eq:split} then reads
$\NAd=(\text{primitive})\oplus(\text{ruling plane})$, with the scalar form
living on the first summand and the hyperbolic form on the second --- the
shape of a N\'eron--Severi group as an orthogonal sum of a primitive part and a
hyperbolic plane, with the caveat that no sign is asserted on the first
summand; see Theorem~\ref{thm:signature}.
\end{remark}

\section{Presented complexes and the determinant}\label{sec:presented}

The generators $E_a$ are not in the domain of $T$: a point mass is not a
compactly supported smooth section, and the diagonal is not trace class
(Remark~\ref{rem:diagonal}).  The natural domain is a category of complexes
that carry their presentation, in the way a cycle carries its
multiplicities.

\begin{definition}\label{def:Corrtr}
$\Corrtr$ is the category whose objects are pairs $(K,\mathfrak f)$ where
$\mathfrak f=(F_n)_{n\in\Z}$ is a finitely supported family of elements of
$\DA$ (the \emph{presentation}) and $K$ is the complex whose $n$-th term is the
smeared object $E_{F_n}$ and whose differentials are elements of $\DA$ acting
by convolution, with $d^{2}=0$; morphisms are morphisms of such presented
complexes.  Set
\begin{equation}\label{eq:s}
s(K,\mathfrak f):=\sum_n(-1)^{n}F_n\in\DA .
\end{equation}
$\Corrtr$ carries $\oplus$ (concatenation of presentations), shift
($\mathfrak f[1]_n=F_{n+1}$), cone
($\Cone(f)_n=F^{L}_n\oplus F^{K}_{n+1}$) and $\star$ (convolution of
presentations, totalized).
\end{definition}

\begin{lemma}\label{lem:slaws}
$s(K\oplus L)=s(K)+s(L)$, $s(K[1])=-s(K)$, $s(\Cone f)=s(L)-s(K)$ for
$f:K\to L$, and $s(K\star L)=s(K)\star s(L)$.
\end{lemma}

\begin{proof}
The first three are immediate from the term descriptions.  For the last, the
$n$-th term of $K\star L$ is $\bigoplus_{i+j=n}E_{F_i\star G_j}$ by
Corollary~\ref{cor:trivial}, so
$s(K\star L)=\sum_n(-1)^{n}\sum_{i+j=n}F_i\star G_j=
\bigl(\sum_i(-1)^{i}F_i\bigr)\star\bigl(\sum_j(-1)^{j}G_j\bigr)$, the
rearrangement being legitimate because the families are finite.
\end{proof}

\begin{proposition}\label{prop:faithful}
The smearing map $\DA\to\{\text{sections}\}$, $F\mapsto E_F$ of
\eqref{eq:smear}, is injective, and so is its composition with the regular
representation $\lambda$ of $L^{1}(C_S)$ on $L^{2}(C_S)$, $\lambda(F)f=F\star
f$.
\end{proposition}

\begin{proof}
Two proofs.  First, by Corollary~\ref{cor:trivial} the section $\sigma$
trivializes the bundle strictly, so the coefficient of $E_F$ at the label $u$
is exactly $F(u)$, with no coefficient at any nontrivial gauge label; hence
$E_F=E_G$ forces $F=G$ pointwise.  Second, $\DA\subseteq L^{1}(C_S)$ and
$\lambda$ is faithful on $L^{1}(C_S)$: if $\lambda(F)=0$ then $F\star G=0$ for
every $G\in L^{1}\cap L^{2}$, so $\widehat F\widehat G=0$ almost everywhere for
every such $G$; taking $\widehat G$ nonvanishing on a prescribed compact set
gives $\widehat F=0$ almost everywhere, hence $F=0$ by Fourier inversion.
Under the trivialization the action of $E_F$ on the carrier is $\lambda(F)$.
This second argument uses nothing about $\DA$ beyond $\DA\subseteq L^{1}(C_S)$,
and is extended to $\Atr$ in Theorem~\ref{thm:Atrcat}(b).
\end{proof}

\begin{remark}\label{rem:cycles}
Proposition~\ref{prop:faithful} says that the presentation is faithfully
recorded by the smeared objects, so that $\Corrtr$ is a category of
\emph{based} objects and $s$ is the analogue of a cycle class with its
multiplicities.  One should not ask, in addition, that the trace descend to
isomorphism classes of the underlying complexes.  No trace of correspondences
has that property, in this theory or in the classical one: on a surface, for
an algebraic correspondence $\Gamma$ and $n\ge2$, the cycles $n\Gamma$ and
$\Gamma$ have the same underlying closed subscheme, hence are isomorphic as
geometric objects, while $\Tr(n\Gamma)=n\Tr(\Gamma)$.  The classical trace is a
function of a cycle --- an element of a free abelian group on subvarieties ---
and not of the underlying object.
\end{remark}

\begin{definition}\label{def:target}
Let $\Picb(\C)$ be the Picard groupoid of based one-dimensional complex lines,
$(\C,s)$ with $s\in\C^{\times}$, and let
\begin{equation}\label{eq:target}
T_A:=\Picb(\BT),\qquad \BT=\C_-\oplus\C_+,
\end{equation}
the Picard groupoid of based invertible $\BT$-modules, equivalently of pairs of
based complex lines, with $\sharp$ exchanging the two factors.  This is the
target of a Knudsen--Mumford determinant for finite-cell $\BT$-modules
\cite{KnudsenMumford}, and $\lambda_{RR}$ below is a based Hermitian
biextension in the sense of Deligne's pairing \cite{Deligne}.  The scalar
target embeds as the first factor, $\Picb(\C)\hookrightarrow T_A$,
$L\mapsto(L,\mathrm{unit})$.  A based line $(\C,s)$ carries the metric
\begin{equation}\label{eq:metric}
\lVert1\rVert:=|s|^{-1/2},\qquad\text{equivalently}\qquad
\lVert1\rVert^{2}=|s|^{-1} .
\end{equation}
The trace with values in the target is
\begin{equation}\label{eq:TrA}
\Tr_A:\Corrtr\to T_A,\qquad
\Tr_A(K):=\bigl(\bigl(\C,e^{T(s(K))}\bigr),\ \mathrm{unit}\bigr),
\end{equation}
and we write $T(K):=T(s(K))\in\C$ for its scalar shadow, and
$\Det^{\mathrm{rel}}_S(F):=e^{T(F)}$ for $F\in\DA$.
\end{definition}

With the convention \eqref{eq:metric}, the line
$\Det^{\mathrm{rel}}_S(F\star G^{\dagger})=(\C,e^{I_A(F,G)})$ has
$\lVert1\rVert^{2}=e^{-\Rel I_A(F,G)}$.  The enlargement of the target from
$\Picb(\C)$ to $\Picb(\BT)$ is needed because of Theorem~\ref{thm:scalartrivial}
below, which shows that the scalar factor $L_S$ alone is trivial on the ruling
plane.

\begin{theorem}[Additivity]\label{thm:additivity}
$\Tr_A(K\oplus L)\cong\Tr_A(K)\otimes\Tr_A(L)$,
$\Tr_A(K[1])\cong\Tr_A(K)^{-1}$, and for a triangle $K\to L\to\Cone\to[1]$ in
$\Corrtr$, $\Tr_A(\Cone)\cong\Tr_A(L)\otimes\Tr_A(K)^{-1}$.  Equivalently
$T(K\oplus L)=T(K)+T(L)$, $T(K[1])=-T(K)$ and $T(\Cone f)=T(L)-T(K)$.
\end{theorem}

\begin{proof}
The three identities for $s$ are Lemma~\ref{lem:slaws}, and $T$ is linear on
$\DA$; multiplicativity of the based line is the identity
$e^{x+y}=e^{x}e^{y}$ under the tensor product of based lines, which multiplies
the distinguished generators.
\end{proof}

\begin{theorem}[Cyclicity and reality]\label{thm:cyclic}
$\Tr_A(K\star L)\cong\Tr_A(L\star K)$ for all $K,L\in\Corrtr$, and
$\Tr_A(K^{\dagger})\cong\overline{\Tr_A(K)}$, i.e.\
$T(K^{\dagger})=\overline{T(K)}$.
\end{theorem}

\begin{proof}
By Lemma~\ref{lem:slaws}, $s(K\star L)=s(K)\star s(L)$ and $s(L\star
K)=s(L)\star s(K)$, and $\DA$ is commutative, so the two agree; apply $T$.
Reality is \eqref{eq:daggerreal} applied to $s(K^{\dagger})=s(K)^{\dagger}$,
the shift signs being preserved since $\dagger$ is degree preserving on
$\Corrtr$.
\end{proof}

Cyclicity is automatic here, since $C_S$ is abelian.

\begin{proposition}[Invariance]\label{prop:invariance}
$\Tr_A$ is invariant under (i) isomorphisms of presented complexes inducing the
identity on presentations, (ii) addition of an elementary contractible presented
complex, i.e.\ one of the form $0\to E_F\xrightarrow{\ \cong\ }E_F\to0$ in
adjacent degrees, and (iii) the operations of Theorem~\ref{thm:additivity}.
\end{proposition}

\begin{proof}
(i) $\Tr_A$ depends only on $\mathfrak f$.  (ii) Such a complex has
$s=F-F=0$, so $T=0$, and $s$ is additive on $\oplus$.  (iii) is
Theorem~\ref{thm:additivity}.
\end{proof}

We do not expect $\Tr_A$ to descend to quasi-isomorphism classes of the
underlying complexes; see Remark~\ref{rem:cycles}.

\section{The trace Banach algebra}\label{sec:banach}

The determinant of the theory is $\exp T$.  In this section we construct a
complete algebra norm for which it is continuous, and we show that the
associated determinant germ is holomorphic and multiplicative, all places
being treated at once.

\begin{definition}\label{def:trnorm}
For $F\in\DA$ let $N(F)$ be the seminorm \eqref{eq:Nsemi} of
Lemma~\ref{lem:Nseminorm}, and put
\begin{equation}\label{eq:trnorm}
\Lambda(F):=\int_{C_S}|F(w)|\bigl(|w|_S^{1/2}+|w|_S^{-1/2}\bigr)\dstar w,
\end{equation}
\[
\lVert F\rVert_{\mathrm{tr}}:=\Lambda(F)+N(F).
\]
Here $N(F)<\infty$ by Lemma~\ref{lem:Nseminorm} and $\Lambda(F)<\infty$ because
$F$ has compact support.  By the inequality
$|w|_S^{1/2}+|w|_S^{-1/2}\ge2$,
\begin{equation}\label{eq:trnormineq}
\lVert F\rVert_{L^{1}}\le\tfrac12\Lambda(F)\le\tfrac12\lVert
F\rVert_{\mathrm{tr}},\qquad
|\lambda_{\mp}(F)|\le\Lambda(F)\le\lVert F\rVert_{\mathrm{tr}} .
\end{equation}
Let $\Atr$ be the completion of $\DA$ for $\lVert\cdot\rVert_{\mathrm{tr}}$.
\end{definition}

\begin{theorem}\label{thm:banach}
\begin{enumerate}
\item[(a)] For all $F,G\in\DA$,
$\lVert F\star G\rVert_{\mathrm{tr}}\le\lVert F\rVert_{\mathrm{tr}}
\lVert G\rVert_{\mathrm{tr}}$ and $\lVert F^{\dagger}\rVert_{\mathrm{tr}}
=\lVert F\rVert_{\mathrm{tr}}$.  Hence $\Atr$ is a commutative Banach
$\dagger$-algebra.  The natural map $\Atr\to L^{1}(C_S)$ is injective, so
$\Atr$ is realized as a subalgebra of $L^{1}(C_S)$ contained in
$\{H\in L^{1}:\Lambda(H)+N(H)<\infty\}$.
\item[(b)] $\lambda_-$ and $\lambda_+$ extend uniquely to continuous linear
functionals on $\Atr$ of norm $\le1$; they remain multiplicative and satisfy
$\lambda_{\mp}(H^{\dagger})=\overline{\lambda_{\pm}(H)}$, so that
$M=(\lambda_-,\lambda_+):\Atr\to\BT$ is a continuous $\dagger$-algebra
homomorphism extending that of Theorem~\ref{thm:tate}.
\item[(c)] $T$ extends uniquely to a continuous linear functional on $\Atr$,
still given by $T(H)=\int\widehat H(\Th+c)\,d\chi$, with $|T(H)|\le N(H)\le
\lVert H\rVert_{\mathrm{tr}}$; and $\Det^{\mathrm{rel}}_S=\exp T$ extends to a
continuous character for the additive structure,
$\Det^{\mathrm{rel}}_S(H+K)=\Det^{\mathrm{rel}}_S(H)\Det^{\mathrm{rel}}_S(K)$.
\item[(d)] For every $H\in\Atr$ and every $r\ge1$,
\begin{equation}\label{eq:powerest}
\bigl|T(H^{\star r})\bigr|\le\lVert H\rVert_{L^{1}}^{\,r-1}N(H)
\le\lVert H\rVert_{L^{1}}^{\,r-1}\lVert H\rVert_{\mathrm{tr}} .
\end{equation}
\item[(e)] For $|z|<1/\lVert F\rVert_{L^{1}}$ the series
\begin{equation}\label{eq:germ}
\log\Delta_F(z):=\sum_{r\ge1}(-1)^{r+1}\frac{z^{r}T(F^{\star r})}{r}
\end{equation}
converges absolutely, is holomorphic in $z$ on that disc, continuous in $F$ for
$\lVert\cdot\rVert_{\mathrm{tr}}$, and equals
\begin{equation}\label{eq:germint}
\log\Delta_F(z)=\int_{\Ch}\Log\bigl(1+z\widehat F(\chi)\bigr)
\bigl(\Th(\chi)+c\bigr)d\chi ,
\end{equation}
$\Log$ the principal branch.
\item[(f)] For $F,G\in\Atr$ and $|z|<\min\bigl(1/\lVert F\rVert_{L^{1}},
1/\lVert G\rVert_{L^{1}}\bigr)$, setting $H:=F+G+zF\star G$,
\begin{equation}\label{eq:multiplicative}
\Delta_F(z)\cdot\Delta_G(z)=\Delta_H(z),\qquad
\Delta_{F^{\dagger}}(\bar z)=\overline{\Delta_F(z)} .
\end{equation}
\item[(g)] For every $p\in S_f$ and every $F\in\Atr$, writing
$a^{(r)}_{\kappa}$ for the unit-block Fourier coefficients of $F^{\star r}$ and
$N^{u}_p(F):=\int|\widehat F|\,c_p(\kappa_p)\log p\,d\chi<\infty$, one has
\[
\Bigl|\sum_{\kappa}c_p(\kappa)a^{(r)}_{\kappa}\Bigr|
\le\lVert F\rVert_{L^{1}}^{\,r-1}(\log p)^{-1}N^{u}_p(F)
\]
uniformly in $r$.  Consequently the unit-packet component of the determinant
germ is holomorphic on $|z|<1/\lVert F\rVert_{L^{1}}$, and it is the
$\widehat W^{u}_p$-component of \eqref{eq:germint}, where
$\widehat W^{u}_p(\chi)=-c_p(\kappa_p)\log p$.
\end{enumerate}
\end{theorem}

\begin{proof}
(a)  For $s=\pm\tfrac12$ write $\lVert F\rVert_{1,s}=\int|F(w)||w|_S^{s}\dstar
w$, so $\Lambda=\lVert\cdot\rVert_{1,1/2}+\lVert\cdot\rVert_{1,-1/2}$.  Since
$w\mapsto|w|_S^{s}$ is a multiplicative character,
$|F\star G(w)||w|^{s}\le\int|F(u)||u|^{s}|G(u^{-1}w)||u^{-1}w|^{s}\dstar u$,
and Fubini gives the weighted Young inequality
$\lVert F\star G\rVert_{1,s}\le\lVert F\rVert_{1,s}\lVert G\rVert_{1,s}$.
Adding the two values of $s$ and using nonnegativity of all four quantities,
$\Lambda(F\star G)\le\Lambda(F)\Lambda(G)$.  Since $(F\star G)^{\wedge}=
\widehat F\widehat G$ and $\lVert\widehat G\rVert_{\infty}\le\lVert
G\rVert_{L^{1}}\le\tfrac12\Lambda(G)$ by \eqref{eq:trnormineq},
\[
N(F\star G)=\int|\widehat F||\widehat G|(|\Th|+|c|)d\chi
\le\tfrac12\Lambda(G)N(F)\le\Lambda(G)N(F).
\]
Adding, $\lVert F\star G\rVert_{\mathrm{tr}}\le\Lambda(G)
(\Lambda(F)+N(F))=\Lambda(G)\lVert F\rVert_{\mathrm{tr}}\le
\lVert F\rVert_{\mathrm{tr}}\lVert G\rVert_{\mathrm{tr}}$; the constant is $1$.
For the involution, $\dstar w$ is inversion invariant and
$|w^{-1}|_S^{s}=|w|_S^{-s}$, so $w\mapsto w^{-1}$ exchanges the two weights
and $\Lambda(F^{\dagger})=\Lambda(F)$; and $(F^{\dagger})^{\wedge}=
\overline{\widehat F}$ (computed in the proof of Theorem~\ref{thm:radical}), so
$N(F^{\dagger})=N(F)$.  For the realization: a
$\lVert\cdot\rVert_{\mathrm{tr}}$-Cauchy sequence $(F_j)$ is $L^{1}$-Cauchy,
hence converges in $L^{1}$ to some $H$, and $\widehat F_j\to\widehat H$
uniformly.  It is also Cauchy in the two complete spaces
$L^{1}\bigl((|w|^{1/2}+|w|^{-1/2})\dstar w\bigr)$ and
$L^{1}\bigl((|\Th|+|c|)d\chi\bigr)$, so it converges there to limits $H'$ and
$\phi$.  A subsequence converges almost everywhere in each, so $H'=H$ and,
since $\widehat F_j\to\widehat H$ uniformly, $\phi=\widehat H$ almost
everywhere.  Hence $\lVert F_j\rVert_{\mathrm{tr}}\to\Lambda(H)+N(H)$ exactly,
and more generally $\lVert F_j-F_k\rVert_{\mathrm{tr}}$ and
$\Lambda(F_j-H)+N(F_j-H)$ tend to $0$.  In particular the map from $\Atr$ to
$L^{1}$ sending the class of $(F_j)$ to $H$ is well defined, isometric for the
norm $\Lambda+N$ extended to its image, and injective: if $H=0$ then
$\lVert F_j\rVert_{\mathrm{tr}}\to0$.

(b)  $|\lambda_{\mp}(F)|=\bigl|\int F|w|^{\mp1/2}\dstar w\bigr|\le
\lVert F\rVert_{1,\mp1/2}\le\Lambda(F)\le\lVert F\rVert_{\mathrm{tr}}$, whence
the unique continuous extension of norm $\le1$.  Multiplicativity, the
involution rule and surjectivity onto $\BT$ pass to the completion, both sides
of each identity being $\lVert\cdot\rVert_{\mathrm{tr}}$-continuous and
agreeing on the dense subalgebra $\DA$.

(c)  The bound is \eqref{eq:plancherel} together with
\eqref{eq:trnorm}; additivity of $T$ gives the character property, and $\exp$
is continuous.

(d)  $(H^{\star r})^{\wedge}=\widehat H^{r}$, so
$|T(H^{\star r})|\le\int|\widehat H|^{r}(|\Th|+|c|)d\chi\le
\lVert\widehat H\rVert_{\infty}^{r-1}N(H)\le\lVert H\rVert_{L^{1}}^{r-1}N(H)$.

(e)  By (d),
$\sum_r|z|^{r}|T(F^{\star r})|/r\le N(F)\sum_r(|z|\lVert F\rVert_{L^1})^{r}/
(r\lVert F\rVert_{L^1})<\infty$ for $|z|\lVert F\rVert_{L^1}<1$, giving
absolute convergence and holomorphy of a locally uniformly convergent power
series.  Continuity in $F$: each $T(F^{\star r})$ is
$\lVert\cdot\rVert_{\mathrm{tr}}$-continuous by (a) and (c), and the tail bound
is uniform on $\lVert\cdot\rVert_{\mathrm{tr}}$-bounded sets.  For
\eqref{eq:germint}, insert $T(F^{\star r})=\int\widehat F^{r}(\Th+c)d\chi$ and
interchange sum and integral; this is legitimate by dominated convergence,
since $\sum_r|z\widehat F|^{r}/r\le|z||\widehat F|/(1-|z|\lVert
F\rVert_{L^1})$ and $|\widehat F|(|\Th|+|c|)$ is integrable, and
$\sum_r(-1)^{r+1}(z\widehat F)^{r}/r=\Log(1+z\widehat F)$ pointwise.

(f)  For each $\chi$, both $1+z\widehat F(\chi)$ and $1+z\widehat G(\chi)$ lie
in the open disc of radius $1$ centred at $1$, hence in the right half plane
with argument in $(-\pi/2,\pi/2)$; so the sum of their principal logarithms has
imaginary part in $(-\pi,\pi)$ and therefore equals the principal logarithm of
the product, $\Log\bigl((1+z\widehat F)(1+z\widehat G)\bigr)=\Log(1+z\widehat
H)$ with $\widehat H=\widehat F+\widehat G+z\widehat F\widehat G$.  Integrating
against $(\Th+c)d\chi$ and exponentiating gives the first identity of
\eqref{eq:multiplicative}.  For the second, $\Th+c$ is real by
Theorem~\ref{thm:symbol}(c) and $(F^{\dagger})^{\wedge}=\overline{\widehat F}$,
so $\int\Log(1+\bar z\overline{\widehat F})(\Th+c)d\chi=
\overline{\int\Log(1+z\widehat F)(\Th+c)d\chi}$.

(g)  By \eqref{eq:Wphat} the unit-packet part of the symbol at $p$ is the
bounded function $\widehat W^{u}_p(\chi)=-c_p(\kappa_p)\log p$, and by
\eqref{eq:conductor} the unit packet of $W_p(H)$ equals
$\int\widehat H\,\widehat W^{u}_p\,d\chi=-(\log p)\sum_{\kappa}
c_p(\kappa)a_{\kappa}(H)$.  Applying this to $H=F^{\star r}$ and estimating as
in (d),
\[
\Bigl|(\log p)\sum_{\kappa}c_p(\kappa)a^{(r)}_{\kappa}\Bigr|
=\Bigl|\int\widehat F^{r}\widehat W^{u}_p\,d\chi\Bigr|
\le\lVert F\rVert_{L^{1}}^{\,r-1}N^{u}_p(F),
\]
and $N^{u}_p(F)<\infty$ by the argument of Lemma~\ref{lem:Nseminorm}, the
integrand being $|\widehat F|$, integrable over the finitely many components
met by $\supp\widehat F$, times a function bounded on them.  Summing the
geometric series gives the holomorphy, and comparison of \eqref{eq:germint}
with the decomposition $\Th=\Wh_{\infty}+\sum_p\Wh_p$ identifies the unit
packet component.
\end{proof}

\begin{theorem}[The trace-class category over $\Atr$]\label{thm:Atrcat}
For $H\in\Atr\subseteq L^{1}(C_S)$ put $E_H:=\lambda(H)$, where $\lambda$ is
the regular representation of $L^{1}(C_S)$ on $L^{2}(C_S)$, $\lambda(H)f=H\star
f$.  Then:
\begin{enumerate}
\item[(a)] \emph{(Boundedness and compatibility)}
$\lVert\lambda(H)\rVert_{\mathrm{op}}\le\lVert H\rVert_{L^{1}}\le
\tfrac12\lVert H\rVert_{\mathrm{tr}}$; $\lambda(H\star K)=\lambda(H)\lambda(K)$;
and $\lambda(H^{\dagger})=\lambda(H)^{*}$.  Hence $H\mapsto E_H$ is a bounded,
multiplicative, $\dagger$-compatible linear map $\Atr\to
\mathcal B(L^{2}(C_S))$.
\item[(b)] \emph{(Faithfulness)}  $H\mapsto E_H$ is injective on $\Atr$.
\item[(c)] \emph{(Extension)}  On $\DA$ the operator $E_H$ is the action of the
smeared object $E_F$ of Corollary~\ref{cor:trivial} on the carrier, so
$F\mapsto E_F$ extends $\lVert\cdot\rVert_{\mathrm{tr}}$-continuously and
injectively from $\DA$ to $\Atr$.
\item[(d)] \emph{(The category)}  Let $\Corrtr(\Atr)$ be defined as in
Definition~\ref{def:Corrtr} with $\DA$ replaced by $\Atr$ throughout:
presentations $\mathfrak f=(H_n)$ with $H_n\in\Atr$, terms $E_{H_n}$, and
differentials in $\Atr$ acting by convolution.  Then $s$,
Lemma~\ref{lem:slaws}, Theorems~\ref{thm:additivity} and \ref{thm:cyclic},
Proposition~\ref{prop:invariance} and $\Tr_A=\exp\circ\,T\circ s$ hold on
$\Corrtr(\Atr)$ verbatim; and $\Corrtr(\Atr)$ is closed under
$\lVert\cdot\rVert_{\mathrm{tr}}$-limits of presentations with a common finite
support: if $\mathfrak f^{(k)}=(H^{(k)}_n)$ are presentations of objects of
$\Corrtr(\Atr)$, all supported in one finite set of degrees, with
$H^{(k)}_n\to H_n$ and $d^{(k)}_n\to d_n$ in
$\lVert\cdot\rVert_{\mathrm{tr}}$, then $(H_n),(d_n)$ is a presentation of an
object of $\Corrtr(\Atr)$ and
$\Tr_A(\mathfrak f^{(k)})\to\Tr_A(\mathfrak f)$.
\end{enumerate}
\end{theorem}

\begin{proof}
(a)  $\lVert H\star f\rVert_{2}\le\lVert H\rVert_{1}\lVert f\rVert_{2}$ is
Young's inequality on the unimodular group $C_S$ with its Haar measure
$\dstar w$; the second inequality is \eqref{eq:trnormineq}.  Multiplicativity
is associativity of convolution on $L^{1}\star L^{2}$.  For the involution,
$\langle H\star f,g\rangle=\langle f,H^{\dagger}\star g\rangle$ by the
substitution $w\mapsto w^{-1}$ and unimodularity, so
$\lambda(H^{\dagger})=\lambda(H)^{*}$.

(b)  This is the second argument of Proposition~\ref{prop:faithful}, which uses
only $\Atr\subseteq L^{1}(C_S)$: if $\lambda(H)=0$ then $\widehat H\widehat
G=0$ almost everywhere for every $G\in L^{1}\cap L^{2}$, so $\widehat H=0$
almost everywhere and $H=0$ in $L^{1}$; composing with the injection
$\Atr\to L^{1}(C_S)$ of Theorem~\ref{thm:banach}(a) gives injectivity on
$\Atr$.

(c)  The identification on $\DA$ is the last sentence of the proof of
Proposition~\ref{prop:faithful}; continuity is (a) and injectivity is (b).

(d)  Every statement listed is an identity in the commutative Banach
$\dagger$-algebra $\Atr$ of Theorem~\ref{thm:banach}(a) involving finitely many
elements: the term descriptions of $\oplus$, shift, cone and $\star$; the
identities of Lemma~\ref{lem:slaws}, which are absolute rearrangements of
finite sums; the invariances of Proposition~\ref{prop:invariance}, which use
only $s(F-F)=0$ and additivity of $s$; and $\Tr_A=\exp\circ\,T\circ s$ with $T$
continuous on $\Atr$ by Theorem~\ref{thm:banach}(c).  None of their proofs uses
compact support or smoothness of the entries, only the algebra structure, the
involution and the trace, so they hold verbatim.  For the closure: the
relations $d_{n+1}\star d_n=0$ are
$\lVert\cdot\rVert_{\mathrm{tr}}$-limits of the relations
$d^{(k)}_{n+1}\star d^{(k)}_n=0$, by submultiplicativity
(Theorem~\ref{thm:banach}(a)), so $(H_n),(d_n)$ is a presented complex over
$\Atr$; and
$\Tr_A(\mathfrak f^{(k)})=\exp T\bigl(\sum_n(-1)^{n}H^{(k)}_n\bigr)\to
\exp T\bigl(\sum_n(-1)^{n}H_n\bigr)=\Tr_A(\mathfrak f)$ by continuity of $T$
and finiteness of the common support.
\end{proof}

\begin{remark}\label{rem:whichcompletion}
Theorem~\ref{thm:Atrcat} is the statement that $\Corrtr$ has an analytic
completion which is again a category of presented complexes, on which every
result of Section~\ref{sec:presented} holds and which is closed under the
limits its norm defines.  Nothing is claimed about limits of presentations
without a common finite support; and the separated completion $\NAd$ of
Theorem~\ref{thm:completion} is a numerical module, not a category completion:
it indexes no objects.
\end{remark}

\begin{remark}\label{rem:whynorm}
The weight $\Lambda$ cannot be dropped.  With the norm
$\lVert F\rVert_{L^{1}}+N(F)$ the spread Tate characters $F_n$ of
\eqref{eq:spread} satisfy $\lVert F_n\rVert_{L^{1}}\le\lVert\vartheta\rVert_1
e^{-na/2}$ and $N(F_n)=O(n^{-1}\log n)$, so their norms tend to zero while
$\lambda_+(F_n)=1$ for every $n$: the Tate characters would be unbounded on the
unit ball and $M$ would have no continuous extension.  With \eqref{eq:trnorm},
$\lVert F_n\rVert_{\mathrm{tr}}\ge\Lambda(F_n)\ge\lambda_+(F_n)=1$, and by
\eqref{eq:trnormineq} the norm of $M$ is at most $1$.  The added term is the
$L^{1}$ norm of the Tate-twisted test.
\end{remark}

\begin{remark}\label{rem:twodets}
Theorem~\ref{thm:banach} concerns the determinant $\exp T$.  It says nothing
about the operator-cutoff germ $z\mapsto\det(1+zR_{\Lambda}(F))$ and its
behaviour as $\Lambda\to\infty$, where a Szeg\H o-type anomaly may occur at the
non-unit shells; that is a different object and it is used nowhere in this
paper.  On the unit block the two germs agree, by part (g) of
Theorem~\ref{thm:banach}.
\end{remark}

\section{Charts, restriction, and the colimit trace}\label{sec:charts}

Everything so far was constructed for a fixed finite set $S$ of places; we
shall call $S$ a \emph{chart}.  In this section we compare different charts
$S\subseteq S'\subseteq S''$.  We write $C_S$, $\DA(S)$, $T_S$, $I_A^{S}$ for
the objects of the previous sections when the chart must be named.

\subsection{The chart system}

\begin{definition}\label{def:chartsystem}
For finite sets of places $S\subseteq S'$ with $\infty\in S$, let
\begin{equation}\label{eq:piTS}
\pi_{S'S}:C_{S'}\longrightarrow C_S
\end{equation}
be the coordinate projection in the presentations of Lemma~\ref{thm:split}: it
is a continuous surjective group homomorphism with \emph{compact} kernel
$K_{S'S}=\prod_{\ell\in S'_f\setminus S_f}\Z_{\ell}^{\times}$, normalized to
total mass one, and $\pi_{S'S}\pi_{S''S'}=\pi_{S''S}$.  We consider the filtered system
\begin{equation}\label{eq:chartsystem}
X:=\bigl(\{C_S\}_S,\{\pi_{S'S}\}_{S\subseteq S'}\bigr),\qquad
Y:=\bigl(\{Y_S\}_S\bigr),\qquad
\mathcal C_X(S):=\DA(S),\qquad \mathcal C_Y(S):=G(S).
\end{equation}
\end{definition}

\begin{theorem}\label{thm:chartsystem}
\begin{enumerate}
\item[(a)] \emph{(Inflation)}  $J_{SS'}:=\pi_{S'S}^{*}:\DA(S)\to\DA(S')$ is an
injective continuous $\dagger$-algebra homomorphism, and
$J_{S'S''}J_{SS'}=J_{SS''}$.
\item[(b)] \emph{(Restriction)}  $R_{S'S}:=(\pi_{S'S})_{*}:\DA(S')\to\DA(S)$,
\begin{equation}\label{eq:RTS}
(R_{S'S}G)(v):=\int_{K_{S'S}}G(\tilde v\,k)\,dk\qquad
\bigl(\tilde v\in\pi_{S'S}^{-1}(v)\ \text{arbitrary}\bigr),
\end{equation}
is a well-defined continuous $\dagger$-algebra homomorphism, and
$R_{S'S}\circ J_{SS'}=\id_{\DA(S)}$.
\item[(c)] \emph{(Functors on objects)}  $E_F\mapsto E_{J_{SS'}F}$ extends to a
monoidal $\dagger$-functor $\Inf_{SS'}:\Corrtr(S)\to\Corrtr(S')$, and
$E_G\mapsto E_{R_{S'S}G}$ to a monoidal $\dagger$-functor $\Rest_{S'S}$ in the
other direction, with $\Rest_{S'S}\Inf_{SS'}=\id$.
\item[(d)] The Tate boundary is strictly invariant: $M\circ J_{SS'}=M$.
\end{enumerate}
\end{theorem}

\begin{proof}
(a)  $\pi_{S'S}$ is a homomorphism, so $\pi_{S'S}^{*}$ preserves convolution up
to the normalization of Haar measure, which is fixed here by giving every
$\Z_{\ell}^{\times}$ total mass one; compactness of the kernel is what makes
$\pi_{S'S}^{*}F$ compactly supported, so $J_{SS'}$ lands in $\DA(S')$, and it
preserves smoothness in the archimedean coordinate, local constancy, and
$\dagger$.  Injectivity follows from (b), and functoriality from
$\pi_{S'S}\pi_{S''S'}=\pi_{S''S}$.

(b)  \emph{Well defined:} the integrand is a function on the coset
$\pi_{S'S}^{-1}(v)$, a translate of the compact group $K_{S'S}$, so
\eqref{eq:RTS} does not depend on the representative.  \emph{Lands in
$\DA(S)$:} $R_{S'S}G$ is supported in $\pi_{S'S}(\supp G)$, which is compact; it
is smooth in the archimedean coordinate because $\pi_{S'S}$ is the identity
there and differentiation passes under the compact integral; and it is locally
constant in the remaining coordinates at the same level as $G$.
\emph{Homomorphism:} pushforward of measures along a surjective homomorphism of
locally compact abelian groups is multiplicative for convolution when the Haar
measures are compatibly normalized, which they are, $\dstar w_{S'}=\dstar
w_S\otimes dk$ with both compact factors of mass one; explicitly, writing
$u=\tilde xk'$ and applying Weil's integration formula,
\[
R(G\star H)(v)=\int_{K}\int_{C_{S'}}G(u)H(u^{-1}\tilde vk)\dstar u\,dk
=\int_{C_S}(RG)(x)(RH)(x^{-1}v)\dstar x .
\]
Compatibility with $\dagger$ is inversion invariance of the Haar measures.
Finally $J_{SS'}F$ is constant on the fibres of $\pi_{S'S}$, so its fibre average
is $F$ itself.

(c)  Both $J$ and $R$ are $\dagger$-algebra homomorphisms, so they carry
presented complexes to presented complexes (apply them to every $F_n$ and to
every differential), preserve $\oplus$, shift, cone and $\star$, and carry
$\dagger$ to $\dagger$; the composite identity is (b) applied termwise.

(d)  For $w\in C_{S'}$ the modulus $|w|_{S'}$ equals the modulus of its archimedean
coordinate, as does $|\pi_{S'S}(w)|_S$; and Haar measure on $C_{S'}$ is
$dt\otimes(\text{probability})$, as on $C_S$.  Hence
\[
\lambda_{\mp}(J_{SS'}F)=\int_{C_{S'}}F(\pi_{S'S}w)|w|_{S'}^{\mp1/2}\dstar w
=\int_{\R}F_{\mathrm{prof}}(t)e^{\mp t/2}\,dt=\lambda_{\mp}(F).\qedhere
\]
\end{proof}

\subsection{The comparison of traces}

\begin{theorem}\label{thm:chartcomp}
Let $S\subseteq S'$ be finite sets of places with $\infty\in S$.  For every
$F\in\DA(S)$,
\begin{equation}\label{eq:chartcomp}
T_{S'}(J_{SS'}F)=T_S(F)+\sum_{\ell\in S'_f\setminus S_f}(\log\ell)
\sum_{k\ne0}\ell^{-|k|/2}F\bigl(\delta_S(\ell^{-k})\bigr),
\end{equation}
where $\delta_S(t):=(t\text{ at }\infty,\ t\text{ at every }p\in S_f)\in C_S$
and the inner sum is finite.
\end{theorem}

\begin{proof}
\emph{Old places.}  Fix $p\in S_f$ and $x=p^{k}u\in\Q_p^{\times}$.  By
\eqref{eq:slices} applied with $S'$ in place of $S$,
$j_{S',p}(x)=(p^{-k}\text{ at }\infty,\ u\text{ at }p,\ p^{-k}\text{ at every
other place of }S')$; applying $\pi_{S'S}$ deletes the new coordinates and leaves
exactly $j_{S,p}(x)$.  So $(J_{SS'}F)\circ j_{S',p}=F\circ j_{S,p}$ as functions
on $\Q_p^{\times}$, and all three ingredients of \eqref{eq:Wp} --- the shell
integral, the unit packet and the value $a(1)$ --- coincide.  The same
computation at $v=\infty$ gives $j_{S',\infty}(x)=(x,1,\dots,1)$ and
$\pi_{S'S}j_{S',\infty}=j_{S,\infty}$, so $W_{\infty}$ agrees.

\emph{New places.}  Fix $\ell\in S'_f\setminus S_f$ and $x=\ell^{k}u$.  As above
\[
j_{S',\ell}(x)=\bigl(\ell^{-k}\text{ at }\infty,\ u\text{ at }\ell,\
\ell^{-k}\text{ at every other place of }S'\bigr),
\]
so
\[
a(x):=(J_{SS'}F)\bigl(j_{S',\ell}(x)\bigr)=F\bigl(\delta_S(\ell^{-k})\bigr),
\]
independent of $u$; hence $a$ is constant on each valuation shell.  For $k=0$
this gives $a(u)=F(\delta_S(1))=F(1)=a(1)$ for every $u\in\Z_{\ell}^{\times}$,
so the unit packet of \eqref{eq:Wp} vanishes identically at $\ell$.  For
$k\ne0$, $\tau_{\ell}(\ell^{k}u)=\ell^{-|k|/2}$ by Lemma~\ref{lem:tau}(i) and
the shell $\ell^{k}\Z_{\ell}^{\times}$ has volume $\log\ell$, so the shell
integral is $(\log\ell)\sum_{k\ne0}\ell^{-|k|/2}F(\delta_S(\ell^{-k}))$.  Only
finitely many $k$ contribute, since $F$ has compact support in the archimedean
coordinate and $|\delta_S(\ell^{-k})|_S=\ell^{-k}$.

Summing over $v\in S'$ and using $(J_{SS'}F)(1)=F(1)$, so that the constant term
$cF(1)$ is the same on both sides by Convention~\ref{conv:c}, gives
\eqref{eq:chartcomp}.
\end{proof}

Thus the trace depends on $S$: enlarging $S$ adds places, and the sum of local
fixed-point contributions acquires the terms of the new places, each prime
entering with its half-density weights.

\begin{proposition}\label{prop:resnotrace}
There are $S\subseteq S'$ and $G\in\DA(S')$ with $T_{S'}(G)\ne T_S(R_{S'S}G)$.
\end{proposition}

\begin{proof}
Take $G:=J_{SS'}F$ for $F\in\DA(S)$ with $F\ge0$ and
$F(\delta_S(\ell^{-1}))>0$ for some $\ell\in S'_f\setminus S_f$.  Then
$R_{S'S}G=F$, while by \eqref{eq:chartcomp}
$T_{S'}(G)=T_S(F)+(\log\ell)\sum_{k\ne0}\ell^{-|k|/2}F(\delta_S(\ell^{-k}))>
T_S(F)=T_S(R_{S'S}G)$.
\end{proof}

The structural reason is visible in \eqref{eq:slices}: the splitting $\varphi$
places the valuation $\ell^{-k}$ at \emph{all} the other coordinates, not at
their fibre average, so the new local trace $W_{\ell}$ samples $F$ at the
discrete points $\delta_S(\ell^{-k})$, which fibre averaging cannot see.

\begin{remark}[What compatibility with change of $S$ holds]\label{rem:restrictionstatus}
Proposition~\ref{prop:resnotrace} is a negative statement and we record it as
such: the trace is \emph{not} compatible with the restriction maps $R_{S'S}$.
What holds is (i) compatibility with inflation up to the explicit correction
\eqref{eq:chartcomp}, a finite sum over the new primes, and (ii) exact
invariance of the colimit trace $T_{\infty}$ under the transition maps
(Theorem~\ref{thm:colimtrace}(c)).  A requirement of the form ``the trace is
compatible with restriction'' is therefore met only in the sense (i)--(ii),
and not literally at a finite chart.
\end{remark}

\subsection{The colimit trace}

\begin{theorem}\label{thm:colimtrace}
For $F\in\DA(S)$ let $R(F):=\sup\{|\log|w|_S|:w\in\supp F\}<\infty$ be its
\emph{radius}.  Then:
\begin{enumerate}
\item[(a)] \emph{(Support window)}  For $\ell\notin S_f$ with $\log\ell>R(F)$
the correction term of \eqref{eq:chartcomp} at $\ell$ vanishes identically.
\item[(b)] \emph{(Stabilization)}  For all finite $S'\supseteq
S_0(F):=S\cup\{\ell:\log\ell\le R(F)\}$ the value $T_{S'}(J_{SS'}F)$ is independent
of $S'$; writing $T_{\infty}(F)$ for this common value,
\begin{equation}\label{eq:Tinf}
T_{\infty}(F)=T_S(F)+\sum_{\ell\notin S_f,\ \log\ell\le R(F)}(\log\ell)
\sum_{k\ne0}\ell^{-|k|/2}F\bigl(\delta_S(\ell^{-k})\bigr),
\end{equation}
a finite sum.
\item[(c)] \emph{(Exact chart invariance)}  $T_{\infty}\circ J_{SS'}=T_{\infty}$
for all $S\subseteq S'$.  Hence on the filtered colimit of $\dagger$-algebras
\begin{equation}\label{eq:colim}
\DA^{\infty}:=\operatorname*{colim}_{S}\DA(S)\qquad
\text{(transition maps }J_{SS'})
\end{equation}
the functional $T_{\infty}$ is well defined, linear and satisfies
$T_{\infty}(H^{\dagger})=\overline{T_{\infty}(H)}$, and
$I_{\infty}(F,G):=T_{\infty}(F\star G^{\dagger})$ is a well-defined Hermitian
form.
\item[(d)] $M$ descends to $\DA^{\infty}$ and the sequence
\eqref{eq:tateseq} holds verbatim over $\DA^{\infty}$.
\item[(e)] \emph{(Window comparison)}  If $R(F\star G^{\dagger})<
\log\ell_{\mathrm{next}}(S)$, where $\ell_{\mathrm{next}}(S)$ is the least
prime not in $S_f$, then $I_{\infty}(F,G)=I_A^{S}(F,G)$ exactly.
\end{enumerate}
\end{theorem}

\begin{proof}
(a)  $\delta_S(\ell^{-k})$ has $|\delta_S(\ell^{-k})|_S=\ell^{-k}$ by
Lemma~\ref{thm:split}, so
$\bigl|\log|\delta_S(\ell^{-k})|_S\bigr|=|k|\log\ell\ge\log\ell>R(F)$ for
$k\ne0$; hence $F$ vanishes there.

(b)  Apply \eqref{eq:chartcomp} for $S\subseteq S'$: by (a) the sum over
$\ell\in S'_f\setminus S_f$ reduces to the finite sum over $\log\ell\le R(F)$,
which is the same for every $S'\supseteq S_0(F)$.

(c)  Both sides equal $T_{S''}(J_{SS''}F)$ for any $S''\supseteq S_0(F)\cup S'$, using
$J_{S'S''}J_{SS'}=J_{SS''}$ and $R(J_{SS'}F)=R(F)$, the chart map preserving the
modulus.  Linearity and reality are inherited from
Proposition~\ref{prop:bilin}.

(d) is Theorem~\ref{thm:chartsystem}(d), together with the fact that a filtered
colimit of $\dagger$-algebras along $\dagger$-algebra maps is a
$\dagger$-algebra, that exactness is preserved by filtered colimits, and that
$M$ is already surjective at each stage.

(e) is (a) applied to $H=F\star G^{\dagger}$ at every $\ell\notin S_f$.
\end{proof}

In other words, $T_S$ is the truncation of $T_{\infty}$ to the places of $S$.
We do not discuss the radical of $I_{\infty}$ on $\DA^{\infty}$;
Theorem~\ref{thm:radical} concerns a fixed $S$.

\section{Effectivity}\label{sec:eff}

\begin{definition}\label{def:eff}
A graph element $\sum_{\gamma,a}b_{\gamma,a}\,\Pi(\cdot)$ is \emph{effective}
if every crossed coefficient $b_{\gamma}$ is a real nonnegative function;
$\Eff_A\subseteq G$ is the class generated by bounded complexes of effective
terms under sums and nonnegative scaling.  A section $F\in\DA$ is
\emph{effective} if $F(u)\ge0$ for every $u\in C_S$, and a presented complex
$(K,\mathfrak f)$ is effective if every $F_n$ is; $\Eff_A^{\mathrm{tr}}
\subseteq\Corrtr$ denotes the class so defined.  Positivity of the Hermitian
form plays no role in these definitions.
\end{definition}

\begin{theorem}\label{thm:cone}
On the fixed-$S$ graph-kernel sector, $\Eff_A$ is:
\begin{enumerate}
\item[(a)] \emph{intrinsic}: an element is effective if and only if all its
half-density row measures $\Pi(b,a)$ are positive;
\item[(b)] closed in the rapid topology and closed under addition and
nonnegative real scaling;
\item[(c)] preserved by $\star$ and by $\dagger$;
\item[(d)] preserved by gauge changes;
\item[(e)] \emph{not} the cone of positive semidefinite operators.
\end{enumerate}
Every $E_w$ is effective, with coefficient $1$ at the single label $(1,w)$.
\end{theorem}

\begin{proof}
(a)  With $\Pi(b,a)=|a|_S^{-1/2}T_{b,a}$ as in \eqref{eq:graphops}, the row
measures of $\Pi(b,a)$ are the measures $b_{\gamma}\,d\mu$ scaled by the
positive factor $|a|_S^{-1/2}$, so nonnegativity of every $b_{\gamma}$ gives
positive row measures.  Conversely, on the conull subset of $X_S$ with nonzero
real coordinate the atoms of the row measures at distinct labels $\gamma$ are
distinct, because the $\Gamma_S$-action is free on that locus --- if
$\gamma x_{\infty}=x_{\infty}$ with $x_{\infty}\ne0$ and $\gamma\in\Q_{>0}$
then $\gamma=1$ --- so positivity of the total row
measure forces nonnegativity of each coefficient almost everywhere, and then
everywhere by continuity.

(b)  Nonnegativity is a closed condition in the rapid seminorms and is stable
under sums and nonnegative scalars pointwise.

(c)  Each crossed coefficient of $\Pi(b,a)\Pi(c,d)=\Pi(b\beta_a(c),ad)$ is a
convergent sum of products of nonnegative functions, hence nonnegative; and
$\Pi(b,a)^{*}=\Pi(\beta_a^{-1}(b^{*}),a^{-1})$ has coefficients
$\overline{b_{\gamma}}\circ\beta_a^{-1}$, nonnegative when the $b_{\gamma}$ are.

(d)  Gauge changes relabel the coefficients.

(e)  The effective permutation kernel exchanging two points is self-adjoint
with eigenvalues $1$ and $-1$, hence effective and not operator-positive;
conversely, with $X=\{1,2\}$, $e$ the identity and $t$ the transposition graph,
$(e-t)^{\dagger}(e-t)=2e-2t$ is a sum of squares but has coefficient $-2$ at
$t$, hence is operator-positive and not effective.
\end{proof}

\begin{proposition}\label{prop:effsmear}
In the trivialization of Corollary~\ref{cor:trivial}, $E_F$ is effective in the
sense of Definition~\ref{def:eff} if and only if $F\ge0$.  Hence the two
notions of effectivity agree along $F\mapsto E_F$, and
$\Eff_A^{\mathrm{tr}}$ is closed under $\oplus$, nonnegative scaling, $\star$
and $\dagger$.
\end{proposition}

\begin{proof}
By Corollary~\ref{cor:trivial} the coefficient of $E_F$ at the label $u$ is
exactly $F(u)$, with no coefficient at a nontrivial gauge label.  By
Theorem~\ref{thm:cone}(a) effectivity is positivity of the half-density row
measures, and the row measure of $E_F$ at $u$ is $F(u)|u|_S^{-1/2}$ times a
positive Haar factor; since $|u|_S^{-1/2}>0$, this is equivalent to $F\ge0$.
Closure under $\oplus$ and nonnegative scaling is pointwise; closure under
$\star$ is $(F\star G)(w)=\int F(u)G(u^{-1}w)\dstar u\ge0$; and since the
coefficients are real, $F\ge0$ gives $F^{\dagger}\ge0$.
\end{proof}

\begin{theorem}\label{thm:properness}
The image of $\Eff_A^{\mathrm{tr}}$ in $N_A$ is a convex cone stable under
$\star$ and $\dagger$, and it is \emph{proper}: if $F$ and $-F$ are both
effective then $F=0$.
\end{theorem}

\begin{proof}
Convexity and stability descend from Proposition~\ref{prop:effsmear} because
$\Nrad$ is a $\star$- and $\dagger$-stable ideal
(Proposition~\ref{prop:ideal}).  Properness in $\DA$ is pointwise: $F(u)\ge0$
and $-F(u)\ge0$ give $F(u)=0$.  By Theorem~\ref{thm:radical}, $\Nrad=\{0\}$, so
$N_A=\DA$ and the quotient map is the identity.
\end{proof}

\begin{theorem}\label{thm:coneproper}
Let $\mathcal E:=\overline{J(\Eff_A^{\mathrm{tr}})}\subseteq\NAd$ be the closed
cone generated by the image of the effective cone.  Then
\begin{equation}\label{eq:properdist}
\mathcal E\cap(-\mathcal E)=\{0\}.
\end{equation}
\end{theorem}

\begin{proof}
$\mathcal E$ is a convex cone: $\Eff_A^{\mathrm{tr}}$ is closed under addition
and nonnegative scaling, $J$ is linear, and the closure of a convex cone is a
convex cone.  Let $x=(\ell,m)\in\mathcal E\cap(-\mathcal E)$ and choose nets
$(F_{\alpha})$, $(G_{\beta})$ in $\Eff_A^{\mathrm{tr}}$ with $J(F_{\alpha})\to
x$ and $J(G_{\beta})\to-x$ in $\DA'\oplus\BT$; nets are needed because the
strong dual of a nuclear LF space need not be metrizable.

\emph{Step 1: the Tate component vanishes.}  The projection $\bar M$ is
continuous, so $m=\lim_{\alpha}M(F_{\alpha})$ and
$-m=\lim_{\beta}M(G_{\beta})$.  For $F\ge0$ one has
$\lambda_{\mp}(F)=\int F|w|_S^{\mp1/2}\dstar w\ge0$, so $M(F)$ lies in the
closed quadrant $\R_{\ge0}^{2}$, and so do both limits.  Hence
$m\in\R_{\ge0}^{2}\cap(-\R_{\ge0}^{2})=\{0\}$, i.e.\ $\lambda_-(F_{\alpha})\to0$
and $\lambda_+(F_{\alpha})\to0$.

\emph{Step 2: the $L^{1}$ norm is controlled by the Tate component on the
cone.}  For every $w$, $|w|_S^{1/2}+|w|_S^{-1/2}\ge2$; hence for $F\ge0$
\begin{equation}\label{eq:amgm}
\lVert F\rVert_{L^{1}}=\int F\dstar w\le\tfrac12\int
F(w)\bigl(|w|_S^{1/2}+|w|_S^{-1/2}\bigr)\dstar w
=\tfrac12\bigl(\lambda_+(F)+\lambda_-(F)\bigr).
\end{equation}
By Step 1, $\lVert F_{\alpha}\rVert_{L^{1}}\to0$, so
$\sup_{\chi}|\widehat F_{\alpha}(\chi)|\to0$.

\emph{Step 3: the distribution component vanishes.}  For every $G\in\DA$,
Lemma~\ref{lem:Nseminorm} gives $|I_A(F_{\alpha},G)|\le
\bigl(\sup|\widehat F_{\alpha}|\bigr)N(G)\to0$.  Since $I_A(F_{\alpha},\cdot)
\to\ell$ pointwise on $\DA$, we get $\ell=0$.  Hence $x=0$.
\end{proof}

\begin{theorem}\label{thm:coneRR}
Let $F,G\in\Eff_A^{\mathrm{tr}}$, i.e.\ $F,G\ge0$ pointwise.
\begin{enumerate}
\item[(a)] \emph{(Positive bidegree)}  $M(F)=(\lambda_-(F),\lambda_+(F))\in
\R_{\ge0}^{2}$, with both entries positive unless $F=0$.  So the cone maps into
the closed positive quadrant of the ruling plane, exactly as effective cycles
have nonnegative bidegree.
\item[(b)] \emph{(Stability)}  $\Eff_A^{\mathrm{tr}}$ is stable under $\oplus$,
nonnegative scaling, $\star$ and $\dagger$, and $M$ carries these to $\oplus$,
scaling, componentwise product and the flip on the quadrant.
\item[(c)] \emph{(Contractivity of the ruling component)}  With the convention
\eqref{eq:metric} and $\lambda_{\mathrm{per}}$ as in
Definition~\ref{def:lper},
\begin{equation}\label{eq:contractive}
\lVert1\rVert^{2}_{\lambda_{\mathrm{per}}(MF,MG)}=
\exp\bigl(-\beta_T(MF,MG)\bigr)
=\exp\bigl(-\lambda_-(F)\lambda_+(G)-\lambda_+(F)\lambda_-(G)\bigr)\le1,
\end{equation}
with equality if and only if $F=0$ or $G=0$.
\item[(d)] \emph{(Scalar component)}  No sign is asserted for $I_A(F,F)$
when $F\ge0$; compare Theorem~\ref{thm:signature}.
\item[(e)] \emph{(Descent)}  The cone descends unchanged to $N_A$
(Theorem~\ref{thm:properness}) and to $\NAd$ as the proper closed cone
$\mathcal E$ (Theorem~\ref{thm:coneproper}).
\end{enumerate}
\end{theorem}

\begin{proof}
(a)  $\lambda_{\mp}(F)\ge0$ for $F\ge0$, and is positive unless $F$ vanishes
almost everywhere, hence everywhere by continuity.  (b) is
Proposition~\ref{prop:effsmear} and Theorem~\ref{thm:tate}; on the positive
quadrant the entries are real, so $\sharp$ is the flip.  (c)
$\beta_T(MF,MG)$ is a sum of products of nonnegative reals by (a); apply
\eqref{eq:metric}.  (e) is Theorems~\ref{thm:properness} and
\ref{thm:coneproper}.
\end{proof}

\begin{remark}\label{rem:twocones}
The cone of Definition~\ref{def:eff} is the operator-kernel analogue of an
effective cycle with nonnegative multiplicities.  It is different from the
$C^{*}$-cone $\overline{\{\sum_iX_i^{\dagger}X_i\}}$, and by
Theorem~\ref{thm:cone}(e) neither contains the other; the lifting of local
positive factorizations is a property of the latter and is not proved here.
\end{remark}

\section{The Riemann--Roch biextension}\label{sec:RR}

\subsection{The biextension}

\begin{definition}\label{def:lper}
On $\BT=\C_-\oplus\C_+$ with the hyperbolic form $\beta_T$ put
\begin{equation}\label{eq:lper}
\lambda_{\mathrm{per}}(b,b'):=\bigl(\C,\ s=\exp\beta_T(b,b')\bigr)
\in\Picb(\C),\qquad\text{so}\quad
\lVert1\rVert^{2}_{\lambda_{\mathrm{per}}(b,b')}=
\exp\bigl(-\Rel\,\beta_T(b,b')\bigr).
\end{equation}
For real $b=(d_1,d_2)$ this is $\lVert1\rVert^{2}=\exp(-2d_1d_2)$, i.e.\
$\lVert1\rVert=\exp(-d_1d_2)$: the ruling line of the classical picture, whose
polarized form is $\Rel\,\beta_T$.  We call $\lambda_{\mathrm{per}}$ the
\emph{intrinsic periodic ruling determinant}; the adjective \emph{periodic} is
justified by Theorem~\ref{thm:ML}(d)--(e), where the Tate functionals enter the
local traces as residues --- at $\infty$ as $\lambda_+$ itself and at $p$ as
its $2\pi/\log p$-periodization \eqref{eq:poisson}.
\end{definition}

\begin{lemma}\label{lem:lperbiext}
$\lambda_{\mathrm{per}}$ is a based Hermitian biextension of $\BT$: biadditivity
of $\beta_T$ gives canonical isometries
$\lambda_{\mathrm{per}}(b+b'',b')\cong\lambda_{\mathrm{per}}(b,b')\otimes
\lambda_{\mathrm{per}}(b'',b')$ and likewise in the second slot, the two orders
agreeing because both send $1\mapsto1\otimes1$; Hermitian symmetry is
$\beta_T(b',b)=\overline{\beta_T(b,b')}$; and
$\lambda_{\mathrm{per}}(b^{\sharp},b'^{\sharp})=
\overline{\lambda_{\mathrm{per}}(b,b')}$.
\end{lemma}

\begin{proof}
All four statements are the corresponding statements for the exponent
$\beta_T$, which is biadditive, Hermitian, and satisfies
$\beta_T(b^{\sharp},b'^{\sharp})=\overline{\beta_T(b,b')}$ by direct
computation from its matrix; based lines whose exponents add have canonically
isomorphic tensor products, matching distinguished generators.
\end{proof}

\begin{definition}\label{def:RR}
For $F,G\in\Atr$ put
\begin{equation}\label{eq:LS}
L_S(F,G):=\Det^{\mathrm{rel}}_S(F\star G^{\dagger})=
\bigl(\C,\ e^{I_A(F,G)}\bigr)\in\Picb(\C),
\end{equation}
and define the \emph{vector-valued Riemann--Roch biextension} by
\begin{equation}\label{eq:lRR}
\lambda_{RR}(F,G):=\bigl(L_S(F,G),\ \lambda_{\mathrm{per}}(MF,MG)\bigr)\in T_A .
\end{equation}
For $x=(\ell,m)\in\NAd$ and $G\in\DA$ put
$\lambda_{RR}(x,G):=\bigl((\C,e^{\ell(G)}),\lambda_{\mathrm{per}}(m,MG)\bigr)$,
and on the ruling plane
$\lambda_{RR}(\iota(b),\iota(b')):=(\mathrm{unit},\lambda_{\mathrm{per}}(b,b'))$.  For
$F\in\DA$ the two definitions agree, since $J(F)=(I_A(F,\cdot),MF)$.  Its
logarithm is the vector-valued numerical pairing
\begin{equation}\label{eq:logRR}
\log\lambda_{RR}(F,G)=\bigl(I_A(F,G),\ \beta_T(MF,MG)\bigr)
=I_{\mathrm{vec}}(F,G).
\end{equation}
\end{definition}

\begin{theorem}\label{thm:scalartrivial}
Fix a finite chart $S$.  Let $(F_{\alpha})$ be any admissible approximation of
$e_+$ and $(G_{\beta})$ any admissible $N$-bounded approximation of $e_-$
(Definition~\ref{def:admissible}).  Then
\begin{equation}\label{eq:trivial}
L_S(F_{\alpha},G_{\beta})=\bigl(\C,e^{I_A(F_{\alpha},G_{\beta})}\bigr)
\longrightarrow(\C,1)=\text{the unit of }\Picb(\C),
\end{equation}
the limit being uniform in the other index, while
$\lambda_{\mathrm{per}}(MF_{\alpha},MG_{\beta})\to
\lambda_{\mathrm{per}}(e_+,e_-)=(\C,e)$, whose metric is
$\lVert1\rVert^{2}=e^{-1}\ne1$.  Consequently the value assigned to the ruling
plane by Definition~\ref{def:RR} is the only value consistent with the dense
test sector, and the scalar determinant $L_S$ cannot carry the intrinsic
ruling determinant, so that a second component is needed.  This is proved for
$L_S$, the determinant of $I_A$; we make no claim about other scalar
determinants that might be built from $\Th_S$.
\end{theorem}

\begin{proof}
Theorem~\ref{thm:rulnullclass}(b) gives $I_A(F_{\alpha},G_{\beta})\to0$
uniformly and $\beta_T(MF_{\alpha},MG_{\beta})\to\beta_T(e_+,e_-)=1$; apply
$\exp$ and Definition~\ref{def:lper}.  The class of admissible pairs is
nonempty by Theorem~\ref{thm:rulnullclass}(c).
\end{proof}

\begin{remark}\label{rem:whyblind}
Three independent reasons, each proved above, explain why the scalar line is
blind to the rulings.  First, the Tate degrees are off the unitary axis:
$\lambda_{\mp}(F)=\widehat F(\mp i/2)$ is a value of an analytic continuation,
and by Theorem~\ref{thm:nosplit}(a) no continuous functional built from values
of $\widehat F$ on the axis can reproduce it, while $I_A(F,G)=\int\widehat
F\overline{\widehat G}(\Th+c)d\chi$ is built entirely from values on the axis.
Second, the symbol overcounts the Tate poles: by \eqref{eq:residues} the residue
of $\Th_S$ at $\pm i/2$ is $\mp i(|S_f|+1)$, one copy per place, so any
normalization of $I_A^{S}$ making the pair $(\lambda_-,\lambda_+)$ appear with
hyperbolic coefficient $1$ is chart-dependent, which a canonical ruling form
cannot be.  Third, the rulings \emph{are} the radical of the scalar component:
Theorem~\ref{thm:rulnull}(a), and a scalar form is by definition zero on its own
radical.
\end{remark}

\begin{remark}\label{rem:notafunctor}
One may ask for a functor $\Phi$ on $\Corrtr$, with values in the ruling
category, which carries $\Tr_A$ to $\lambda_{\mathrm{per}}$.  If the value of
$\Phi$ on a presented complex $K$ depends only on $M(s(K))\in\BT$, no such
functor exists: by Theorem~\ref{thm:separate} the two one-term complexes with
presentations $F$ and $F'$ have $M(s(K))=M(s(K'))$ and
$T(F)-T(F')\ne0$, so $\Tr_A(K)\ne\Tr_A(K')$.  This is not a pathology of the
present theory: on a surface the trace $\Tr(D)=D\cdot\Delta$ of a
correspondence is not a function of its bidegree, and no functor to the ruling
plane carries it.  What is meaningful, and what is proved next, is that the
Riemann--Roch determinant on the mixed sector \emph{extends} the intrinsic
ruling determinant, the ruling plane being a canonical quotient of $N_A$ and a
canonical subobject of $\NAd$.
\end{remark}

\begin{theorem}\label{thm:RR}
With the target $T_A=\Picb(\BT)$ of Definition~\ref{def:target} and the metric
convention \eqref{eq:metric}:
\begin{enumerate}
\item[(a)] \emph{(Biextension laws)}  $\lambda_{RR}$ is a based Hermitian
biextension: additive in the first slot, conjugate-additive in the second,
Hermitian under exchange, with the interchange and associativity identities;
and $\dagger$-compatible,
$\lambda_{RR}(F^{\dagger},G^{\dagger})=\overline{\lambda_{RR}(F,G)}$ with
$\sharp$ acting on the second component.
\item[(b)] \emph{(Convergence and continuity)}  Both components are defined and
continuous on $\Atr\times\Atr$ for $\lVert\cdot\rVert_{\mathrm{tr}}$; on
$\NAd\times\DA$ the pairing is separately continuous; and the determinant germ
$z\mapsto\Delta_F(z)$ is holomorphic on $|z|<1/\lVert F\rVert_{L^{1}}$ and
exactly multiplicative.
\item[(c)] \emph{(Principal invariance)}
$\lambda_{RR}(a\star F\star a^{-1},G)\cong\lambda_{RR}(F,G)$ for every
invertible multiplier $a$ whose commutators remain in the trace domain.
\item[(d)] \emph{(Restriction)}  Along the chart maps $J_{SS'}$ the second
component is strictly invariant,
$\lambda_{\mathrm{per}}(MJ_{SS'}F,MJ_{SS'}G)=\lambda_{\mathrm{per}}(MF,MG)$; the
first transforms by the exact comparison \eqref{eq:chartcomp}; and on the
colimit $\DA^{\infty}$ with $T_{\infty}$ both components are chart-invariant.
\item[(e)] \emph{(Extension of the ruling determinant; the equality formula)}
In the canonical splitting \eqref{eq:split},
\[
\lambda_{RR}\bigl|_{\iota(\BT)\times \iota(\BT)}=(\mathrm{unit},\ 
\lambda_{\mathrm{per}}),\qquad
\lambda_{RR}\bigl|_{N_A^{\mathrm{dist,prim}}\times\DA^{\mathrm{prim}}}
=(L_S,\ \mathrm{unit}),
\]
and for every $F,G\in\DA$, writing $J(F)=(I_A(F,\cdot),0)+\iota(MF)$,
\begin{equation}\label{eq:weilformula}
\lambda_{RR}(F,G)=\lambda_{RR}\bigl((I_A(F,\cdot),0),G\bigr)\otimes
\lambda_{\mathrm{per}}(MF,MG)
\end{equation}
in $T_A$; equivalently, in logarithms,
$I_{\mathrm{vec}}(F,G)=(I_A(F,G),0)+(0,\beta_T(MF,MG))$.  So the Riemann--Roch
determinant restricts on the ruling plane to $\lambda_{\mathrm{per}}$, and the
rulings enter a mixed intersection only through their Deligne pairing.  This
has the shape of Weil's decomposition $D\cdot D'=P\cdot P'+d_1d_2'+d_2d_1'$,
with two differences: \eqref{eq:weilformula} is a pair of components and not a
sum, and its first component is $I_A(F,G)$ for the whole of $F$ and $G$, not a
pairing of primitive parts.  The analogy is formal.
\item[(f)] \emph{(Local metric comparison)}  Fix $q\in S_f$ and let
$L_S^{(q)}(F,G):=(\C,e^{W_q(F\star G^{\dagger})})$ be the local factor of
$L_S$, and write $L^{(q)}(F):=L_S^{(q)}(F,F)$ for its diagonal.  With the
convention \eqref{eq:metric} and the normalization dictionary of
Remark~\ref{rem:normdict}, $L_S^{(q)}$ is the tangential determinant line with
$\lVert1\rVert^{2}=\exp(-\Rel\,W_q(F\star G^{\dagger}))$, and the cross-effect
$P(F,G):=L^{(q)}(F+G)\otimes L^{(q)}(F)^{-1}\otimes L^{(q)}(G)^{-1}$ has
$\lVert1\rVert=\exp(-\Rel\,W_q(F\star G^{\dagger}))$.
Moreover the all-amplitude limit holds at the level of the symbol: for
$H:=F\star F^{\dagger}$, $c=0$ and $H_N:=H/N$,
\begin{equation}\label{eq:allamp}
N\sum_{r\ge1}(-1)^{r+1}\frac{T(H_N^{\star r})}{r}=T(H)+O(1/N)
=I_A(F,F)+O(1/N),
\end{equation}
the error being at most $\lVert H\rVert_{L^{1}}\lVert
H\rVert_{\mathrm{tr}}\bigl/\bigl(2N(1-\lVert H\rVert_{L^{1}}/N)\bigr)$ for
$N>\lVert H\rVert_{L^{1}}$.
\item[(g)] \emph{(Signature)}  The scalar component is indefinite with
infinite index on both sides (Theorem~\ref{thm:signature}); in particular
\eqref{eq:weilformula} is an identity of based lines and carries no
inequality of Hodge-index type.
\end{enumerate}
\end{theorem}

\begin{proof}
(a)  \emph{Additivity.}  $I_A$ is additive in $F$ and conjugate-additive in $G$
(Proposition~\ref{prop:bilin}), and so is $\beta_T(M\cdot,M\cdot)$, $M$ being
linear; based lines with adding exponents give canonical isometries in each
component, both sides being $(\C,\cdot)$ with matching distinguished
generators.  \emph{Interchange and associativity.}  The two composites
expressing $\lambda_{RR}(F+F',G+G')$ as a tensor product of four pieces, built
in the two possible orders, both send $1$ to $1\otimes1\otimes1\otimes1$ and
have equal exponents, the exponent being biadditive; so they agree.
\emph{Hermitian.}  $I_A(G,F)=\overline{I_A(F,G)}$ by
Proposition~\ref{prop:bilin} and $\beta_T(m',m)=\overline{\beta_T(m,m')}$.
\emph{Involution.}  $(F^{\dagger})^{\wedge}=\overline{\widehat F}$ and $\Th+c$
is real, so $I_A(F^{\dagger},G^{\dagger})=\overline{I_A(F,G)}$ by
\eqref{eq:IAfourier}; and
$\beta_T(MF^{\dagger},MG^{\dagger})=\beta_T((MF)^{\sharp},(MG)^{\sharp})=
\overline{\beta_T(MF,MG)}$ by Lemma~\ref{lem:lperbiext}.

(b)  First component:
$|I_A(F,G)|=|T(F\star G^{\dagger})|\le N(F\star G^{\dagger})\le
\lVert F\star G^{\dagger}\rVert_{\mathrm{tr}}\le\lVert F\rVert_{\mathrm{tr}}
\lVert G\rVert_{\mathrm{tr}}$ by Theorem~\ref{thm:banach}(a),(c), and $\exp$ is
continuous.  Second component: by Theorem~\ref{thm:banach}(b), $M$ is
continuous $\Atr\to\BT$ and $\beta_T$ is a continuous bilinear form, so
$(F,G)\mapsto\lambda_{\mathrm{per}}(MF,MG)$ is continuous.  On
$\NAd\times\DA$ the pairing $((\ell,m),G)\mapsto(\ell(G),\beta_T(m,MG))$ is
separately continuous by Theorem~\ref{thm:completion}.  On the ruling plane the
value is Definition~\ref{def:RR}'s, consistent with the test sector along every
admissible approximation by Theorem~\ref{thm:rulnullclass}; no continuous
extension of $I_A$ to $\NAd\times\NAd$ is claimed
(Remark~\ref{rem:noextension}).  Holomorphy and multiplicativity of the germ are
Theorem~\ref{thm:banach}(e),(f).

(c)  $T(a\star F\star a^{-1})=T(F)$ by cyclicity
(Theorem~\ref{thm:cyclic}).  Since $\DA$ is commutative the conjugation is
trivial there; for a noncommutative enlargement the commutator condition
would have to be checked.  The second component is unchanged because
$M(a\star F\star a^{-1})=M(a)M(F)M(a)^{-1}=M(F)$ in the commutative $\BT$.

(d)  Second component: Theorem~\ref{thm:chartsystem}(d).  First component:
\eqref{eq:chartcomp} applied to $H=F\star G^{\dagger}$, using that $J_{SS'}$ is a
$\dagger$-algebra map, so $J_{SS'}(F\star G^{\dagger})=(J_{SS'}F)\star
(J_{SS'}G)^{\dagger}$.  Colimit: Theorem~\ref{thm:colimtrace}(c),(d).

(e)  By Theorem~\ref{thm:splitting}, $J(F)=(I_A(F,\cdot),0)+\iota(MF)$ inside
$\NAd$.  Evaluating Definition~\ref{def:RR} on each summand,
$\lambda_{RR}((I_A(F,\cdot),0),G)=\bigl((\C,e^{I_A(F,G)}),
\lambda_{\mathrm{per}}(0,MG)\bigr)=(L_S(F,G),\mathrm{unit})$ since
$\beta_T(0,\cdot)=0$; and
$\lambda_{RR}(\iota(MF),G)=\bigl((\C,e^{0}),\lambda_{\mathrm{per}}(MF,MG)\bigr)
=(\mathrm{unit},\lambda_{\mathrm{per}}(MF,MG))$ since the first component of
$\iota(m)$ is $0$.  The tensor product of these two elements of
$T_A=\Picb(\C)^{2}$ multiplies the distinguished generators componentwise,
giving $(L_S(F,G),\lambda_{\mathrm{per}}(MF,MG))=\lambda_{RR}(F,G)$, which is
\eqref{eq:weilformula}.  The two displayed restrictions are the same two
computations, read on $\iota(\BT)\times \iota(\BT)$, where the first component is the
unit, and on $N_A^{\mathrm{dist,prim}}\times\DA^{\mathrm{prim}}$, where $MG=0$
so the second component is the unit.  Taking logarithms gives the additive
form.

(f)  The identification of $L_S^{(q)}$ with the tangential line is
\eqref{eq:metric} together with the normalization dictionary of
Remark~\ref{rem:normdict}: $\DA$ is commutative, so $F\star G^{\dagger}=
G^{\dagger}\star F$, and the substitution $F\mapsto|\cdot|_S^{1/2}F$ identifies
the unnormalized involution with the half-density one.  For the cross-effect, the exponent of $P(F,G)$ is
\[
W_q\bigl((F+G)\star(F+G)^{\dagger}\bigr)-W_q(F\star F^{\dagger})
-W_q(G\star G^{\dagger})=W_q(F\star G^{\dagger})+W_q(G\star F^{\dagger})
=2\,\Rel\,W_q(F\star G^{\dagger}),
\]
by biadditivity of $(F,G)\mapsto F\star G^{\dagger}$ and
$W_q(H^{\dagger})=\overline{W_q(H)}$, which is
Proposition~\ref{prop:bilin} applied place by place through
\eqref{eq:plancherelv} and the reality of $\Wh_q$ (Theorem~\ref{thm:symbol}(a));
so $\lVert1\rVert^{2}_{P(F,G)}=\exp(-2\Rel W_q(F\star G^{\dagger}))$ by
\eqref{eq:metric}.  For \eqref{eq:allamp}: by
\eqref{eq:powerest}, $|T(H_N^{\star r})|\le\lVert H_N\rVert_{L^1}^{r-1}
\lVert H_N\rVert_{\mathrm{tr}}$ with $\lVert H_N\rVert_{L^1}=\lVert
H\rVert_{L^1}/N$ and $\lVert H_N\rVert_{\mathrm{tr}}=\lVert
H\rVert_{\mathrm{tr}}/N$.  The term $r=1$ is $T(H_N)=T(H)/N$, so
$N\,T(H_N)=T(H)=I_A(F,F)$ exactly, and
\[
N\Bigl|\sum_{r\ge2}(-1)^{r+1}\frac{T(H_N^{\star r})}{r}\Bigr|
\le N\sum_{r\ge2}\frac{(\lVert H\rVert_{L^1}/N)^{r-1}
(\lVert H\rVert_{\mathrm{tr}}/N)}{r}
\le\frac{\lVert H\rVert_{\mathrm{tr}}\lVert H\rVert_{L^1}}
{2N\bigl(1-\lVert H\rVert_{L^1}/N\bigr)},
\]
valid for $N>\lVert H\rVert_{L^1}$.  Note that $L_S$ itself is not the local
factor: $S$ always contains $\infty$, so
$L_S=L_S^{(\infty)}\otimes\bigotimes_{p\in S_f}L_S^{(p)}$ carries an
archimedean factor that a one-place construction does not.

(g) is Theorem~\ref{thm:signature}.
\end{proof}

\begin{remark}\label{rem:whatRRis}
Once $\lambda_{RR}$ is defined as the pair
$(L_S,\lambda_{\mathrm{per}}\circ(M\times M))$, part (e) is formal, and
\eqref{eq:weilformula} in logarithms is biadditivity.  The content lies in the
results on which the definition rests: Theorem~\ref{thm:rulings}, that the two rulings are objects of the
theory, explicit limits of spread Tate characters with $\DA$ acting on them by
the characters $\lambda_{\pm}$; Theorem~\ref{thm:splitting}, that the
completion splits canonically as a $\DA$-module, so that the ruling plane is a
subobject and not only a quotient; Theorem~\ref{thm:scalartrivial}, that the
analytic determinant $L_S$ is trivial on that subobject, which is why $L_S$
alone does not suffice; and the laws (a)--(d), (f).
\end{remark}

\subsection{Base change along the Tate boundary}

\begin{theorem}\label{thm:basechange}
Let $\DA^{+}$ be the unitization of $\DA$ and $M:\DA^{+}\to\BT$ the unital
extension of $(\lambda_-,\lambda_+)$, with $1\mapsto(1,1)$.  Then:
\begin{enumerate}
\item[(a)] $M$ is a surjective unital $\dagger$-algebra homomorphism, carrying
$\dagger$ to $\sharp$.
\item[(b)] \emph{(Base change)}  $\Phi_{\det}:=-\otimes_{\DA^{+}}\BT$, from
bounded complexes of based free $\DA^{+}$-modules with differentials over
$\DA^{+}$ --- the presented complexes of Definition~\ref{def:Corrtr} written in
matrix form --- to bounded complexes of based free $\BT$-modules, is exact,
strictly monoidal on free complexes, carries $\dagger$ to $\sharp$, carries
$\Delta$ to the unit, and has Euler class $M\circ s$.
\item[(c)] $M$ is \emph{not} a homological epimorphism: with $I:=\ker M$, and
with $\operatorname{Tor}$ and $I^{2}$ taken in the purely algebraic sense over
the ring underlying $\DA^{+}$ --- no completed tensor product is used ---
\begin{equation}\label{eq:tor}
\operatorname{Tor}_1^{\DA^{+}}(\BT,\BT)\cong I/I^{2},
\end{equation}
and $I/I^{2}$ surjects onto $\C^{2}$; hence $\operatorname{Tor}_1\ne0$ and is
at least two-dimensional.
\end{enumerate}
\end{theorem}

\begin{proof}
(a)  Multiplicativity and surjectivity of $(\lambda_-,\lambda_+)$ on $\DA$ are
Theorem~\ref{thm:tate}; the unital extension is multiplicative because
$M(1)=(1,1)$ is the unit of $\BT$; and
$M(F^{\dagger})=(\overline{\lambda_+(F)},\overline{\lambda_-(F)})=M(F)^{\sharp}$.

(b)  A based free module $(\DA^{+})^{n}$ satisfies
$(\DA^{+})^{n}\otimes_{\DA^{+}}\BT=\BT^{n}$, with basis carried to basis; a
differential given by a matrix over $\DA^{+}$ is carried to the matrix of its
$M$-images, and $d^{2}=0$ is preserved because $M$ is a ring homomorphism.  The
functor $-\otimes_{\DA^{+}}\BT$ is right exact always and exact on free
modules, so a degreewise split short exact sequence of free complexes goes to a
degreewise split short exact sequence.  Strict monoidality:
$(\DA^{+})^{m}\otimes_{\DA^{+}}(\DA^{+})^{n}=(\DA^{+})^{mn}$ maps to
$\BT^{mn}=\BT^{m}\otimes_{\BT}\BT^{n}$ on the nose with bases matching, and the
associativity and unit constraints are identities on both sides; $\Delta$ has
rank one and goes to $\BT$.  The Euler class statement is linearity of $M$
applied to \eqref{eq:s}.

(c)  For any ring $R$ and any surjection $R\to R/I$, tensoring
$0\to I\to R\to R/I\to0$ with $R/I$ gives
$\operatorname{Tor}_1(R,R/I)=0$, hence
$\operatorname{Tor}_1(R/I,R/I)=\ker(I\otimes_RR/I\to R/I)=I\otimes_RR/I=I/I^{2}$,
the displayed map being zero since $I\cdot(R/I)=0$.  Apply this with
$R=\DA^{+}$ and $R/I=\BT$.

For the surjection onto $\C^{2}$, extend the symbol to $\DA^{+}$ by
$(a\cdot1+F)^{\wedge}:=a+\widehat F$, the unit having symbol $1$, so that
$M(a1+F)=\bigl((a+\widehat F)(-i/2),(a+\widehat F)(i/2)\bigr)$; and define
$\partial_{\pm}:I\to\C$ by $\partial_{\pm}(a1+F):=\partial_{\tau}\widehat
F(\pm i/2)$, using that $\widehat F$ extends to an entire function of $\tau$ on
each component and evaluating on $\mathfrak c_0$.  These are linear, and they
kill $I^{2}$: if $X=a1+F$ and $Y=b1+G$ lie in $I$ then $\widehat X(\pm
i/2)=\widehat Y(\pm i/2)=0$, so $(X\star Y)^{\wedge}=\widehat X\widehat Y$ has
a double zero at $\pm i/2$ and $\partial_{\pm}(X\star Y)=\widehat X(\pm
i/2)\partial_{\pm}Y+\widehat Y(\pm i/2)\partial_{\pm}X=0$; by bilinearity they
kill all of $I^{2}$.

They are jointly surjective.  Restrict to the sections supported on
$\mathfrak c_0$ with profile $f\in C_c^{\infty}(\R)$, for which
$\widehat F(\tau)=\int f(t)e^{-i\tau t}dt$, so that $\widehat F(\mp
i/2)=\int f(t)e^{\mp t/2}dt$ and
$\partial_{\pm}(F)=-i\int tf(t)e^{\pm t/2}dt$.  The four functionals
$f\mapsto\int fe^{t/2},\ \int fe^{-t/2},\ \int tfe^{t/2},\ \int tfe^{-t/2}$ on
$C_c^{\infty}(\R)$ are linearly independent, because the four locally
integrable functions $e^{t/2},e^{-t/2},te^{t/2},te^{-t/2}$ are linearly
independent on $\R$: a nontrivial relation would make a nonzero function
$(\alpha+\beta t)e^{t/2}+(\gamma+\delta t)e^{-t/2}$ vanish identically, which is
impossible on inspecting $t\to\pm\infty$.  Hence the joint map
$C_c^{\infty}(\R)\to\C^{4}$ given by the four functionals is surjective, so the
preimage of $\{0\}\times\{0\}\times\C^{2}$ maps onto $\C^{2}$ under the last
two coordinates; and that preimage consists exactly of the profiles with
$\widehat F(\pm i/2)=0$, that is, of elements of $I$.  So
$(\partial_+,\partial_-):I\to\C^{2}$ is surjective and kills $I^{2}$, whence
$I/I^{2}\twoheadrightarrow\C^{2}$.
\end{proof}

Thus derived base change along the Tate boundary does not commute with
composition without a correction; but the correction relevant to
Theorem~\ref{thm:RR}(e) is not a Tor determinant.  It is the direct sum
decomposition \eqref{eq:split} of the completion, in which the primitive and
ruling contributions simply add.

\begin{proposition}\label{prop:contactblind}
Let $\Phi$ be any functor defined on $G$ by taking the determinant line of the
derived diagonal contact of its objects.  Then $\Phi$ is constant on the
finite-place data of the graph sector, and therefore does not determine
$\Tr_A$.
\end{proposition}

\begin{proof}
By Theorem~\ref{thm:shadow}(c) and Remark~\ref{rem:notorsion}, at a finite
place the contact complex of a twist $E_w$ with $v_p(w_p)\ne0$ is the single
one-dimensional class $\C[0]$ of the fixed point $x_p=0$, with no differential
and independent of $w_p$; such a complex has determinant line canonically $\C$
and no torsion.  For a unit twist $\HH_0$ is \eqref{eq:D2}, of infinite
dimension, so no determinant line is defined.  Meanwhile $\Tr_A$ separates classes by
Theorem~\ref{thm:separate}.
\end{proof}

The arithmetic content of $\Tr_A$ lives, as Remark~\ref{rem:notorsion} explains,
in the half-density normalization of the point class, not in any
complex-theoretic invariant of it: the number $\tau_p(\sigma)$ is the torsion of
the \emph{smooth model} $[\R[z]\xrightarrow{(\sigma-1)z}\R[z]]$, which the
totally disconnected coefficient algebra does not carry.  The functor
$\Phi_{\det}$ of Theorem~\ref{thm:basechange} is not of the contact type: it is
base change along a ring map.

\section{The classical model}\label{sec:classical}

The formalism of Sections~\ref{sec:cat}--\ref{sec:dagger} has a model in
perfect kernels on a curve.  The comparison is not used elsewhere in the
paper; it shows that the formal structure is the classical one.

\begin{theorem}\label{thm:classical}
Let $C$ be a smooth projective geometrically connected curve of genus $g$ over
a finite field $k$, and let $h$ be a theta characteristic, $h^{2}\cong\omega_C$.
Theta characteristics exist over the algebraic closure, and each is defined
over a finite extension of $k$; we assume, replacing $k$ by such an extension
and $C$ by its base change if necessary, that $h$ is defined over $k$.  This is
an assumption on the field of definition, not a property of $C$ proved here.  Write
$Y=C\times C$, $H=h\boxtimes h=p_1^{*}h\otimes p_2^{*}h$, let $\sigma$ be the
swap, $K^{\vee}=R\!\Hom(K,\mathcal O_Y)$ and
$D_h(K)=\sigma^{*}(K^{\vee})\otimes H[1]$.  Let $\phi:\Lambda\to\Aut(C)$ be a
homomorphism from an abelian group $\Lambda$, and let $G_{\Lambda}$ be the
differential graded category of bounded complexes of finite sums of objects
$E_{\lambda}$, $\lambda\in\Lambda$, with the composition and involution of
Sections~\ref{sec:cat}--\ref{sec:dagger}, i.e.\ $G$ with its external class
group replaced by $\Lambda$.  Then the additive functor
\[
\Phi:G_{\Lambda}\longrightarrow\Perf(C\times C),\qquad
\Phi(E_{\lambda}):=\mathcal O_{\Gamma_{\phi(\lambda)}},
\]
exact for cones and shifts, satisfies
\begin{enumerate}
\item[(a)] $\Phi(K\star L)\cong\Phi(K)\star\Phi(L)$, the composition on the
right being pull--tensor--push over $C^{3}$;
\item[(b)] $\Phi(\Delta)\cong\mathcal O_{\Delta}$, a two-sided unit;
\item[(c)] for every $\lambda$, writing $f=\phi(\lambda)$ and
$i_{f}:C\to C\times C$, $x\mapsto(x,f(x))$,
\begin{equation}\label{eq:dualitygraph}
D_h\bigl(\Phi(E_{\lambda})\bigr)\ \cong\
\Phi(E_{\lambda^{-1}})\otimes p_1^{*}h\otimes p_2^{*}h^{-1}
\ =\ i_{f^{-1}*}\bigl(h\otimes(f^{-1})^{*}h^{-1}\bigr):
\end{equation}
the transposed graph carrying the canonical half-density twist.  The twist is
trivial exactly on the diagonal, so $D_h(\mathcal O_{\Delta})\cong
\mathcal O_{\Delta}$;
\item[(d)] $\chi(L\Delta^{*}\mathcal O_{\Delta})=2-2g=\Delta\cdot\Delta$, and
$\chi(L\Delta^{*}\Phi(E_{\lambda}))=0$ whenever $\phi(\lambda)$ is a
fixed-point-free automorphism.
\end{enumerate}
\end{theorem}

\begin{proof}
Throughout we use the standard six-functor formalism for perfect kernels on
smooth proper schemes --- derived pullback, proper pushforward, the projection
formula and derived base change \cite{SGA6} --- together with the fundamental
local isomorphism for a regular closed immersion $i:Z\hookrightarrow X$ of
codimension $c$ with normal bundle $N$,
\begin{equation}\label{eq:fli}
R\!\Hom_{\mathcal O_X}\bigl(i_{*}\mathcal O_Z,\mathcal O_X\bigr)
\ \cong\ i_{*}\bigl(\det N\bigr)[-c],
\end{equation}
for which see \cite[Ch.~III, \S7]{HartshorneRD}.

(a)  For automorphisms $f,g$ of $C$ the relevant fibre product
$\{(x,y,z):y=f(x),\,z=g(y)\}$ is isomorphic to $C$, and $p_{13}$ identifies it
with $\Gamma_{g\circ f}$; both projections from the fibre product are
isomorphisms, so the intersection is transverse, there is no excess and no
derived correction, and $\mathcal O_{\Gamma_f}\star\mathcal O_{\Gamma_g}\cong
\mathcal O_{\Gamma_{g\circ f}}$.  Since $\phi$ is a homomorphism,
$\Phi(E_{\lambda}\star E_{\mu})=\Phi(E_{\lambda\mu})$.  On complexes both sides
are the term-by-term totalization with the same Koszul signs, so the
isomorphism is degreewise and compatible with the differentials.

(b)  $\phi(1)=\id$ and $\Gamma_{\id}=\Delta$; and $\mathcal O_{\Delta}$ is a
two-sided unit for the composition, by derived base change.

(c)  By Definition~\ref{def:dagger}, $E_{\lambda}^{\dagger}=E_{\lambda^{-1}}$,
so we must compute $D_h(\mathcal O_{\Gamma_f})$.  The map $i_f$ is a regular
closed immersion of codimension one; the automorphism $\id\times f^{-1}$ of
$C\times C$ carries $\Gamma_f$ isomorphically onto $\Delta$, so, under the
identification $\Gamma_f\cong C$ by the first projection, the normal bundle of
$\Gamma_f$ is that of the diagonal, namely $T_C$.  By \eqref{eq:fli},
\[
\mathcal O_{\Gamma_f}^{\vee}\cong i_{f*}\bigl(T_C\bigr)[-1].
\]
Next, $\sigma\circ i_f=i_{f^{-1}}\circ f$, because
$\sigma(x,f(x))=(f(x),x)=i_{f^{-1}}(f(x))$; and $\sigma$ is an involution, so
$\sigma^{*}=\sigma_{*}$ and, $f$ being an isomorphism,
\[
\sigma^{*}i_{f*}(G)\cong i_{f^{-1}*}\bigl((f^{-1})^{*}G\bigr)
\qquad\text{for every line bundle }G\text{ on }C .
\]
Finally $i_{f^{-1}}^{*}H=h\otimes(f^{-1})^{*}h$, so by the projection formula
\[
D_h(\mathcal O_{\Gamma_f})
=\sigma^{*}\bigl(\mathcal O_{\Gamma_f}^{\vee}\bigr)\otimes H[1]
\cong i_{f^{-1}*}\Bigl((f^{-1})^{*}T_C\otimes h\otimes(f^{-1})^{*}h\Bigr)
\cong i_{f^{-1}*}\Bigl(h\otimes(f^{-1})^{*}\bigl(T_C\otimes h\bigr)\Bigr).
\]
Since $h^{2}\cong\omega_C$ we have $T_C\otimes h\cong\omega_C^{-1}\otimes
h\cong h^{-1}$, whence
$D_h(\mathcal O_{\Gamma_f})\cong i_{f^{-1}*}\bigl(h\otimes
(f^{-1})^{*}h^{-1}\bigr)$, which is \eqref{eq:dualitygraph}; the second
equality there is the projection formula applied to
$\mathcal O_{\Gamma_{f^{-1}}}\otimes p_1^{*}h\otimes p_2^{*}h^{-1}$.  For
$f=\id$ the twist is $h\otimes h^{-1}\cong\mathcal O_C$, so
$D_h(\mathcal O_{\Delta})\cong\mathcal O_{\Delta}$.  This is the only place
where the theta characteristic is used.

The half-density twist in \eqref{eq:dualitygraph} cannot be removed by
renormalizing $\Phi$.  If one sets $\Psi(E_{\lambda}):=i_{f*}(h^{a}\otimes
f^{*}h^{b})$ for fixed exponents $a,b$, the computation above gives
$D_h(\Psi(E_{\lambda}))\cong i_{f^{-1}*}\bigl((f^{-1})^{*}h^{-a-1}\otimes
h^{1-b}\bigr)$, while $\Psi(E_{\lambda^{-1}})=i_{f^{-1}*}\bigl(h^{a}\otimes
(f^{-1})^{*}h^{b}\bigr)$; equality would force $a=1-b$ and $b=-a-1$, i.e.\
$0=-2$.  So $\dagger$ and $D_h$ agree on objects only up to the transposed
half-density, which is the content of \eqref{eq:dualitygraph}.  Without that
twist the identification fails already for the Frobenius, where the cycle-level
inverse of $\Gamma_F$ is $q^{-1}\Gamma_F^{t}$ and not $\Gamma_F^{t}$; the
present functor is defined on automorphisms, where $\deg\phi(\lambda)=1$ and no
such factor occurs.

(d)  $L\Delta^{*}\mathcal O_{\Delta}\cong\mathcal O_C\oplus\Omega^{1}_C[1]$,
the self-intersection of a regular embedding of codimension one with the excess
retained; so $\chi=(1-g)-(g-1)=2-2g$, which is $\deg T_C=\Delta\cdot\Delta$.  If
$\phi(\lambda)$ has no fixed point then $\Gamma_{\phi(\lambda)}\cap\Delta=
\emptyset$, the derived restriction is the zero complex, and its Euler
characteristic is $0$.  No Lefschetz fixed point formula is used as an input:
the value is read off the empty intersection.
\end{proof}

\begin{remark}\label{rem:classicalscope}
Three comments on the scope of this comparison.  First, $\Phi$ is a comparison for the
formalism, on a surrogate category built from an abstract abelian group
$\Lambda$ with a homomorphism to $\Aut(C)$; it is not a functor out of $G$
itself, whose external group is $C_S$, and no homomorphism $C_S\to\Aut(C)$ is
constructed or needed.  The statement is that the formalism of
Sections~\ref{sec:cat}--\ref{sec:dagger} has a model in perfect kernels on a
curve.  Second, $\dagger$ is antilinear on morphisms, being the Hilbert adjoint,
while $D_h$ is a $\C$-linear duality; so \eqref{eq:dualitygraph} compares the
two involutions on objects and on the composition law, not their antilinear
structures.  Third, no comparison with non-invertible
correspondences such as $\Gamma_{F^{k}}$, of bidegree $(q^{k},1)$, is provided:
$G$ contains only graphs of invertible classes, and for a map of degree $>1$
the transpose is not the inverse.
\end{remark}

\begin{remark}\label{rem:effcone}
The effective cone of Definition~\ref{def:eff} is \emph{not} identified with
the cone of effective graph cycles, and no such identification is used.  With
$X=\{1,2\}$, $e$ the identity and $t$ the transposition graph,
$(e-t)^{\dagger}(e-t)=2e-2t$ is a sum of squares with coefficient $-2$ at $t$,
while $t$ is effective with eigenvalues $1$ and $-1$; containment fails in both
directions, and passage to the numerical quotient does not repair it.
Section~\ref{sec:RR} never uses a sum-of-squares cone.
\end{remark}

\begin{remark}\label{rem:invertible}
Every generator of $G$ is the graph bimodule of an \emph{invertible} external
class, and the theory is constructed on that sector.  Three of the four
ingredients of the classical argument are nevertheless present.  The
\emph{bidegree} is supplied by Theorem~\ref{thm:tate}: every generator has
\begin{equation}\label{eq:bidegree}
M(E_a)=\bigl(|a|_S^{-1/2},|a|_S^{+1/2}\bigr)\in\BT ,
\end{equation}
the half-density normalization of the classical bidegree $(1,q)$.  The
\emph{normalized relation} holds on the nose: $E_a\star E_a^{\dagger}=E_a\star
E_{a^{-1}}=E_1=\Delta$ by \eqref{eq:starmult} and Definition~\ref{def:dagger},
with $\beta_T(ME_a,ME_a)=2$ independent of $a$, matching $d_1d_2=1$; at cycle
level the identity is $\Gamma_F\star\Gamma_F^{t}=q\,\Delta$, and the
half-density divides by $q^{1/2}$ at each slot.  The \emph{rulings} are objects
of the theory, as the canonical plane $\iota(\BT)\subseteq\NAd$ of
Theorems~\ref{thm:rulings} and \ref{thm:splitting}.  These three live in three
different objects --- $\Delta$ and the $E_a$ in $G$, the smeared objects
$E_F$ in $\Corrtr$, the rulings in $\NAd$ --- joined by the comparison maps
$F\mapsto E_F$ and $J$; we do not place all three in a single category.  What is not present is a \emph{non-invertible} correspondence as an
object.
\end{remark}

\section{The signature of the form}\label{sec:signature}

\begin{theorem}\label{thm:signature}
For $F\in\DA$,
\begin{equation}\label{eq:IAF}
I_A(F,F)=T\bigl(F\star F^{\dagger}\bigr)
=\int_{\Ch}|\widehat F(\chi)|^{2}\bigl(\Th(\chi)+c\bigr)d\chi ,
\end{equation}
a real number.  Let $m:=\sup_{\Ch}\Th<\infty$
(Theorem~\ref{thm:symbol}(c)).  Then:
\begin{enumerate}
\item[(a)] $I_A$ has infinite negative index, for every $c$;
\item[(b)] $I_A$ has infinite positive index if and only if $c>-m$, and then
$I_A$ is indefinite with infinite index on both sides;
\item[(c)] $m>0$; in fact
\begin{equation}\label{eq:mbound}
m\ \ge\ \log\pi-\psi(\tfrac14)+\sum_{p\in S_f}(\log p)
\frac{2p^{-1/2}}{1-p^{-1/2}}\ \ge\ \log\pi+\gamma+3\log2+\frac{\pi}{2}
\ \approx\ 5.372 .
\end{equation}
In particular for $c=0$ the hypothesis of (b) holds.
\item[(d)] The $I_A$-orthogonal complement of $\DA^{\mathrm{prim}}$ is zero,
and there is no $\DA$-module splitting of the Tate boundary map
(Theorem~\ref{thm:nosplit}).
\end{enumerate}
\end{theorem}

\begin{proof}
Formula \eqref{eq:IAF} is \eqref{eq:plancherel} applied to $H=F\star
F^{\dagger}$, whose symbol is $|\widehat F|^{2}$, together with the reality of
$\Th$; the value is real because both $|\widehat F|^{2}$ and $\Th+c$ are.

(a) and (b).  Fix a component $\mathfrak c$ and work with profiles: for
$\varphi\in C_c^{\infty}(\R)$ let $F[\varphi]$ be the test with profile
$\varphi$ on $\mathfrak c$ and $0$ elsewhere, so that
$\widehat{F[\varphi]}(\mathfrak c,\tau)=\widehat\varphi(\tau)$ and
$I_A(F[\varphi],F[\varphi])=\int_{\R}|\widehat\varphi(\tau)|^{2}
(\Th(\mathfrak c,\tau)+c)\frac{d\tau}{2\pi}$.  Note that no $\widehat\varphi$
is compactly supported: by Paley--Wiener each is entire of exponential type.
The argument therefore proceeds by \emph{concentration}, not by disjoint
supports.  Two operations stay inside $C_c^{\infty}(\R)$ and act on
$\widehat\varphi$ by the two motions we need: $\varphi(t)\mapsto
e^{iat}\varphi(t)$ translates $\widehat\varphi$ by $a$, and
$\varphi(t)\mapsto\varphi(t/\delta)$ dilates it.  Fix $\varphi_0\in
C_c^{\infty}(\R)$ with $\lVert\widehat\varphi_0\rVert_{L^{2}}=1$;
$\widehat\varphi_0$ is Schwartz.  Put
$\varphi_{a,\delta}(t):=\delta^{-1/2}e^{iat}\varphi_0(t/\delta)$, so that
$\widehat\varphi_{a,\delta}(\tau)=\delta^{1/2}\widehat\varphi_0(\delta(\tau-a))$
has $L^{2}$ norm $1$ and, for any $\rho>0$,
$\int_{|\tau-a|>\rho}|\widehat\varphi_{a,\delta}|^{2}d\tau=
\int_{|s|>\delta\rho}|\widehat\varphi_0(s)|^{2}ds\to0$ as $\delta\to\infty$.
Using the growth bound \eqref{eq:logbound} on the fixed component
$\mathfrak c$, $|\Th(\mathfrak c,\tau)|\le A_{\mathfrak c}+B_S\log(2+|\tau|)$,
\begin{equation}\label{eq:weightedtail}
\int_{|\tau-a|>\rho}\bigl(|\Th(\mathfrak c,\tau)|+|c|\bigr)
|\widehat\varphi_{a,\delta}(\tau)|^{2}d\tau
\le\int_{|s|>\delta\rho}\Bigl(A_{\mathfrak c}+|c|+B_S\log\bigl(2+|a|+|s|/\delta\bigr)\Bigr)
|\widehat\varphi_0(s)|^{2}ds\ \longrightarrow\ 0
\end{equation}
as $\delta\to\infty$, because $\widehat\varphi_0$ is Schwartz and the weight
grows only logarithmically.  So the \emph{weighted} tail, not merely the
$L^{2}$ tail, can be made arbitrarily small outside any window of any fixed
width.  Below we normalize $\lVert\widehat\varphi_0\rVert^{2}=
\int|\widehat\varphi_0|^{2}d\tau/2\pi=1$ and measure all tails with
$d\tau/2\pi$; write $w(\tau):=|\Th(\mathfrak c,\tau)|+|c|$.

\emph{A Gram estimate.}  Let $\psi_1,\psi_2,\dots\in C_c^{\infty}(\R)$ with
$\lVert\widehat\psi_j\rVert=1$, and let $I_1,I_2,\dots$ be pairwise disjoint
intervals such that $\sup_{I_j}w\le L_j$,
$\int_{I_j^{c}}w|\widehat\psi_j|^{2}\frac{d\tau}{2\pi}\le\varepsilon_j$, and,
for $j<k$, $\int_{I_k}w|\widehat\psi_j|^{2}\frac{d\tau}{2\pi}\le\eta_k$.  Then
for $j<k$, splitting $\int w|\widehat\psi_j||\widehat\psi_k|$ over $I_j$,
$I_k$ and the rest, and applying the Cauchy--Schwarz inequality to each piece,
\begin{equation}\label{eq:gram}
|I_A(F[\psi_j],F[\psi_k])|\le\int w|\widehat\psi_j||\widehat\psi_k|
\frac{d\tau}{2\pi}\le(L_j\varepsilon_k)^{1/2}+(\eta_kL_k)^{1/2}
+(\varepsilon_j\varepsilon_k)^{1/2}.
\end{equation}
Each term contains a quantity that is chosen at step $k$, when
$\psi_1,\dots,\psi_{k-1}$ are already fixed: the tail $\varepsilon_k$ of
$\psi_k$, and the mass $\eta_k$ of the earlier functions in the new window.
The construction is therefore sequential.  At step $k$ we first choose the
window $I_k$, disjoint from $I_1,\dots,I_{k-1}$, so that the finitely many
fixed functions $\widehat\psi_j$, $j<k$, have weighted mass at most $\eta_k$
in it (possible either by placing $I_k$ far out, since each
$w|\widehat\psi_j|^{2}$ is integrable, or by taking $I_k$ short, since these
functions are bounded on bounded sets); then we choose
$\psi_k=\varphi_{a_k,\delta_k}$ with $\delta_k$ so large, by
\eqref{eq:weightedtail}, that its weighted tail outside $I_k$ is at most
$\varepsilon_k$ and $\int_{I_k}|\widehat\psi_k|^{2}\frac{d\tau}{2\pi}\ge\tfrac12$.

\emph{(a).}  By Theorem~\ref{thm:symbol}(c), $\Th(\mathfrak c,\cdot)+c$ is
continuous and tends to $-\infty$.  Put $K_0:=\max(2|c|,1)$.  At step $k$ pick
a point $a_k$, so far out that $\Th(\mathfrak c,a_k)+c\le-\tfrac32K_0$, that
the window below is disjoint from the earlier ones, and that the mass condition
holds with $\eta_k:=K_0\,2^{-4k-12}$ for every window of radius at most $1$
around $a_k$ (possible because $\int_{|\tau|>R}w|\widehat\psi_j|^{2}\to0$ as
$R\to\infty$); put $K_k:=-\tfrac23(\Th(\mathfrak
c,a_k)+c)\ge K_0$, and by continuity choose $\rho_k\in(0,1]$ with
$-2K_k\le\Th(\mathfrak c,\cdot)+c\le-K_k$ on $I_k=(a_k-\rho_k,a_k+\rho_k)$.
Then $w\le2K_k+2|c|\le3K_k=:L_k$ on $I_k$.  Take
$\varepsilon_k:=K_0\,2^{-4k-12}$.  With $u_k:=F[\psi_k]/\sqrt{K_k}$ and
$G_{jk}:=I_A(u_j,u_k)$,
\[
G_{kk}\le K_k^{-1}\bigl(-\tfrac12K_k+\varepsilon_k\bigr)\le-\tfrac12+2^{-16},
\]
and by \eqref{eq:gram}, for $j<k$, using $K_j,K_k\ge K_0$,
\[
|G_{jk}|\le\sqrt3\Bigl(\frac{\varepsilon_k}{K_k}\Bigr)^{1/2}
+\sqrt3\Bigl(\frac{\eta_k}{K_j}\Bigr)^{1/2}
+\Bigl(\frac{\varepsilon_j\varepsilon_k}{K_jK_k}\Bigr)^{1/2}
\le(2\sqrt3+1)\,2^{-2k-6}\le2^{-j-k-2}.
\]
For a finite combination $\sum_jc_ju_j$,
\[
I_A\Bigl(\sum c_ju_j,\sum c_ju_j\Bigr)\le\bigl(-\tfrac12+2^{-16}\bigr)\sum|c_j|^{2}
+\sum_{j\ne k}2^{-j-k-2}|c_j||c_k|
\le\Bigl(-\tfrac12+2^{-16}+\tfrac1{12}\Bigr)\sum|c_j|^{2},
\]
using $\sum_{j\ne k}2^{-j-k}|c_j||c_k|\le(\sum_j4^{-j})\sum|c_j|^{2}=\tfrac13
\sum|c_j|^{2}$.  So $I_A$ is negative definite on the span of the $u_j$, which
is infinite dimensional because the same inequality shows that the $u_j$ are
linearly independent.  Only $\Th+c\to-\infty$ was used, so this holds for
every $c$.

\emph{(b).}  If $c\le-m$ then $\Th+c\le0$ everywhere and $I_A\le0$, so the
positive index is $0$.  If $c>-m$, pick $\mathfrak c$ and a window
$I=(a-\rho,a+\rho)$ with $\Th(\mathfrak c,\cdot)+c\ge\kappa>0$ on $I$, possible
since $\Th$ is continuous with supremum $m$, and put $L:=\sup_Iw<\infty$.  Run
the sequential construction with all windows $I_k$ inside $I$ (at step $k$ a
short interval in $(a+2^{-k}\rho,a+2^{-k+1}\rho)$, short enough for the mass
condition), $L_k:=L$ and $\varepsilon_k=\eta_k:=\kappa^{2}2^{-4k-12}/(L+\kappa)$.
Then $I_A(F[\psi_k],F[\psi_k])\ge\tfrac12\kappa-\varepsilon_k\ge\tfrac14\kappa$
and, by \eqref{eq:gram}, $|I_A(F[\psi_j],F[\psi_k])|\le3\kappa2^{-2k-6}\le
\kappa2^{-j-k-4}$ for $j<k$; the computation of (a) with the signs reversed
gives an infinite-dimensional positive definite subspace.  With (a) this makes
$I_A$ indefinite with infinite index on both sides.

(c)  Evaluate $\Th$ at $\tau=0$ on the principal component.  Every finite-place
term of \eqref{eq:Wphat} is $2(\log p)\Rel(\zeta/(1-\zeta))$ with
$\zeta=p^{-1/2}$ real and positive, hence positive and equal to
$(\log p)\cdot2p^{-1/2}/(1-p^{-1/2})$; and the archimedean term is
\eqref{eq:maxarch}, with $\psi(\tfrac14)=-\gamma-3\log2-\pi/2$.  This is
\eqref{eq:mbound}.

(d) is Theorem~\ref{thm:nosplit}.
\end{proof}

We note two consequences.  No statement of Hodge-index type holds on $N_A$ as
it stands: the form is indefinite with infinite
index on both sides, and the primitive sector has no orthogonal complement, so
the shape ``positive on the primitive part of a hyperbolic-plus-definite
decomposition'' is not available.  Any argument requiring a sign must either
change the form --- for instance by composing with a positive operator --- or
restrict to a sector on which $\Th+c$ has a sign, and by (a) such a sector is
never all of $\DA$.  And the sign character depends on the normalizing constant
only in the following way: for $c>-m$ the form is indefinite with infinite index
on both sides, and for $c\le-m$ it is negative definite, which makes every
class negative, rulings included.

\appendix

\section{Auxiliary: the rulings in the limit over \texorpdfstring{$S$}{S}}\label{app:mertens}

\emph{Status.}  This appendix is auxiliary.  Its proof uses the von Mangoldt
function $\Lambda(N)$ and Mertens' theorem, that is, arithmetic input beyond
the product formula, which the construction of
Sections~\ref{sec:carrier}--\ref{sec:signature} does not use.  It is therefore
not part of that construction, nothing in the main text depends on it, and
any statement of the main text that is advertised as free of arithmetic input
beyond the product formula remains so.

Theorem~\ref{thm:scalartrivial} concerns a fixed finite $S$.  In the colimit of
Theorem~\ref{thm:colimtrace} the situation changes, and the reason is
arithmetic: the proof of the next theorem uses Mertens' theorem.

\begin{theorem}\label{thm:mertens}
Work in the colimit chart $\DA^{\infty}$ with $T_{\infty}$,
$I_{\infty}(F,G)=T_{\infty}(F\star G^{\dagger})$ and $c=0$.  Let $\vartheta\in
C_c^{\infty}((0,\infty))$ be real with $\vartheta\ge0$, $\int\vartheta=1$ and
$\supp\vartheta\subseteq[a,b]$, and let $F_n,G_m$ be the spread Tate characters
\eqref{eq:spread}.  Then
\begin{equation}\label{eq:mertenslimit}
\lim_{n,m\to\infty}I_{\infty}(F_n,G_m)=1=\beta_T(e_+,e_-),\qquad
\lim_{n\to\infty}I_{\infty}(F_n,F_n)=0=\beta_T(e_+,e_+),
\end{equation}
the first limit being uniform in the order of the two limits, indeed
$I_{\infty}(F_n,G_m)=1+O(\max(n,m)^{-1})$, and both limits being independent of
the choice of $\vartheta$ in the stated class.  Hence in the colimit chart the
scalar line $L_{\infty}(F,G):=(\C,e^{I_{\infty}(F,G)})$ reproduces
$\lambda_{\mathrm{per}}$ on the rulings, and the vector-valuedness of
$\lambda_{RR}$ at a finite chart is the finite-chart shadow of a single scalar
hyperbolic global pairing.
\end{theorem}

\begin{proof}
\emph{Step 0: the profiles.}  Since $\vartheta$ is real, the profile of
$G_m^{\dagger}$ at $t$ is $\overline{g_m(-t)}=m^{-1}e^{-t/2}\vartheta(t/m)=f_m(t)$,
so $G_m^{\dagger}=F_m$ and therefore $I_{\infty}(F_n,G_m)=T_{\infty}(F_n\star
F_m)$ and $I_{\infty}(F_n,F_n)=T_{\infty}(F_n\star F_n^{\dagger})$.  All four
sections are supported on the principal component and are constantly $1$ on the
compact part; since the compact part is a group of total mass one, the
convolution of the two compact factors is again its indicator, and in the
coordinate $t=\log|w|_S$ convolution is ordinary additive convolution of
profiles.  Writing $\vartheta_n(u):=n^{-1}\vartheta(u/n)$, a probability density on
$[na,nb]$,
\begin{equation}\label{eq:profileconv}
\text{profile of }F_n\star F_m:\ h_{n,m}(t)=e^{-t/2}\phi_{n,m}(t),\qquad
\phi_{n,m}:=\vartheta_n*\vartheta_m ,
\end{equation}
a smooth probability density on $[(n+m)a,(n+m)b]$, with
$\int|\phi'_{n,m}|=\int|\vartheta'_n*\vartheta_m|\le\lVert\vartheta'_n\rVert_{L^{1}}=
\lVert\vartheta'\rVert_{L^{1}}/n$, and by symmetry
$\int|\phi'_{n,m}|\le\lVert\vartheta'\rVert_{L^{1}}/\max(n,m)$.

\emph{Step 1: the finite places sum to a Mertens sum.}  For a section $H$
supported on the principal component with profile $h$, formula \eqref{eq:Wp}
gives, exactly as in the proof of Theorem~\ref{thm:symbol}(a) --- the
coefficient at the shell $k$ being $h(-k\log p)$, since $j_p(p^{k}u)$ has
archimedean coordinate $p^{-k}$ and all compact coordinates units, on which $H$
is constant ---
\begin{equation}\label{eq:Wpprofile}
W_p(H)=(\log p)\sum_{k\ne0}p^{-|k|/2}h(-k\log p).
\end{equation}
Apply this to $H=F_n\star F_m$ and $h=h_{n,m}$.  Since $\phi_{n,m}$ is
supported in $(0,\infty)$, only $k=-j$ with $j\ge1$ contributes, and there
$h(j\log p)=p^{-j/2}\phi_{n,m}(j\log p)$ while $p^{-|k|/2}=p^{-j/2}$; so
$W_p(F_n\star F_m)=(\log p)\sum_{j\ge1}p^{-j}\phi_{n,m}(j\log p)$, and summing
over all primes, with $N=p^{j}$ and $\Lambda(N)=\log p$,
\begin{equation}\label{eq:mertenssum}
\sum_pW_p(F_n\star F_m)=\sum_{N\ge2}\frac{\Lambda(N)}{N}\,\phi_{n,m}(\log N).
\end{equation}
The sum is finite: $\phi_{n,m}(\log N)=0$ unless $\log N\in[(n+m)a,(n+m)b]$.
In particular $T_{\infty}(F_n\star F_m)$ is the finite sum of
Theorem~\ref{thm:colimtrace}(b), and \eqref{eq:mertenssum} is its finite-place
part.

\emph{Step 2: partial summation.}  Let $A(x):=\sum_{N\le x}\Lambda(N)/N$.  By
Mertens' theorem \cite{Mertens}, which is elementary, $A(x)=\log x+E(x)$ with
$|E(x)|\le C_0$ for $x\ge1$.  Since $\phi_{n,m}$ is smooth with compact support
in $(0,\infty)$, Riemann--Stieltjes integration by parts on $[1,\infty)$, with
vanishing boundary terms, gives
\[
\sum_N\frac{\Lambda(N)}{N}\phi_{n,m}(\log N)
=\int_1^{\infty}\phi_{n,m}(\log x)\,dA(x)
=-\int_1^{\infty}A(x)\,\phi'_{n,m}(\log x)\,\frac{dx}{x} ,
\]
which after the substitution $x=e^{t}$ and $A(e^{t})=t+E(e^{t})$ equals
\[
-\int_0^{\infty}\bigl(t+E(e^{t})\bigr)\phi'_{n,m}(t)\,dt ,
\]
which is $\int_0^{\infty}\phi_{n,m}(t)dt-\int_0^{\infty}E(e^{t})
\phi'_{n,m}(t)dt$, using $-\int t\phi'=\int\phi$.  The first term is $1$; the
second is bounded by $C_0\int|\phi'_{n,m}|\le
C_0\lVert\vartheta'\rVert_{L^{1}}/\max(n,m)$.  Hence
\begin{equation}\label{eq:step2}
\sum_pW_p(F_n\star F_m)=1+O\bigl(\max(n,m)^{-1}\bigr),
\end{equation}
with an implied constant depending only on $\vartheta$, uniform in the other index,
so that the double limit is $1$ in either order.

\emph{Step 3: the archimedean place is negligible.}  The profile $h_{n,m}$
vanishes for $t<0$ and on a neighbourhood of $t=0$, its support starting at
$(n+m)a>0$; so \eqref{eq:Winf} reduces to $W_{\infty}(F_n\star
F_m)=\int_0^{\infty}K(t)h_{n,m}(t)dt$.  On $t\ge(n+m)a$ one has
$K(t)\le(1-e^{-2a})^{-1}e^{-t/2}$, so
\[
|W_{\infty}(F_n\star F_m)|\le(1-e^{-2a})^{-1}\int
e^{-t}\phi_{n,m}(t)\,dt\le(1-e^{-2a})^{-1}e^{-(n+m)a}\longrightarrow0
\]
exponentially.  Adding to \eqref{eq:step2} gives the first limit of
\eqref{eq:mertenslimit}.

\emph{Step 4: the self-pairing.}  Put $H':=F_n\star F_n^{\dagger}$ and let $\eta$ denote its profile.  From
$(f_n*f_n^{\dagger})(t)=\int f_n(u)\overline{f_n(u-t)}du$ with
$f_n(u)=n^{-1}e^{-u/2}\vartheta(u/n)$, the substitution $u=ns$ gives
$\eta(t)=n^{-1}e^{t/2}\int_{\R}e^{-ns}\vartheta(s)\vartheta(s-t/n)\,ds$, supported in
$|t|\le n(b-a)$.  On the support one has $s\ge a$ and $s-t/n\ge a$, whence
\begin{equation}\label{eq:hprime}
|\eta(t)|\le C_{\vartheta}n^{-1}e^{-na}e^{-|t|/2},\qquad
|\eta(t)|+|\eta'(t)|\le C'_{\vartheta}n^{-1}e^{-na}e^{-|t|/2}\quad(n\ge1):
\end{equation}
for $t>0$ one uses $e^{-ns}\le e^{-na-t}$, which follows from $s\ge a+t/n$,
against the prefactor $e^{t/2}$; for $t\le0$ one uses simply $e^{-ns}\le
e^{-na}$ with the prefactor $e^{t/2}=e^{-|t|/2}$; and the derivative bound adds
only $\tfrac12\eta$ and a term carrying $\vartheta'$.

At the finite places, \eqref{eq:Wpprofile} gives
$|W_p(H')|\le(\log p)\sum_{k\ne0}p^{-|k|/2}|\eta(-k\log p)|$.  Write
$\eta(t)=e^{t/2}B(t)$ with
$B(t)=n^{-1}\int_{\R}e^{-ns}\vartheta(s)\vartheta(s-t/n)\,ds$.  The two regimes of the
support give
\[
|B(t)|\le C_{\vartheta}\,n^{-1}e^{-na}e^{-t}\ \ (t>0),\qquad
|B(t)|\le C_{\vartheta}\,n^{-1}e^{-na}\ \ (t\le0),
\]
the first from $s\ge a+t/n$ on the support and the second from $s\ge a$.
Hence for $k=-j<0$, where $t=j\log p>0$,
$p^{-j/2}|\eta(j\log p)|=p^{-j/2}p^{j/2}|B(t)|\le
C_{\vartheta}n^{-1}e^{-na}p^{-j}$; and for $k=j>0$, where $t=-j\log p\le0$,
$p^{-j/2}|\eta(-j\log p)|=p^{-j/2}p^{-j/2}|B(t)|\le
C_{\vartheta}n^{-1}e^{-na}p^{-j}$.  Summing,
\[
\sum_p|W_p(H')|\le 2C_{\vartheta}n^{-1}e^{-na}
\sum_{N\le e^{n(b-a)}}\frac{\Lambda(N)}{N}
\le 2C_{\vartheta}n^{-1}e^{-na}\bigl(n(b-a)+C_0\bigr)\longrightarrow0
\]
exponentially, by Mertens' theorem again.  At the real place, split
\eqref{eq:Winf} at $t=1$: near $0$ one has $K(t)\le1/t$ and the bracket
$\eta(t)+\eta(-t)-2\eta(0)e^{-t/2}$ is $O(t\sup|\eta'|)$, so
$\int_0^{1}K|\cdots|\le C\sup|\eta'|$; for $t\ge1$ one has
$K(t)\le1.2\,e^{-t/2}$ and \eqref{eq:hprime} gives an integrable product.
With $|\eta(0)|\le C_{\vartheta}n^{-1}e^{-na}$, all
three contributions are $O(n^{-1}e^{-na})$.  Hence $I_{\infty}(F_n,F_n)\to0$.

The limits $1$ and $0$ do not involve $\vartheta$; the error terms are
$O_{\vartheta}(\max(n,m)^{-1})$ and $O_{\vartheta}(e^{-na})$.
\end{proof}

\begin{remark}\label{rem:mertensmeaning}
The value $1$ in \eqref{eq:mertenslimit} is the logarithmic density of prime
powers, $\sum_{N\le x}\Lambda(N)/N\sim\log x$.  The structure constant of the
hyperbolic ruling pairing is, in the colimit chart, that density.  The theorem
does not repair the finite chart: at finite $S$ the primes carrying the value
$1$ lie outside $S$, so $I_A^{S}$ does vanish on the rulings and
Theorem~\ref{thm:scalartrivial} stands, which is why $\lambda_{RR}$ is
vector-valued.  The chart dependence is the residue overcount
\eqref{eq:residues} seen from the other side.
\end{remark}

\begin{thebibliography}{99}

\bibitem{CCWeilArch}
A.~Connes and C.~Consani, \emph{Weil positivity and trace formula: the
archimedean place}, Selecta Math. (N.S.) \textbf{27} (2021), no.~4, Paper
No.~77.

\bibitem{CuntzQuillen}
J.~Cuntz and D.~Quillen, \emph{Excision in bivariant periodic cyclic
cohomology}, Invent. Math. \textbf{127} (1997), 67--98.

\bibitem{DLMF}
NIST Digital Library of Mathematical Functions, \texttt{https://dlmf.nist.gov/},
Chapter 5 (Gamma function), \S5.4(iii) and \S5.9; also E.~T.~Whittaker and
G.~N.~Watson, \emph{A Course of Modern Analysis}, 4th ed., Cambridge, 1927,
\S12.3.

\bibitem{HartshorneRD}
R.~Hartshorne, \emph{Residues and Duality}, Lecture Notes in Math. \textbf{20},
Springer, 1966; Chapter III, \S7.

\bibitem{Brylinski}
J.-L.~Brylinski, \emph{Algebras associated with group actions and their
homology}, Brown University preprint, 1987.

\bibitem{Deligne}
P.~Deligne, \emph{Le d\'eterminant de la cohomologie}, in: Current Trends in
Arithmetical Algebraic Geometry, Contemp. Math. \textbf{67}, Amer. Math. Soc.,
1987, 93--177.

\bibitem{GetzlerJones}
E.~Getzler and J.~D.~S.~Jones, \emph{The cyclic homology of crossed product
algebras}, J. Reine Angew. Math. \textbf{445} (1993), 161--174; \S3.

\bibitem{KnudsenMumford}
F.~Knudsen and D.~Mumford, \emph{The projectivity of the moduli space of stable
curves I: Preliminaries on ``det'' and ``Div''}, Math. Scand. \textbf{39}
(1976), 19--55.

\bibitem{MacLane}
S.~Mac~Lane, \emph{Homology}, Grundlehren der mathematischen Wissenschaften
\textbf{114}, Springer, 1963; Chapter VIII.

\bibitem{Mertens}
F.~Mertens, \emph{Ein Beitrag zur analytischen Zahlentheorie},
J. Reine Angew. Math. \textbf{78} (1874), 46--62.

\bibitem{MeyerCyclic}
R.~Meyer, \emph{Local and Analytic Cyclic Homology}, EMS Tracts in Mathematics
\textbf{3}, European Mathematical Society, 2007.

\bibitem{Nistor}
V.~Nistor, \emph{Group cohomology and the cyclic cohomology of crossed
products}, Invent. Math. \textbf{99} (1990), 411--424; \S\S3--4.

\bibitem{SGA6}
P.~Berthelot, A.~Grothendieck and L.~Illusie, \emph{Th\'eorie des intersections
et th\'eor\`eme de Riemann--Roch (SGA 6)}, Lecture Notes in Math. \textbf{225},
Springer, 1971.

\bibitem{Schaefer}
H.~H.~Schaefer, \emph{Topological Vector Spaces}, Graduate Texts in Math.
\textbf{3}, Springer, 1971.

\bibitem{Treves}
F.~Tr\`eves, \emph{Topological Vector Spaces, Distributions and Kernels},
Academic Press, 1967.

\bibitem{BrodzkiLykova}
J.~Brodzki and Z.~A.~Lykova, \emph{Excision in cyclic type homology of
Fr\'echet algebras}, Bull. London Math. Soc. \textbf{33} (2001), 283--291.

\bibitem{Tate}
J.~Tate, \emph{Fourier analysis in number fields and Hecke's zeta-functions},
thesis, Princeton, 1950; reprinted in: J.~W.~S.~Cassels and A.~Fr\"ohlich
(eds.), \emph{Algebraic Number Theory}, Academic Press, 1967, 305--347.

\bibitem{Weil1948}
A.~Weil, \emph{Vari\'et\'es ab\'eliennes et courbes alg\'ebriques}, Hermann,
Paris, 1948.

\bibitem{Weil1952}
A.~Weil, \emph{Sur les ``formules explicites'' de la th\'eorie des nombres
premiers}, Comm. S\'em. Math. Univ. Lund (1952), 252--265.

\bibitem{Wodzicki}
M.~Wodzicki, \emph{Excision in cyclic homology and in rational algebraic
K-theory}, Ann. of Math. (2) \textbf{129} (1989), 591--639.

\end{thebibliography}

\end{document}
